\documentclass[11pt]{article}
\usepackage[a4paper,margin=19mm]{geometry}
\usepackage{fontspec}
\setmainfont{lmroman10-regular.otf}
\usepackage{amsmath,amssymb,amsthm,mathtools,amscd,xcolor,graphicx,hyperref}
\hypersetup{hidelinks}
\numberwithin{equation}{section}
\newtheorem{theorem}{Theorem}[section]
\newtheorem{lemma}[theorem]{Lemma}
\newtheorem{proposition}[theorem]{Proposition}
\newtheorem{corollary}[theorem]{Corollary}
\newtheorem{claim}[theorem]{Claim}
\newtheorem{fact}[theorem]{Fact}
\newtheorem{conjecture}[theorem]{Conjecture}
\theoremstyle{definition}
\newtheorem{definition}[theorem]{Definition}
\newtheorem{example}[theorem]{Example}
\newtheorem{remark}[theorem]{Remark}
\newtheorem{question}[theorem]{Question}
\theoremstyle{plain}
\newcommand{\C}{\mathbb C}
\newcommand{\R}{\mathbb R}
\newcommand{\Q}{\mathbb Q}
\newcommand{\Z}{\mathbb Z}
\newcommand{\N}{\mathbb N}
\newcommand{\eee}{\mathrm e}
\newcommand{\cclass}[1]{\textsf{#1}}
\newcommand{\cproblem}[1]{\textsc{#1}}
\newcommand{\ie}{i.\,e.}
\newcommand{\Ie}{I.\,e.}
\newcommand{\eg}{e.\,g.}
\newcommand{\Eg}{E.\,g.}
\newcommand{\phd}{Ph.\,D.}
\newcommand{\msc}{M.\,S.}
\newcommand{\bsc}{B.\,S.}
\newcommand{\expref}[2]{\hyperref[#2]{#1~\ref{#2}}}
\newcommand{\expeqref}[2]{\hyperref[#2]{#1~\eqref{#2}}}
\newcommand{\exprefsilent}[2]{\hyperref[#2]{#1}}
\usepackage{eprint}
\usepackage{amsmath}
\usepackage{amsthm}
\usepackage{latexsym}
\usepackage{xspace}
\usepackage{amssymb}
\usepackage{fancybox}
\usepackage{mathtools}
\usepackage{thm-restate}
\usepackage{dsfont}
\usepackage{cite}

\newcommand{\genericdomain}{\mathcal{X}}
\newcommand{\domain}{\mathcal{D}}
\newcommand{\reals}{\mathbb{R}}
\newcommand{\E}{\mathbb{E}}
\newcommand{\LL}{\mathbb{L}}
\newcommand{\fun}[3]{ #1 : #2 \to #3}
\newcommand{\sign}{\mathrm{sign}}
\newcommand{\sgn}{\mathrm{sign}}
\newcommand{\cF}{\mathcal{F}}
\renewcommand{\P}{\mathcal{P}}

\newcommand{\opt}{\mathrm{opt}}
\newcommand{\Maj}{\mathrm{Maj}}
\newcommand{\poly}{\mathrm{poly}}
\newcommand{\quasipoly}{\mathrm{quasipoly}}
\newcommand{\polylog}{\mathrm{polylog}}

\newcommand{\key}{\mathrm{key}}
\newcommand{\up}{\mathrm{up}}
\newcommand{\ukn}{\mathrm{ukn}}
\newcommand{\pure}{\mathrm{pure}}

\newcommand{\ED}{\mathsf{ ED}}
\newcommand{\omb}{\mathsf{ ODDMAXBIT}}
\newcommand{\SURJ}{\mathsf{ SURJ}}
\newcommand{\SDU}{\mathsf{ SDU}}
\newcommand{\SE}{\mathsf{ IST}}
\newcommand{\dumSE}{\mathsf{ dSE}}
\newcommand{\EA}{\mathsf{ EA}}
\newcommand{\Surjectivity}{\mathsf{Surjectivity}}
\newcommand{\DIST}{\mathsf{DIST}}
\newcommand{\dumSURJ}{\mathsf{dSURJ}}
\newcommand{\dumDIST}{\mathsf{dDIST}}
\newcommand{\psicomb}{\psi_{\OMB \circ \OR}}
\newcommand{\psiin}{\psi_{\OR}}

\newcommand{\OR}{\mathsf{ OR}}
\newcommand{\bbI}{\mathbb{I}}
\newcommand{\ORinv}{{\OR^{-1}(+1)}}
\newcommand{\GAPMAJ}{\mathsf{GAPMAJ}}
\newcommand{\MAJ}{\mathsf{MAJ}}
\newcommand{\NOR}{\mathsf{ NOR}}
\newcommand{\XOR}{\mathsf{PAR}}
\newcommand{\THR}{\mathsf{THR}}
\newcommand{\bits}{\{-1,1\}}
\newcommand{\zbits}{\{0,1\}}
\newcommand{\B}{\{0,1\}}
\newcommand{\zbn}{\zbits^n}
\newcommand{\bn}{\bits^n}
\newcommand{\isafunc}{: \bn \rightarrow \bits}
\newcommand{\sse}{\subseteq}

\newcommand{\ceq}{\subseteq}

\newcommand{\promcomp}{G^{\le N}}
\newcommand{\symcomp}{G^{\mathsf{prop}}}
\newcommand{\symsymcomp}{\tilde{G}^{\mathsf{prop}}}
\newcommand{\ls}{\star}
\newcommand{\eat}[1]{}
\newcommand{\trans}{{\ensuremath{\mathsf{T}}}}           
\newcommand{\F}{\mathbb{F}}
\newcommand{\Ot}{\tilde{O}}
\newcommand{\Omegat}{\tilde{\Omega}}
\newcommand{\tO}{\tilde{O}}
\newcommand{\tOmega}{\tilde{\Omega}}
\newcommand{\tTheta}{\tilde{\Theta}}
\newcommand{\help}{\mathfrak{h}}
\newcommand{\func}{G}
\newcommand{\alg}{\mathcal{A}}
\newcommand{\cB}{\mathcal{B}}
\newcommand{\PP}{\mathbf{PP}}
\newcommand{\PPcc}{\mathbf{PP}^{\text{cc}}}
\newcommand{\UPP}{\mathbf{UPP}}
\newcommand{\UPPcc}{\mathbf{UPP}^{\text{cc}}}
\newcommand{\cP}{\mathcal{P}}
\newcommand{\cV}{\mathcal{V}}
\newcommand{\cQ}{\mathcal{Q}}
\newcommand{\cT}{\mathcal{T}}
\newcommand{\fail}{\bot}
\newcommand{\cX}{\mathcal{X}}
\newcommand{\Ind}{\mathbb{I}}
\newcommand{\eps}{\varepsilon}
\newcommand{\bx}{\mathbf{x}}
\newcommand{\by}{\mathbf{y}}
\newcommand{\br}{\mathbf{r}}
\newcommand{\bv}{\mathbf{v}}
\newcommand{\bs}{\mathbf{s}}
\newcommand{\idx}{\textsc{index}\xspace}
\newcommand{\universe}{{\mathcal U}}
\newcommand{\mnote}[1]{\normalmarginpar \marginpar{\tiny #1}}
\newcommand{\edit}[1]{{{#1}}}
\newcommand{\etal}{{\textit{et al.}}}
\newcommand{\wb}{\mathrm{wb}}
\newcommand{\half}{\frac{1}{2}}
\newcommand{\Znk}{Z_{n,k}}
\newcommand{\argmin}[1]{\underset{#1}{\mathrm{argmin}}}
\newcommand{\argmax}[1]{\underset{#1}{\mathrm{argmax}}}
\newcommand{\summ}{\displaystyle \sum}
\newcommand{\intt}{\displaystyle\int}
\newcommand{\var}{\text{Var}}
\newcommand{\nchoosek}[2]{\left(\begin{array}{*{20}c}#1\\#2\end{array}\right)}  %
\newcommand{\ignore}[1]{}
\newcommand{\provisionallyremove}[1]{}

\newcommand{\ba}{\mathbf{a}}
\newcommand{\be}{\mathbf{e}}
\newcommand{\bw}{\mathbf{w}}
\newcommand{\bb}{\mathbf{b}}
\newcommand{\bu}{\mathbf{u}}

\newcommand{\bz}{\mathbf{z}}
\newcommand{\bc}{\mathbf{c}}
\newcommand{\bd}{\mathbf{d}}
\newcommand{\bh}{\mathbf{h}}
\newcommand{\bmu}{\boldsymbol{\mu}}
\newcommand{\balpha}{\boldsymbol{\alpha}}
\newcommand{\bbeta}{\boldsymbol{\beta}}
\newcommand{\btau}{\boldsymbol{\tau}}
\newcommand{\bxi}{\boldsymbol{\xi}}
\newcommand{\blambda}{\boldsymbol{\lambda}}
\newcommand{\bepsilon}{\boldsymbol{\epsilon}}
\newcommand{\bsigma}{\boldsymbol{\sigma}}
\newcommand{\btheta}{\boldsymbol{\theta}}
\newcommand{\bomega}{\boldsymbol{\omega}}
\newcommand{\Lcal}{\mathcal{L}}
\newcommand{\Ocal}{\mathcal{O}}
\newcommand{\Acal}{\mathcal{A}}
\newcommand{\Gcal}{\mathcal{G}}
\newcommand{\Ccal}{\mathcal{C}}
\newcommand{\Xcal}{\mathcal{X}}
\newcommand{\Jcal}{\mathcal{J}}
\newcommand{\Dcal}{\mathcal{D}}
\newcommand{\Fcal}{\mathcal{F}}
\newcommand{\Hcal}{\mathcal{H}}
\newcommand{\Rcal}{\mathcal{R}}
\newcommand{\Ncal}{\mathcal{N}}
\newcommand{\Scal}{\mathcal{S}}
\newcommand{\Pcal}{\mathcal{P}}
\newcommand{\Qcal}{\mathcal{Q}}
\newcommand{\Wcal}{\mathcal{W}}
\newcommand{\field}{\mathbb{F}}
\newcommand{\adeg}{\widetilde{\operatorname{deg}}}
\newcommand{\ubdeg}{\widetilde{\operatorname{ubdeg}}}
\newcommand{\dpdeg}{\widetilde{\operatorname{dpdeg}}}
\newcommand{\tV}{\tilde{V}}
\newcommand{\AND}{\mathsf{AND}}
\newcommand{\PARITY}{\mathsf{PARITY}}
\newcommand{\OMB}{\mathsf{OMB}}
\newcommand{\GAPAND}{\mathsf{GapAND}}
\newcommand{\Ucal}{\mathcal{U}}
\newcommand{\Ld}{\tilde{L}}
\newcommand{\uloss}{\ell^\star}
\newcommand{\loss}{\mathcal{L}}
\newcommand{\losst}{\ell_t}
\newcommand{\norm}[1]{\left\|#1\right\|}
\newcommand{\inner}[1]{\langle#1\rangle}
\newcommand{\COMPAO}{\psi_{\textbf{AND-OR}}}
\renewcommand{\sgn}{\operatorname{sgn}}
\title{Tight threshold degree of the Surjectivity function}
\author{Mark Bun, Justin Thaler}
\date{doi:10.4086/toc.2021.v017a007; CC BY 3.0}
\begin{document}
\maketitle
\noindent\textit{Editorial note.} Complete original mathematical paper, all supporting author proofs and original published bibliography retained. Native frontmatter/biographies remain in raw source/publication PDF. Existing article interface adapts XeLaTeX presentation only.\par\medskip
\noindent\textit{Editorial note.} One primary construction family is counted by this task. Related consequences, special cases and restatements remain complete supporting context. Original theorem/environment names and all exact assumptions retained.\par\medskip
\noindent\textit{Editorial note.} Exact original CC BY3.0 retained. Authors own mathematics; no proof review, mathematical repair, added assumptions, content hashes or cross-version comparison.\par\medskip
\noindent\textit{Allowed standard prerequisites (editorial).}\begin{itemize}
\item Original dual-polynomial or affine-code definitions and cited standard tools. All original degree/error/family/depth or field/dimension/minimum-distance/soundness assumptions retained literally.. Complete source-local correction/amplification or sampling/agreement construction supplied, without collector reproof or certification.
\end{itemize}\medskip
\section{Introduction}
The \emph{threshold degree} of a Boolean function $f\colon \bits^n \to \bits$, denoted $\deg_{\pm}(f)$, 
is the least degree of a real polynomial $p$ that sign-represents $f$, \ie, $p(x) \cdot f(x) > 0$
for all $x \in \bits^n$.
A closely related notion is the $\eps$-approximate degree of $f$, denoted $\adeg_{\eps}(f)$, which is the least degree of a 
real polynomial $p$ such that $|p(x)-f(x)| \leq \eps$ for all $x \in \bits^n$. 

The parameter setting $\eps=1$ is a degenerate case: $\adeg_1(f)=0$ because 
the constant 0 function approximates any Boolean $f$ to error $\eps=1$. 
However, as soon as $\eps$ is strictly less than 1, $\eps$-approximate degree is a highly non-trivial
notion with a rich mathematical theory.
In particular, it is easily seen that 
$$\deg_{\pm}(f)=\lim_{\eps \nearrow 1} \adeg_{\eps}(f).$$ In other words, threshold degree is equivalent to the notion of $\eps$-approximate
degree when $\eps$ is permitted to be \emph{arbitrarily} close to (but strictly less than) 1.\footnote{It is known that for any $d > 0$, there are functions of threshold degree $d$
	that cannot be approximated by degree $d$ polynomials to error better than $1-2^{-\tilde{\Omega}(n^d)}$~\cite{podolskii2007perceptrons},
	and this bound is tight~\cite{bvdw}. Hence, threshold degree is also equivalent to the notion of $\eps$-approximate degree for some value of $\eps$
	that is \emph{doubly}-exponentially close to $1$.}

In this paper, we are concerned with proving %
lower bounds on 
the %
$\eps$-approximate degree  %
 when either:

\begin{itemize}
	\item  $\eps$ is \emph{arbitrarily} close to 1, or
	\item $\eps$ is \emph{exponentially} close to 1 (\ie, $\eps=1-2^{-n^{1-\delta}}$ 
	for some constant $\delta > 0$).
\end{itemize}
The former parameter regime captures threshold degree, while we refer to the latter as  \emph{large-error}
approximate degree.
%
While the approximate and threshold degree of a function $f$ capture simple statements about its approximability by polynomials, these quantities relate intimately to the complexity of computing $f$ in concrete computational models. Specifically, the query complexity models $\UPP^{\text{dt}}$ and $\PP^{\text{dt}}$, and the communication models $\UPP^{\text{cc}}, \PP^{\text{cc}}$, are all defined (see \expref{Section}{s2}) as natural analogs of the Turing machine class $\PP$, which in turn captures probabilistic computation with arbitrarily small advantage over random guessing. Briefly, both $\UPP^{\text{dt}}$ and $\PP^{\text{dt}}$ capture the minimum number of queries required to evaluate a Boolean function with probability at least $1/2$ on every input;  $\UPP^{\text{cc}}, \PP^{\text{cc}}$ do the same for two-party communication problems.
It is known that the threshold degree of $f$ is equivalent to its complexity $\UPP^{\text{dt}}(f)$, while the logarithm of a fundamental matrix-analytic analog of threshold degree known as \emph{sign-rank} characterizes $\UPP^{\text{cc}}$. \label{s1}
Similarly, large-error approximate degree characterizes the query complexity measure $\PP^{\text{dt}}$, in the following 
sense: for any $d > 0$, $\adeg_{1-2^{-d}}(f) \geq \Omega(d) \Longleftrightarrow \PP^{\text{dt}}(f) \geq \Omega(d)$. As described in \expref{Section}{sec:applications}, $\PP^{\text{dt}}$ is also closely related to the notion of \emph{threshold weight}. 
\expref{Section}{s2} elaborates on these models and their many applications in learning theory, circuit complexity, and cryptography.
\medskip \noindent \textbf{Our results in a nutshell.}
We prove two results about the threshold degree and large-error approximate degree
of functions in AC$^0$.\footnote{AC$^0$ is the non-uniform class of sequences of functions computed by polynomial-size %
Boolean circuits of constant depth.} 

First, we prove a tight $\tilde{\Omega}(n^{1/2})$ lower bound
on the threshold degree (\ie,  $\UPP^{\text{dt}}$ complexity) of a natural function called $\mathsf{SURJECTIVITY}$,
which is computed by a depth-three %
circuit with logarithmic bottom fan-in.
This matches the previous best threshold degree lower bound for any AC$^0$
function, due to Sherstov~\cite{sherstov15}. Our analysis is 
much simpler than
Sherstov's, which takes up the bulk of a (70+)-page manuscript~\cite{sherstov15}.
%
An additional advantage of our analysis is that our lower bound on the threshold degree of $\mathsf{SURJECTIVITY}$ ``lifts'' to give a lower bound for the communication analog  $\UPP^{\text{cc}}$ as well.
In particular, we obtain an $\Omega(n^{1/2})$ 
$\UPP^{\text{cc}}$
lower bound for a related AC$^0$ function; this improves over the previous best $\UPP^{\text{cc}}$
lower bound for AC$^0$, of $\Omega(n^{2/5})$~\cite{bt16}. 

Second, we give nearly optimal bounds on the large-error approximate degree (and hence, $\PP^{\text{dt}}$ complexity) of AC$^0$. For any constant $\delta > 0$, we show that there is an AC$^0$ function with $\eps$-approximate degree
$\Omega(n^{1-\delta})$, where $\eps=1-2^{-\Omega(n^{1-\delta})}$. This result lifts to an analogous $\PP^{\text{cc}}$ lower bound. 

To summarize our results succinctly:
\begin{itemize}
	\item We prove an $\tilde{\Omega}(n^{1/2})$ lower bound on the $\UPP$ complexity of $\mathsf{SURJECTIVITY}$ in the query setting,
	and of a related AC$^0$ function in the communication setting.
	\item We prove an $\Omega(n^{1-\delta})$ lower bound on the $\PP$ complexity of some AC$^0$ circuit of depth $O(1/\delta)$, in both the query and communication settings.
\end{itemize}

\expref{Table}{tab:tab} compares our new lower bounds for AC$^0$ to the long line of prior %
work   %
with similar goals.




\medskip \noindent \textbf{Context and prior work.}
The study of both large-error approximate degree and threshold degree
has led to many breakthrough results in theoretical computer science,
especially in the algorithmic and complexity-theoretic study of
constant-depth %
circuits.
For example, threshold degree upper bounds are at the core of many of the fastest known PAC learning algorithms.
This includes the notorious case of polynomial-size %
CNF formulas on $n$ variables, for which the fastest known algorithm~\cite{ksdnf} runs in time $\exp(\tilde{O}(n^{1/3}))$
owing to an $\tilde{O}(n^{1/3})$ upper bound on the threshold degree of any such formula. This upper bound \label{sec:MP}
is tight, matching a classic $\Omega(n^{1/3})$ lower bound of Minsky and Papert~\cite{mp} for the following read-once CNF:
$\AND_{n^{1/3}} \circ \OR_{n^{2/3}}$ (here, we use \label{s:footyfootfoot}
subscripts to clarify the number of inputs on which a function is defined). 

In complexity theory, breakthrough results
of Sherstov \cite{sherstovmajmaj, patmat} and Buhrman et al.~\cite{bvdw}
used lower bounds on large-error approximate degree
to show that there are AC$^0$ functions with polynomial $\PP^{\text{cc}}$
complexity. 
One notable implication  of these results is that Allender's~\cite{allender} classic simulation
of AC$^0$ functions by depth-three majority circuits is optimal. (This resolved an open problem of Krause and Pudl{\'{a}}k~\cite{kp}.)
A subsequent, related breakthrough of Razborov and Sherstov~\cite{razborovsherstov} used Minsky and Papert's lower bound 
on the threshold degree of $\AND_{n^{1/3}} \circ \OR_{n^{2/3}}$ to prove the first polynomial $\UPPcc$
lower bound for a function in AC$^0$, answering an old open question of Babai et al.~\cite{babaietal}. 

\medskip

These breakthrough  lower bounds raised the intriguing possibility that AC$^0$ functions could be \emph{maximally} hard for
the $\UPP^{\text{cc}}$ and $\PP^{\text{cc}}$ communication models, as well as for related complexity measures. Nevertheless, the quantitative parameters achieved in these %
papers   %
are far from actually showing that this is the case. Indeed, the following basic questions about the complexity of AC$^0$ remain open.

\label{s:openproblem1}

\medskip \noindent \textbf{Open Problem 1.} Is there an AC$^0$ function $F \colon \bits^{n \times n} \to \bits$  with 
$\UPP^{\text{cc}}$ complexity $\Omega(n)$?

\medskip \noindent \textbf{Open Problem 2.} Is there an AC$^0$ function $F \colon \bits^{n \times n} \to \bits$  with 
$\PP^{\text{cc}}$ complexity  $\Omega(n)$?

\medskip 
An affirmative answer to either question would be tight: \emph{Every} function $F \colon \bits^n \times \bits^n \to \bits$
has  $\UPP^{\text{cc}}$ and  $\PP^{\text{cc}}$ complexity at most $n$. %
Obtaining an affirmative answer to Open Problem 1 is
harder than for Open Problem 2, since $\UPP^{\text{cc}}(f) \leq \PP^{\text{cc}}(f)$ for all $f$. 



Guided by these open problems, a sequence
of %
papers  %
has established quantitatively stronger and more general lower bounds for AC$^0$ functions \cite{bt13, bt14, sherstov14, sherstov15, dlweight, bt16, adegsurj}.
In addition to making partial progress toward resolving these questions, the techniques developed in %
this body of work  %
have found fruitful applications in new domains.
For example, Bouland et al.~\cite{bchtv} built on techniques from
%
some of these papers  %
\cite{bt13, bt14, sherstov14, bt16} to resolve several 
old open questions about the relativized power of statistical zero knowledge proofs and their variants.
%
As another example, our recent 
%
papers
\cite{adegsurj, bkt} built on the same line of work to
resolve or nearly resolve a number of longstanding open questions in quantum query complexity.
Finally, large-error and threshold degree lower bounds on AC$^0$ functions
have recently proved instrumental in the development of cryptographic secret-sharing
schemes with reconstruction procedures in AC$^0$ \cite{viola, bogicalp, cheng2017near}. We thus believe that the new techniques developed in this
%
article  %
will find %
further applications, perhaps in unexpected areas.


Prior to our work, the best known result toward a resolution of Open Problem 1 was an $\Omega(n^{2/5})$ lower
bound on $\UPP^{\text{cc}}$ complexity of an AC$^0$ function~\cite{bt16}, while 
the best known result toward Open Problem 2
was an $\Omega(n^{1/2})$ 
bound on the $\PP^{\text{cc}}$ complexity of a very complicated AC$^0$ circuit~\cite{sherstov15}. \label{sec:previoussec}

\label{s:openquestion1}

\subsection{Our results in detail}


\subsubsection{Resolving the threshold degree of
  \texorpdfstring{$\mathsf{SURJECTIVITY}$}{SURJECTIVITY}}
\label{surjoverview}
\noindent \textbf{Surjectivity and its history.} Let $R$ be a power of 2 and $n=N \log R$. The function $\mathsf{SURJECTIVITY}_{R, N}$ %
($\SURJ_{R, N}$ for short) is defined as follows. Given an input in $\bits^n$, $\SURJ_{R, N}$ interprets the input as a list of $N$ numbers $(s_1, \dots, s_N)$
from a range $[R]\coloneqq \{1, \dots, R\}$, and %
evaluates to $-1$ if and only if every element of the range $[R]$ appears at least once in the list.\footnote{As is standard, we associate $-1$ with logical TRUE and $+1$ with logical FALSE throughout.} $\SURJ_{R, N}$ is computed by an AC$^0$ circuit
of depth three and logarithmic bottom fan-in, since it is equivalent to the $\AND_R$ (over all range items $r \in [R]$) of the $\OR_N$ (over all inputs $i \in [N]$) of ``Is input $s_i$ equal to $r$?'', where the quoted question is computed by a conjunction of width $\log R$ over the input bits.

$\SURJ_{R, N}$ has been studied extensively in the contexts of quantum query complexity and approximate degree. Beame and Machmouchi~\cite{beame} showed that computing $\SURJ_{R, N}$ for $R = N/2 + 1$ requires $\tilde{\Omega}(n)$ quantum queries, making it the only known AC$^0$ function with  linear quantum query complexity. Meanwhile,
the $(1/3)$-approximate degree of $\SURJ_{R, N}$ was recently shown to be $\tilde{\Theta}(R^{1/4} \cdot N^{1/2})$. The lower
bound is from our prior work~\cite{bkt}, while the upper bound
was shown by Sherstov~\cite{algopoly}, with a different proof given in~\cite{bkt}. 
In particular, when $R=N/2$, $\adeg_{1/3}(\SURJ_{R, N}) = \tilde{\Theta}(N^{3/4})$. 
Our prior %
work  %
\cite{adegsurj, bkt} built directly on the
%
approximate-degree %
lower bound for $\SURJ_{R, N}$ to 
give near-optimal lower bounds on the $(1/3)$-approximate degree of AC$^0$ (see \expref{Section}{s:hardnessamps} for details).

\medskip \noindent \textbf{Our result.} 
In spite of the progress described above, the  threshold degree $\SURJ_{R, N}$
remained open. For $R < N/2$, an upper bound of $\tilde{O}(\min\{R, N^{1/2}\})$ follows from standard techniques (we prove
this in \expref{Section}{app:ub} for completeness). The best known lower bound was $\Omega(\min\{R, N^{1/3}\})$, obtained by a reduction
to Minsky and Papert's threshold degree lower bound for $\AND_{n^{1/3}} \circ \OR_{n^{2/3}}$.
In this work, we settle the threshold degree of $\SURJ_{R, N}$, showing that the known upper bound is tight up to logarithmic
factors.


\begin{theorem} \label{thm:surj}
	For $R < N/2$, the threshold degree of $\SURJ_{R, N}$ is
	$\tilde{\Theta}(\min\{R, N^{1/2}\})$. In particular, if $R=N^{1/2}$, $\deg_{\pm}(\SURJ_{R, N}) = \tilde{\Theta}(N^{1/2})$. 
\end{theorem}


In addition to  
resolving a natural question in its own right, \expref{Theorem}{thm:surj} matches the best prior threshold degree lower bound for AC$^0$,
previously proved in~\cite{sherstov15} for a much more complicated  function computed by a circuit of strictly greater depth. 
Furthermore, with some extra effort, our lower bound for $\SURJ_{R, N}$ extends to give an $\tilde{\Omega}(n^{1/2})$ lower
bound on the $\UPP^{\text{cc}}$ complexity of a related AC$^0$ function, yielding progress on Open Question 1
%
(see \expref{Section}{s:openquestion1}).   %
In contrast, Sherstov's  $\Omega(n^{1/2})$ threshold degree lower bound
for AC$^0$~\cite{sherstov15} is not known to extend to $\UPP^{\text{cc}}$ complexity. As stated in \expref{Section}{sec:previoussec}, the best previous $\UPP^{\text{cc}}$ lower bound for an AC$^0$ function was $\Omega(n^{2/5})$.

\begin{corollary} \label{cor:signrank}
	There is an AC$^0$
	function $F \colon \bits^{n \times n} \to \bits$ 
	such that $\UPP^{\text{cc}}(F) \geq \tilde{\Omega}(n^{1/2})$.
\end{corollary} 


\subsubsection{\texorpdfstring{AC$^0$ has nearly maximal $\PP^{\text{cc}}$
    complexity}{AC0 has nearly maximal PPcc complexity}}
\label{s:discintro}
In our second result, for any constant $\delta > 0$, we exhibit an AC$^0$ function $f \colon \bits^n \to \bits$ with $\adeg_{\eps}(f) = \Omega(n^{1-\delta})$
for some $\eps= 1-2^{-\Omega(n^{1-\delta})}$. This is a major strengthening
of our %
earlier results \cite{adegsurj, bkt},
which
%
gave  %
a similar result for $\eps=1/3$.
By combining this large-error approximate-degree %
lower bound
with a ``query-to-communication lifting theorem'' for $\PP$~\cite{patmat}, we obtain an $\Omega(n^{1-\delta})$ bound on the $\PP^{\text{cc}}$ complexity
of an AC$^0$ function, nearly resolving Open Question 2 from the previous section. 

\begin{theorem} \label{discthm} \label{thm:discintro}
	For any constant $\delta > 0$, there is an AC$^0$ function $F \colon \bits^{n \times n} \to \bits$ with $\PP^{\text{cc}}(F) = \Omega(n^{1-\delta})$.
\end{theorem}

\label{s:disc}


The best previous lower bound for the $\PP^{\text{cc}}$ complexity of an AC$^0$ function was  $\Omega(n^{1/2})$~\cite{sherstov15}.

%
%

\begin{table}
	{\small
		\centering
		\setlength\tabcolsep{1.7pt}
		\begin{tabular}{|c|c|c|c|c|c|c|}
			\hline
			Reference & $\PP^{\text{dt}}$ & $\PP^{\text{cc}}$   &$\UPP^{\text{dt}}$    &  $\UPP^{\text{cc}}$ & Circuit\\
			& and $\log(\text{threshold weight}$) & $\approx \log(1/\text{discrepancy})$ & $=$ threshold degree  & $\approx \log(\text{sign-rank})$ &	Depth\\
			\hline
			\cite{mp} & --- & --- & $\Omega(n^{1/3})$ & --- & 2\\
			\hline
			\cite{kp}& $\Omega(n^{1/3})$ & ---  & --- & --- & 3 \\
			\hline
			\cite{forster} & --- & $\Omega(\log^k(n))$ & --- & $\Omega(\log^k(n))$ & $O(k)$\\
			\hline
			\cite{osnewbounds} & $\Omega(n^{1/3} \log^k n)$ & --- & $\Omega(n^{1/3} \log^k(n))$ & --- & $O(k)$\\
			\hline
			\cite{sherstovmajmaj}& --- &  $\Omega(n^{1/5})$ & --- & ---  & 3\\
			\hline
			\cite{bvdw, patmat} & --- &  $\Omega(n^{1/3})$  & ---  & --- & 3\\
			\hline
			\cite{razborovsherstov} & --- & --- & --- & $\Omega(n^{1/3})$ & 3\\
			\hline
			\cite{bt14} & $\Omega(n^{2/5})$ & $\Omega(n^{2/5})$ & --- & ---  & 3\\
			\hline
			\cite{sherstov14}&$\Omega(n^{1/2-\delta})$   & $\Omega(n^{1/2-\delta})$ 
			& $\Omega(n^{1/2-\delta})$ & ---& $O(1/\delta)$\\
			\hline
			\cite{sherstov15} & $\Omega(n^{3/7})$ & ---& $\Omega(n^{3/7})$ &--- & 3\\
			\hline
			\cite{sherstov15} & $\Omega(n^{1/2})$  & $\Omega(n^{1/2})$ & $\Omega(n^{1/2})$ &--- & 4\\
			\hline
			\cite{bt16} & --- & --- & --- & $\Omega(n^{2/5})$ & 3\\
			\hline
			\cite{dlweight} & $\Omega(n^{1/2 - \delta})$  & $\Omega(n^{1/2 - \delta})$ & ---  & --- & 3\\
			\hline %
			This paper & $\tilde{\Omega}(n^{1/2})$  &$\tilde{\Omega}(n^{1/2})$  & $\tilde{\Omega}(n^{1/2})$ & --- & 3\\
			\hline
			This paper & ---  & --- & --- & $\tilde{\Omega}(n^{1/2})$ & 7\\
			\hline
			This paper & $\Omega(n^{1 - \delta})$  & $\Omega(n^{1 - \delta})$ & ---  & --- & $O(1/\delta)$\\  %
			\hline
		\end{tabular}
		\caption{Comparison of our new bounds for AC$^0$ to prior work in roughly chronological order. The circuit depth column lists the depth of the Boolean circuit used to exhibit the bound, $\delta$ denotes an arbitrarily small positive constant, and $k$ an arbitrary positive integer. All Boolean circuits are of polynomial size.
		}
		\label{tab:tab}
	}
\end{table}

\section{Algorithmic and complexity-theoretic applications} \label{sec:applications}
\label{s2}
To introduce the applications of our results, we begin by defining the query complexity quantities $\UPP^{\text{dt}}$ and $\PP^{\text{dt}}$ and the communication complexity quantities $\UPP^{\text{cc}}$ and $\PP^{\text{cc}}$.

\medskip \noindent \textbf{Query models.} 
In randomized query complexity, an algorithm aims to evaluate a known Boolean function $f$ on an unknown input $x \in \{-1, 1\}^n$ by reading as few bits of $x$ as possible. We say that the \emph{query cost} of a randomized algorithm is the maximum number of bits it queries for any input $x$.

\begin{itemize}
	\item $\UPP^{\text{dt}}$ considers ``unbounded error'' randomized algorithms, which means that on any input $x$, the algorithm outputs $f(x)$ with probability strictly greater than $1/2$. 
	$\UPP^{\text{dt}}(f)$ is the minimum query cost of any 
	unbounded-error %
	algorithm for $f$.
	\item $\PP^{\text{dt}}(f)$ captures ``large'' (rather than unbounded) error algorithms. If
	a randomized query algorithm outputs $f(x)$ with probability $1/2 + \beta$ for all $x$, then the $\PP$-cost
	of the algorithm 
	is the sum of the query cost and $\log(1/\beta)$. The $\log(1/\beta)$ term here is a proxy for the number of bits of randomness used by the algorithm. Hence,
	the $\PP$-cost of a query algorithm is a lower bound on the runtime, assuming both a query and a random coin toss cost at least one time step.
	$\PP^{\text{dt}}(f)$ is the minimum $\PP$-cost of any randomized query algorithm for $f$.
\end{itemize}

\noindent \textbf{Communication models.} 
$\UPPcc$ and $\PPcc$
consider the standard two-party setup
where Alice holds an input $x$ and Bob holds an input $y$, and they run a private-coin randomized communication protocol to compute a function $f(x, y)$, while minimizing the number of bits they  exchange. In direct analogy to the query complexity measures above,
we say that the \emph{communication cost} of a randomized protocol is the maximum number of bits Alice and Bob exchange on any input $(x, y)$.

\begin{itemize}
	\item $\UPP^{\text{cc}}(f)$~\cite{paturisimon} is the minimum communication cost of any randomized protocol that outputs $f(x, y)$ with probability strictly greater than $1/2$ on all inputs $(x, y)$.
	\vspace{-1mm}
	\item $\PP^{\text{cc}}(f)$~\cite{babaietal} is the minimum $\PP$-cost of a protocol for $f$, where the $\PP$-cost of a protocol that outputs $f(x, y)$ with probability $1/2 + \beta$ for all $(x, y)$ is the sum of the communication cost and $\log(1/\beta)$.
\end{itemize}

We now give an overview of the applications of \expref{Theorem}{thm:discintro} and \expref{Corollary}{cor:signrank}. Further technical background and details are given in \expref{Section}{s:apps}.

\vspace{-3mm}
\subsection{Applications of \expref{Theorem}{thm:discintro}}

$\PPcc$ is known to be equivalent to two measures of central importance in 
learning theory and communication complexity, namely \emph{margin complexity}~\cite{linial2009learning} and \emph{discrepancy}~\cite{klauck2001lower}.
Hence,
\expref{Theorem}{thm:discintro} implies that AC$^0$ has nearly maximal complexity under both measures. Below,
we highlight four additional applications. 

\begin{itemize}
	\item \textbf{Communication Complexity}. 
	The $\PP^{\text{cc}}$ communication model can efficiently simulate 
	almost every two-party communication model, including $\mathbf{P}$ (\ie, deterministic communication), $\mathbf{BPP}$ (randomized communication), $\mathbf{BQP}$ (quantum), and $\mathbf{P}^{\mathbf{NP}}$.
	The only well-studied exceptions are $\UPPcc$, and communication analogs of the polynomial hierarchy (the latter of which we do not know how to prove lower bounds against).
	Hence, in showing that AC$^0$ has essentially maximal $\PP^{\text{cc}}$  complexity, we subsume or nearly subsume 
	all previous results on the communication complexity of AC$^0$.
	\vspace{-1mm}
	\item \textbf{Cryptography}. Bogdanov et al.~\cite{viola} observed that for any
	$f \colon \bits^n \to \bits$ and $d>0$, 
	if one shows that $\adeg_{\eps}(f) \geq d$, then one obtains a scheme
	for sharing a secret bit $b \in \bits$ among $n$ parties
	such that any subset of $d$ shares provides no reconstruction
	advantage, yet applying $f$ to all $n$ shares yields
	$b$ with probability at least $1/2 + \eps/2$.  They combined
	this with known approximate-degree %
	lower bounds for AC$^0$
	functions to get secret sharing schemes with reconstruction procedures in AC$^0$.
	Via this connection, an immediate corollary of \expref{Theorem}{thm:discintro} is a nearly 
	optimal secret sharing scheme in AC$^0$: for any desired constant $\delta > 0$,
	any subset of $n^{1-\delta}$ shares provides no reconstruction
	advantage, yet all $n$ shares can be successfully reconstructed (by applying an AC$^0$ function) with probability
	$1-2^{-n^{1-\delta}}$. 
	\vspace{-1mm}
	\item \textbf{Learning Theory}.
	Valiant~\cite{evolve} introduced the \emph{evolvability} model to quantify how complex functions can emerge through an evolutionary process over realistic population sizes within realistic time
	periods. Feldman~\cite{feldmanevolve} showed that the ``weak evolvability''
	of a class of functions $\mathcal{F} = \{\phi_1, \dots, \phi_{|\mathcal{F}|}\}$
	is characterized by the $\PPcc$ complexity of the function $F(x, y)=\phi_x(y)$. 
	Hence, a consequence of \expref{Theorem}{thm:discintro} 
	is that there are AC$^0$ functions that are nearly maximally
	hard to evolve. That is, for any constant $\delta > 0$,
	there are AC$^0$ functions that require either $2^{n^{1-\delta}}$ generations, or
	populations of size $2^{n^{1-\delta}}$ to evolve,
	even if one only wants evolution to produce a function that has advantage
	just $2^{-n^{1-\delta}}$ over random guessing. 
	See~\cite{feldmanevolve} for additional details.
	
	\vspace{-1mm}
	\item \textbf{Circuit Complexity}.  If $\PPcc(f)\geq d$, then $f$ is not computable by Majority-of-Threshold circuits of size $2^{\Omega(d)}$~\cite{nisan}.
	Hence, by showing that AC$^0$ has nearly maximal $\PPcc$ complexity, we show
	that there are AC$^0$ functions that are not computed by Majority-of-Threshold circuits of size $2^{n^{1-\delta}}$.
	That is, AC$^0$ has essentially no non-trivial simulation by Majority-of-Threshold circuits (in contrast,
AC$^0$ can be efficiently simulated by depth-three Majority circuits~\cite{allender}).
	
	
	
\end{itemize}


\subsection{Applications of \expref{Corollary}{cor:signrank}} 

As indicated in \expref{Section}{s1}, $\UPPcc(F)$ is known to be characterized by (the logarithm of)
of the \emph{sign-rank} of the matrix $[F(x, y)]_{x, y \in \bits^{n \times n}}$~\cite{paturisimon}.\footnote{The sign-rank of a matrix $M$
	with entries in $\{\pm1\}$
	is the least rank of a real matrix $M'$ that agrees in sign with $M$ entry-wise.}
Hence,  \expref{Corollary}{cor:signrank} implies an $\exp(\tilde{\Omega}(n^{1/2}))$ lower bound on the sign-rank of AC$^0$ function.
Below, we highlight two additional applications of  \expref{Corollary}{cor:signrank}, based on the following connections
between communication complexity, circuit complexity, and learning theory. 

In communication complexity, $\UPP^{\text{cc}}$ is the most powerful two-party model against which we know how to prove lower bounds. 
In circuit complexity, if $\UPP^{\text{cc}}(f) \geq d$, then $f$ cannot be computed
by Threshold-of-Majority circuits of size $2^{\Omega(d)}$~\cite{forstertm}. (Threshold-of-Majority circuits represent the most powerful class of threshold circuits against which we can
prove superpolynomial lower bounds.)
In learning theory, it is commonly assumed that data can be classified by a halfspace in many dimensions; the $\UPP^{\text{cc}}$-complexity of a concept class precisely captures how many dimensions are needed.
To connect this to a previously mentioned example, Klivans and Servedio~\cite{ksdnf} observed that an upper bound of $d$ on the $\UPPcc$
complexity of a concept class $\mathcal{C}$ yields a PAC learning for $\mathcal{C}$ running in time $2^{O(d)}$. They 
used this result to give a $2^{\tilde{O}(n^{1/3})}$-time algorithm for PAC-learning CNFs. This remains the state-of-the-art algorithm for this fundamental problem.
Accordingly, \expref{Corollary}{cor:signrank} has the following implications.

\begin{itemize}
	\item \textbf{Circuit Complexity.} There are AC$^0$ functions that are not computable by Threshold-of-Majority Circuits of size $2^{\tilde{\Omega}(n^{1/2})}$.
	
	\item \textbf{Learning Theory.}  $\UPP^{\text{cc}}$-based learning algorithms cannot learn AC$^0$
	in time better than $2^{\tilde{\Omega}(n^{1/2})}$. 
\end{itemize}



\vspace{-4mm}
\section{Techniques} 
\subsection{The \texorpdfstring{$\mathsf{SURJECTIVITY}$}{SURJECTIVITY}
     Lower Bound}
\label{s:surjoverview}
For a function $f_n$, let $f^{\leq N}$ denote the partial function obtained by
restricting $f$ to the domain of inputs of Hamming weight at most $N$. 
The $\eps$-approximate degree of $f^{\leq N}$, denoted $\adeg_{\eps}(f^{\leq N})$, is the least degree of a 
real polynomial $p$ such that 
\begin{equation} \label{dupy} |p(x)-f(x)| \leq \eps \text{ for all inputs } x \text{ of Hamming weight at most } N. \end{equation}
Note that \expeqref{Property}{dupy} allows $p$ to \emph{behave arbitrarily} on inputs $x$ of Hamming weight more than $N$.
Similarly, the threshold degree of $f^{\leq N}$ is the least degree of a 
real polynomial $p$ such that 
\begin{equation*} p(x) \cdot f(x) > 0 \text{ for all inputs } x \text{ of Hamming weight at most } N. \end{equation*}
Our prior work~\cite{adegsurj} showed that the $\eps$-approximate (respectively, threshold) degree of $\SURJ_{R, N}$
is \emph{equivalent} to the $\eps$-approximate (respectively, threshold) degree of $(\AND_R \circ \OR_N)^{\leq N}$.
Hence, the main technical result underpinning our 
threshold-degree %
lower bound for $\SURJ$
is the following theorem about the threshold degree of $(\AND_R \circ \OR_N)^{\leq N}$ (we have made no effort to optimize the logarithmic factors).

\begin{restatable}{theorem}{surjthmintro}
	\label{surjthmintro} \label{surjthm}
	Let $R=N^{1/2}$. Then $\deg_{\pm}\left(\left(\AND_R \circ \OR_N\right)^{\leq N}\right) = \Omega(N^{1/2}/\log^{3/2} N)$.
\end{restatable}

\noindent \textbf{Discussion.}
\expref{Theorem}{surjthmintro} is a substantial strengthening of the classic result
of Minsky and Papert~\cite{mp} mentioned above, %
which established that the total function $\mathsf{MP}_{N^{1/2}, N} \coloneqq  \AND_{N^{1/2}} \circ \OR_N$ on $n=N^{3/2}$
inputs has threshold
degree $\Omega(N^{1/2})$. %
\expref{Theorem}{surjthmintro} establishes that Minsky and Papert's lower bound holds even under the promise that the $n$-bit input to $\mathsf{MP}_{N^{1/2}, N}$
has Hamming weight at most $N=n^{2/3}$. That is, any polynomial that sign-represents $\AND_{n^{1/3}} \circ \OR_{n^{2/3}}$ on inputs
of Hamming weight at most $n^{2/3}$
has degree $\tilde{\Omega}(n^{1/3})$, \emph{even when $p$ is allowed to behave arbitrarily on inputs of Hamming weight larger than $n^{2/3}$}. 

\medskip 
\noindent \textbf{Proof overview for \expref{Theorem}{surjthmintro} and comparison to prior work.}
Like much recent work on %
%
approxi\-mate-degree and threshold-degree %
 lower bounds, our proof makes 
use of \emph{dual polynomials}. A dual polynomial is a dual solution to a certain linear program capturing the approximate or threshold degree of any function,
and acts as a certificate of the high approximate or threshold degree of the function.

A dual polynomial that witnesses the fact that $\deg_{\pm}(f_M) \geq d$ is a 
function $\psi \colon \bits^M \to \bits$ satisfying three 
%
conditions:%
\begin{itemize}
	\item $\psi(x) \cdot f(x) \geq 0$ for all $x \in \bits^M$. If $\psi$ satisfies this condition, we say $\psi$
	agrees in sign with $f$. 
	\item  $\sum_{x \in \bits^M} |\psi(x)| = 1$. If $\psi$ satisfies this condition, it is said to have $\ell_1$-norm equal to 1.
	\item For all polynomials $p \colon \bits^M \to \R$ of degree at most $d$, $\sum_{x \in \bits^M} p(x) \cdot \psi(x) = 0$. 
	If $\psi$ satisfies this condition, it is said to have \emph{pure high degree} at least $d$.
\end{itemize}
A dual witness for the fact that $\adeg_{\eps}(f_M) \geq d$ is similar, except that the first condition is replaced
with:
\begin{itemize}
	\item $\sum_{x \in \bits^M}\psi(x) \cdot f(x) > \eps$. If $\psi$ satisfies this condition, it is said to be \emph{$\eps$-correlated} with $f$.
	If $\psi(x) \cdot f(x) < 0$, we say that $\psi$ \emph{makes an error} at $x$.
\end{itemize}

Sherstov~\cite{sherstov15} reproved Minsky and Papert's result by constructing an explicit dual witness for $\mathsf{MP}_{N^{1/2}, N}$, via
a two-step process. First, Sherstov started with a dual witness $\psi_{\textsf{base}}$ for the fact that 
$$\adeg_{\eps}(\mathsf{MP}_{N^{1/2}, N}) =\Omega(N^{1/2}), \text{ for } \eps=1-2^{-N^{1/2}}.$$
The function $\psi_{\textsf{base}}$ was introduced
in our prior work~\cite{bt14}, where it was constructed by combining a dual witness for $\AND_{N^{1/2}}$ with
a dual witness for $\OR_{N}$ via a technique called dual block composition~\cite{shizhu, lee, sherstovhalfspaces1} (see \expref{Section}{s:dualblockcomposition} for details of this 
very important technique for combining dual witnesses). 

Unfortunately, $\psi_{\textsf{base}}$ falls short of witnessing Minsky and Papert's threshold-degree %
lower bound
because it makes errors on some inputs. In the second step of Sherstov's construction~\cite{sherstov15}, he adds in a correction term that zeros out the errors of $\psi_{\textsf{base}}$, without disturbing
the sign of $\psi_{\textsf{base}}$ on any other inputs, and without lowering its pure high degree.

\expref{Theorem}{surjthm} asserts that $\mathsf{MP}_{N^{1/2}, N}^{\leq N}$ satisfies the same threshold-degree %
lower bound as 
$\mathsf{MP}_{N^{1/2}, N}$ itself. To prove \expref{Theorem}{surjthm}, we need to construct a dual witness $\psi$ that not only reproves Minsky and Papert's
classic lower bound for $\mathsf{MP}_{N^{1/2}, N}$, but also satisfies the extra condition that:
\begin{align} \label{howhowhow} \psi(x)=0 \text{ for all inputs } x \text{ of Hamming weight more than } N\,.\end{align}
(Here, the Hamming weight of $x$ is the number of $-1$ entries.) To accomplish this, we apply a novel strategy that can be thought
of as a three-step process. First, like Sherstov, we start with $\psi_{\textsf{base}}$. Second, we modify $\psi_{\textsf{base}}$
to obtain a dual witness $\psi'_{\textsf{base}}$ that places significant mass on all inputs of Hamming weight at most $d$, for some $d=\tilde{\Omega}(N^{1/2})$
(details of the construction of $\psi'_{\textsf{base}}$ are described two paragraphs hence).
More specifically, we ensure that $\psi'_{\textsf{base}}$  satisfies:
%
\begin{align} \label{noreallythispaperisntgood} |\psi'_{\textsf{base}}(x)| \gg n^{-d}  \text{ for all inputs } x \text{ of Hamming weight at most } d\,. \end{align}
%
We refer to this property by saying that $\psi'_{\textsf{base}}$ is ``smooth'' or ``large'' on all inputs of Hamming weight at most $d$.
Note that, in modifying $\psi_{\textsf{base}}$ to obtain $\psi'_{\textsf{base}}$, we do \emph{not} correct the errors that $\psi_{\textsf{base}}$ makes, nor do we ensure that $\psi'_{\textsf{base}}$
is supported on inputs of Hamming weight at most $N$.

Third, we add in a correction term, very different than Sherstov's correction term, that not only zeros out the errors of $\psi'_{\textsf{base}}$,
but also zeros out any mass it places on inputs of Hamming weight more than $N$. While the general technique we use to construct
this correction term appeared in our prior %
work \cite{adegsurj, bkt}, the novelty in our construction and analysis is two-fold.
First, the technique was used in our prior work only to zero out mass placed on inputs of Hamming weight more than $N$ (\ie, to ensure that
\expeqref{Equation}{howhowhow} is satisfied), not to correct errors. Second, and more importantly, 
we crucially exploit the largeness of $\psi'_{\textsf{base}}$ on inputs of Hamming weight at most $d$ to
ensure that the correction term does not disturb the sign of $\psi'_{\textsf{base}}$ on any inputs other 
than those on which it is deliberately being zeroed out. This is what enables us 
to obtain a %
threshold-degree  %
lower bound, whereas our
%
in our prior work we %
were only able to obtain %
$\eps$-approximate-degree %
lower bounds for $\eps$ bounded away from 1. 

\medskip
Our ``smoothing followed by correction'' approach appears to be significantly more generic than the correction technique of~\cite{sherstov15}. For example, prior work of Bouland et al.~\cite{bchtv} proved an $\Omega(n^{1/4})$ lower bound on the threshold degree of a certain function denoted $\mathsf{GAPMAJ}_{n^{1/4}} \circ \mathsf{PTP}_{n^{3/4}}$, and used this result to give an oracle separating the complexity classes $\textsf{SZK}$ and $\textsf{UPP}$, thereby answering
an open question of Watrous from 2002. Our techniques can be used to give a 
much simpler proof of this result of this %
lower bound on threshold degree. %
We expect that our technique will find additional applications in the future. 

\medskip \noindent \textbf{Details of the smoothing step.}
As stated above, the dual witness $\psi_{\textsf{base}}$ from our prior work
does not %
have %
the property we need
%
(see \expeqref{Equation}{noreallythispaperisntgood})  %
of being ``large'' on all inputs of Hamming weight at most
$d=\tilde{\Omega}(N^{1/2})$. 

Fortunately, we observe that although $\psi_{\textsf{base}}$ is \emph{not} large on all inputs of Hamming weight at most $d$, it
\emph{is} large on one very special input of low Hamming weight, namely the $\textsf{ALL-FALSE}$ input. 
That is, $\psi_{\textsf{base}}(\mathbf{1}) \geq 2^{-d}$. 
So we just need a way to ``bootstrap'' this largeness property on $\mathbf{1}$ to 
a largeness property on all inputs of Hamming weight at most $d$.
Put another way, we need to be able to treat other inputs of Hamming weight at most $d$ as if they actually have Hamming weight 0.
But
$\mathsf{MP}_{N^{1/2}, N} \coloneqq  \AND_{N^{1/2}} \circ \OR_N$ has a property that enables precisely this: we can fix the inputs to any constant fraction $c$ of the $\OR$ gates to an arbitrary value in $\OR^{-1}(-1)$,
and the remaining function of the unrestricted inputs is $\AND_{(1-c) \cdot R} \circ \OR_N$. This  is ``almost'' the same
function as $\AND_R \circ \OR_N$; we have merely
slightly reduced the top fan-in, which does not substantially lower the threshold degree of the resulting function. 

We exploit the above observation
to achieve the following: for each input $x$ of Hamming weight at most $d$, we build a dual witness $\nu_x$
targeted at $x$ (\ie, that essentially treats $x$ as if it is the $\textsf{ALL-FALSE}$ input). We do this as follows. Let $T$ be the set of all $\OR$ gates
that are
fed one or more $-1$s by $x$, and let $S \subseteq [N^{1/2} \cdot N]$ be the union of the inputs
to each of the $\OR$ gates in $T$. Let $\psi_{\textsf{base}}$ be the dual witness for $\AND_{N^{1/2} - |T|} \circ \OR_N$ given in our
%
earlier paper~\cite{bt14}.  %
We let 
\begin{align*} \nu_x(y)=\begin{cases} \psi_{\textsf{base}}(y_{\bar{S}}) & \text{ if } y_{S} = x_S\\
0 & \text{ otherwise,} \end{cases}\end{align*} 
%
where $y_{\bar{S}}$ denotes the set of all the coordinates of $y$ other than those in $S$.

The dual witness $\psi'_{\textsf{base}}$ is then defined to be the average of the $\nu_x$'s, over all inputs $x$ of Hamming weight at most $d$. 
This averaged dual witness $\psi'_{\textsf{base}}$ has all of the same useful properties as $\psi_{\textsf{base}}$, and additionally satisfies the key requirement captured by \expeqref{Equation}{noreallythispaperisntgood}.

\subsection{Extension to \texorpdfstring{$\UPPcc$}{UPPcc}: 
Proof of \expref{Corollary}{cor:signrank}}
\label{igivethefrigup}
Building on the celebrated framework of Forster~\cite{forster}, Razborov and Sherstov~\cite{razborovsherstov} developed
techniques to translate threshold-degree %
lower bounds into sign-rank lower bounds. Specifically, they showed that, in order for a threshold-degree %
lower bound of the form $\deg_{\pm}(f_n) \geq d$ to translate into a $\UPPcc$ lower bound for a related function $F$,
it suffices for the threshold-degree %
lower bound for $f_n$ to be exhibited by a dual witness $\phi$
satisfying the following smoothness condition:
\begin{align} |\phi(x)| \geq 2^{-O(d)} \cdot 2^{-n} \text{ for all but a } 2^{-\Omega(d)} \text{ fraction of inputs } x \in \bits^n\,. \label{rscondition}\end{align} 
Note
that this is a different smoothness condition than the one satisfied by the dual witness $\psi'_{\textsf{base}}$ discussed above for $\mathsf{MP}_{N^{1/2}, N}$
%
(\expeqref{Equation}{noreallythispaperisntgood}): %
on inputs $x$ of Hamming weight at most $d$,
$|\psi'_{\textsf{base}}(x)|$ is always at least $n^{-d} \gg 2^{-d} \cdot 2^{-n}$, whereas on inputs $x$ of Hamming weight more than $d$, $|\psi'_{\textsf{base}}(x)|$ may be 0.
In words, $|\psi'_{\textsf{base}}(x)|$ is \emph{very large} on inputs $x$ of Hamming weight at most $d$, but may not be large at all on inputs of larger Hamming weight.
In contrast, \expeqref{Equation}{rscondition} requires a dual witness to be ``somewhat large'' (within a $2^{-O(d)}$ factor of uniform) on \emph{nearly all} inputs. 


In summary, our construction of a dual witness 
for $\mathsf{MP}^{\leq N}_{N^{1/2}, N}$ that is sketched in the previous subsection is not sufficient to apply
Razborov and Sherstov's framework
to $\SURJ_{R, N}$, for two reasons.
First, the dual witness we construct for  $\mathsf{MP}_{N^{1/2}, N}^{\leq N}$ 
is not smooth in the sense of \expeqref{Equation}{rscondition}, as it is only ``large'' on inputs of Hamming weight at most $d$.
Second, 
to apply Razborov and Sherstov's framework to $\SURJ_{R, N}$, we actually need to give a smooth dual witness for $\SURJ_{R, N}$ itself, not for $\mathsf{MP}^{\leq N}_{N^{1/2}, N}$. Note that
$\SURJ_{R, N}$ is defined over the domain $\bits^n$ where $n=N \log R$, while $\mathsf{MP}_{N^{1/2}, N}^{\leq N}$ is defined over subset of $\bits^{NR}$ 
consisting of inputs of Hamming weight at most $N$.

We address both of the above issues as follows. First, we show how to turn our dual witness $\mu$ for $\mathsf{MP}_{N^{1/2}, N}^{\leq N}$
into a dual witness $\hat{\sigma}$
for the fact that $\deg_{\pm}\left(\SURJ_{R, N}\right)\geq d$, such that $\hat{\sigma}$ inherits the ``largeness'' property of $\mu$ on inputs of Hamming weight at most $d$. 
Second, we transform $\hat{\sigma}$ into a dual witness $\tau$ for the fact that $\deg_{\pm}\left(\SURJ_{R, N} \circ \AND_{\log^2 n} \circ \PARITY_{\log^3 n}\right)\geq d$, such
that $\tau$ satisfies the smoothness condition given in \expeqref{Equation}{rscondition}. 
We conclude that $\SURJ_{R, N} \circ \AND_{\log^2 n} \circ \PARITY_{\log^3 n}$ can be transformed into
a related function $F$ (on $\tilde{O}(n)$ inputs, and which is also in AC$^0$) that has sign-rank $\exp(\tilde{\Omega}(n^{1/2}))$.

\label{precedingsection}
\subsection{The \texorpdfstring{$\PPcc$}{PPcc} bound: 
Proof of \expref{Theorem}{thm:discintro}}
\label{s:hardnessamps}
As mentioned in \expref{Section}{s:disc}, the core of \expref{Theorem}{thm:discintro} is to exhibit an
AC$^0$ function $f$ such that $\adeg_{\eps}(f) = \Omega(n^{1-\delta})$
for some $\eps= 1-2^{-\Omega(n^{1-\delta})}$. To accomplish this, we prove a hardness
amplification theorem that should be understood in the context of a weaker result from our prior work~\cite{adegsurj}.

As stated in \expref{Section}{s:surjoverview}, for $\eps=1/3$, our prior work~\cite{adegsurj} showed how to take any Boolean function $f_n$ in AC$^0$ 
with $\eps$-approximate degree $d$ and transform it into a related function $g$ on roughly the same number of variables, such that $g$ is still in AC$^0$, and $g$ has significantly higher $\eps'$-approximate degree for some $\eps' \approx 1/3$. %
This was done in a two-step process. First, we showed
that in order to construct a ``harder'' function $g$, it is sufficient to identify an AC$^0$ function $G$ defined on $\poly(n)$ inputs
such that for some $\ell=n\cdot\polylog(n)$, $\adeg_{\eps'}(G^{\leq \ell}) \gg d$.
\footnote{This step was also
	used in the analysis of $\SURJ_{R, N}$ outlined in \expref{Section}{precedingsection} above, where $G$ was the function $\AND_R \circ \OR_N$.} 
Second, we exhibited such a $G$. 
In our prior %
work~\cite{adegsurj, bkt}, for general functions $f_n$, 
the function $G$ was $f_n \circ \AND_r \circ \OR_{m'}$, where $r=10 \log n$, and $m'=\Theta(n/d)$. 

We would like to prove a similar result, but we require that $G$ have larger $\eps'$-approximate degree than $f_n$, \emph{where $\eps'$ is exponentially closer to 1
	than is $\eps$ itself}. 
Unfortunately, the definition of $G$ from our
%
earlier papers  %
\cite{adegsurj, bkt} does not necessarily result in such a function. For example, if $f_n = \OR_n$
(or any polylogarithmic DNF for that matter), then the function $G=f_n \circ \AND_r \circ \OR_{m'}$ is also a DNF of polylogarithmic width, and 
it is not hard to see that all such DNFs have $\eps$-approximate degree at most $\polylog(n)$ for some $\eps=1-1/n^{\polylog(n)}$. 

To address this situation, we change the definition of $G$. Rather than defining  $G\coloneqq f_n \circ \AND_r \circ \OR_{m'}$, we
define $G=\GAPMAJ_t  \circ f_z \circ \AND_r \circ \OR_{m}$ for appropriately chosen settings of the parameters $t$, $z$, $r$, and $m$. (For simplicity of exposition, some of these parameters are described differently here than in the proof in \expref{Section}{s:discproof}.)
Here, $\GAPMAJ_t$ denotes any function evaluating to 1 on inputs of Hamming weight at most $t/3$, $-1$ on inputs of Hamming weight
at least $2t/3$, and taking any value in $\{-1, 1\}$ on all other inputs (such functions are also called \emph{approximate majorities}, and 
it is known that there are approximate majorities computable in AC$^0$). $\GAPMAJ$ has also played an important role
in related prior work \cite{bchtv, adegsurj}.

In order to show that $\adeg_{\eps'}(G^{\leq \ell}) \gg \adeg_{\eps}(f_n)$ for an $\eps'$ that is exponentially closer to 1 than is $\eps$,
we require a more delicate construction of a dual witness than our
%
%
papers \cite{adegsurj, bkt}.  %
After all, %
those papers   %
 only required a dual witness for $G^{\leq \ell}$ with correlation at least $1/3$ with $G^{\ell}$, while we require
a dual witness achieving correlation with $G^{\leq \ell}$ that is exponentially close to 1. 
Roughly speaking, %
in \cite{adegsurj, bkt} we %
%
were able to get away
with exclusively using the simple and clean technique called dual block composition (described in \expref{Section}{s:dualblockcomposition}) for constructing
dual witnesses for block composed functions by combining dual witnesses for the constituent functions. %
In this %
article,  %
we instead use a closely related but more involved
construction introduced by Sherstov~\cite{sherstovdp}. (Sherstov introduced his construction to prove that approximate degree satisfies a type of direct-sum theorem.)

More specifically, suppose that for some positive integer $k$, $f_z$ has $\eps(z)$-approximate degree at least $d(z)=z^{k/(k+1)}$, where $\eps(z)=1-2^{-z^{k/(k+1)}}$. 
In our definition of $G$, we set intermediate parameters $t=n^{1/(k+2)}$, $z=n^{(k+1)/(k+2)}$, $r=10 \log n$, and $m=n^{2/(k+2)}$, and we build a dual witness
for $G^{\leq \ell}$ via a multi-step construction.

In Step 1, we take dual witnesses $\psi_{f_z}$, $\psi_{\AND_r}$, and $\psi_{\OR_m}$ for  $f_z$, $\AND_r$, and $\OR_{m}$
respectively, and we combine them 
using the technique of Sherstov~\cite{sherstovdp}, to give a dual witness $\gamma$ for $f_z \circ \AND_r \circ \OR_{m}$ satisfying
the following %
conditions: %
$\gamma$ has pure high degree 
at least $D(n)=n^{(k+1)/(k+2)} = d(n)^{(k+1)/k} \gg d(n)$, and $\gamma$'s correlation with  with $f_z \circ \AND_r \circ \OR_{m}$ is $\eps''\approx\eps(z)$. 
That is, $\gamma$ witnesses the fact that the $\eps''$-approximate degree of $f_z \circ \AND_r \circ \OR_{m}$ is much larger than 
the $\eps(n)$-approximate degree of $f_n$ itself. 

%
%
This step of the construction is in contrast to our prior work, which constructed a dual witness for $f_n \circ \AND_r \circ \OR_{m'}$
via direct dual block composition of $\psi_{f_n}$, $\psi_{\AND_r}$, and $\psi_{\OR_m}$. Direct dual block composition does
not suffice for us because it
would yield a dual witness with significantly worse correlation with $f_z \circ \AND_r \circ \OR_{m}$ than $\eps(z)$.

While achieving correlation $\eps'' \approx \eps(z)$ is an improvement over what would obtain
from direct dual block composition, it is still significantly farther from $1$ than is $\eps(n)$, \ie, $1-\eps'' \gg 1-\eps(n)$. And we ultimately
need to construct a dual witness for $G^{\leq \ell}$ that is significantly \emph{closer} to 1 than is $\eps(n)$.
To address this issue,
in Step 2 of our construction, we use dual block composition to turn $\gamma$ into a dual witness $\eta$ for $G=\GAPMAJ_t  \circ f_z \circ \AND_r \circ \OR_{m}$
satisfying the following %
conditions: %
$\eta$ has  the same pure high degree as $\gamma$, and moreover $\eta$ has
correlation at least $\eps'=1-2^{-\Omega(n^{(k+1)/(k+2)})}$ with $G$. 

However, after Step 2, we are still not done, because $\eta$ places some mass on
inputs of Hamming weight as large as $t \cdot z \cdot r \cdot m \gg \ell$. Hence $\eta$ is only a dual
witness to the high $\eps'$-approximate degree of $G$, not the high $\eps'$-approximate degree of $G^{\leq \ell}$ (recall that any
dual witness %
for $G^{\leq \ell}$, must evaluate to 0 on all inputs of Hamming weight larger than $\ell$, %
as described in %
\expeqref{Equation}{howhowhow}). 
Nonetheless, as in our prior work \cite{adegsurj, bkt}, we are able to argue that $\eta$ places \emph{very little} mass on inputs of Hamming weight more than
$\ell$, and thereby invoke techniques from %
these papers %
to zero out this mass. 
The reason this final step of the argument is not immediate
from %
\cite{adegsurj, bkt} is as follows. Although %
these papers %
developed
a precise understanding of how much mass is placed on inputs of Hamming weight
more than $\ell$ by dual witnesses constructed via basic dual block composition, 
the dual witness $\gamma$ for $f_z \circ \AND_r \circ \OR_{m}$ that we constructed in Step 1 
was \emph{not} built by invoking pure dual block composition. Our key observation is that Sherstov's technique that
we invoked to construct $\gamma$ is ``similar enough'' to vanilla dual block composition that the precise understanding
of dual block composition developed in our prior work can be brought to bear on our dual witness $\eta$.

In summary, there are two main technical contributions in our proof of \expref{Theorem}{thm:discintro}. The first is the identification
of a hardness amplification construction for $\eps$-approximate degree that not only amplifies the degree against
which the lower bound holds, but also the error parameter $\eps$. The second is constructing a dual polynomial
to witness the claimed lower bound, using techniques more involved and delicate than the vanilla dual block composition technique that
sufficed in %
\cite{adegsurj, bkt}.



\vspace{-2mm}
\section{Subsequent work and discussion}

Subsequent to our work, Sherstov and Wu~\cite{SherstovWu19} have made major progress toward resolving Open Problem 1 by showing nearly optimal threshold degree and sign-rank lower bounds for AC$^0$. Specifically, for every $k \ge 1$, they exhibit a family of depth-$k$ AC$^0$ circuits with threshold degree $\tilde{\Omega}(n^{(k-1)/(k+1)})$. This generalizes Minsky and Papert's lower bound of $\Omega(n^{1/3})$ on the threshold degree of DNF, as well as our lower bound of $\tilde{\Omega}(n^{1/2})$ for the depth-3 $\mathsf{SURJECTIVITY}$ function. Sherstov and Wu, moreover, show that for any positive constant $\delta > 0$ there is a family of AC$^0$ circuits with depth $O(1/\delta)$ and sign-rank $\exp\left(\tilde{\Omega}(n^{1-\delta})\right)$. This gives an almost optimal improvement to our sign-rank lower bound of $\exp\left(\tilde{\Omega}(n^{1/2})\right)$ on an AC$^0$ function.

As in our proof of \expref{Theorem}{thm:discintro} (see \expref{Theorem}{hardnessamp} in \expref{Section}{s:discproof}), as well as our %
earlier paper~\cite{adegsurj},
Sherstov and Wu obtain their %
lower bound on the threshold degree of %
 AC$^0$ by recursively applying a new hardness amplification theorem. Their hardness amplification theorem shows how to convert a function $f_z$ into a new function $g_n$, computable by circuits with slightly higher depth and roughly the same size, but with polynomially larger threshold degree. Again as in the proof of \expref{Theorem}{thm:discintro}, in order to obtain such a $g$, it suffices to construct a function $G$ with $\deg_{\pm}(G^{\le n}) \gg \deg_{\pm}(f)$. Starting from a function $f_z$ with threshold degree $z^{(k-1)/(k+1)}$, the function $G$ that they identify as sufficient for this purpose is $G = f_z \circ \mathsf{MP}_{r, r^2}$, where $z = n^{(k+1)/(k+3)}$ and $r = n^{2/(k+3)}$. When $f_z$ is a trivial function, this recovers our lower bound of $\tilde{\Omega}(n^{1/2})$ for $\mathsf{SURJECTIVITY}$. Hence, their construction in full can be viewed as a generalization of our \expref{Theorem}{surjthm} that is amenable to recursive application. 
This requires several technical new ideas in the construction of the dual witness. However, we remain optimistic that the simplicity of our analysis for 
$\mathsf{SURJECTIVITY}$ will nonetheless lead to future applications of our techniques.

Sherstov and Wu's sign-rank lower bound follows from a similar high-level (though more technically demanding) strategy, where they show that \emph{smooth} threshold degree also obeys such a hardness amplification theorem.

While these new results resolve the most glaring question raised in the initial version of this %
paper, %
a number of interesting directions remain for further study. A common feature of our large-error %
approximate-degree %
lower bound and Sherstov and Wu's %
threshold-degree %
and sign-rank lower bounds for AC$^0$ is that, in order to obtain lower bounds of the form $\Omega(n^{1-\delta})$, we must consider functions computed by circuits of depth $\Theta(1/\delta)$. This contrasts with the situation for %
bounded-error %
approximate degree~\cite{adegsurj}, where a lower bound of $\Omega(n^{1-\delta})$ can be obtained at depth only $O(\log(1/\delta))$. Can one show that there are AC$^0$ functions $f$ of depth $O(\log(1/\delta))$ with $\adeg_{\eps}(f) = \Omega(n^{1-\delta})$ for $\eps = 1 - 2^{-\Omega(n^{1-\delta})}$ or with $\deg_{\pm}(f) = \Omega(n^{1-\delta})$?  There is a common underlying reason why our construction and Sherstov and Wu's construction both require circuits of depth $\Theta(1/\delta)$ and not $\Theta(\log(1/\delta))$:  a component of the hardness amplifier in both constructions (in our case, $\GAPMAJ_{n^{1/(k+1)}}$,
and in Sherstov and Wu's case, the top gate of $\mathsf{MP}_{r, r^2}$) is used to amplify error but does not amplify degree. In contrast, in the construction of 
~\cite{adegsurj} for lower bounding bounded-error approximate degree, up to a logarithmic factor, all of the hardness amplifier is used to amplify degree. 


\medskip

We would also like to highlight the question of proving sublinear \emph{upper bounds} on the threshold degree of AC$^0$. 
Given the surprising $O(R^{1/4} \cdot N^{1/2})$ upper bound on the $(1/3)$-approximate degree of $\SURJ_{R, N}$ from
recent %
work %
\cite{algopoly, bkt}, we have begun to seriously entertain the possibility that for every function $f$ computable by AC$^0$ of depth $k$, there is some constant $\delta(k)>0$ such that the threshold degree (and possibly even $(1/3)$-approximate degree) of $f$ is $O(n^{1-\delta})$. 
Unfortunately, we cannot currently even show that this is true for 
depth-three %
circuits of quadratic size. Any progress in this direction would be very interesting,
and we believe that such progress would likely lead to new circuit lower bounds.


\medskip \noindent \textbf{Outline for the rest of the paper.}
\expref{Section}{sec:prelims} covers technical preliminaries. \expref{Section}{s:surj}
contains the proof of our tight threshold-degree %
lower bound for $\mathsf{SURJECTIVITY}$ (\expref{Theorem}{thm:surj})
and its extension to a sign-rank lower bound for a related function in AC$^0$ (\expref{Corollary}{cor:signrank}). 
\expref{Section}{s:discproof} proves our nearly optimal bound on the discrepancy of AC$^0$ (\expref{Theorem}{thm:discintro}).
\expref{Section}{s:apps} elaborates on applications of these results in communication complexity, circuit complexity, learning theory, and cryptography.

\section{Preliminaries}
\label{sec:prelims}
\subsection{Notation}
For a natural number $N$, let $[N] = \{1, 2, \dots, N\}$ and $[N]_0 = \{0, 1, 2, \dots, N\}$. All logarithms in this paper are taken in base $2$ unless otherwise noted via
the notation $\ln$, which refers to base $e$. 
For $x \in \R$, define $\sgn(x) = -1$ if $x < 0$ and 1 otherwise. For any set $S$, the notation $x \sim S$ means
that $x$ is drawn uniformly at random from $S$. 

As mentioned in \expref{Section}{s:footyfootfoot}, we sometimes use subscripts to indicate the number of variables on which a function is defined. 
For example, we denote $\OR \colon \bits^n \to \bits$ by $\OR_n$.

For any input $x \in \bits^n$, $|x|=|\{i \colon x_i=-1\}|$ denotes the Hamming weight of $x$.
$\mathcal{H}_n$ denotes the set $\bits^n$, and 
for any $d \leq n$, we define $\mathcal{H}_n^{\leq d} \coloneqq  \{x \in \bits^n \colon |x| \leq d\}$.
Similarly,  $\mathcal{H}_n^{> d} \coloneqq  \{x \in \bits^n \colon |x| > d\}$.
Given a Boolean function $f_n$, we denote by $f_n^{\leq d}$ the \emph{partial} function obtained
by restricting the domain of $f_n$ to $\mathcal{H}_n^{\leq d}$.




\subsection{Threshold degree and its dual formulation}

\begin{definition} \label{def:tdeg}
	Let $D \subseteq \{-1, 1\}^n$, and let $f$ be a function mapping $D$ to $\{-1,1\}$.
	The \emph{threshold degree} of $f$, denoted $\deg_{\pm}(f)$, is the least degree of a real polynomial $p \colon \bits^n \to \R$ such that $p(x) \cdot f(x) > 0 $ when $x \in \mathcal{D}$.
	No constraint is placed on $p(x)$ for any $x \in (\bits^n \setminus D)$. 	
\end{definition}


In this paper, we make essential use of a (standard) dual formulation of threshold degree. To describe
this dual formulation, we need to introduce some terminology.

\begin{definition}
	Let $\psi \colon \bits^n \to \R$ be any real-valued function on the Boolean hypercube. 
	Given another function $p \colon \bits^n \to \R$, we let $\langle \psi, p\rangle \coloneqq  \sum_{x \in \bits^n} \psi(x) \cdot p(x)$,
	and refer to $\langle \psi, p\rangle$ as the \emph{correlation} of $\psi$ with $p$. 
	If $\langle \psi, p\rangle = 0$ for all polynomials $p$ of degree at most $d$, we say that $\psi$ has \emph{pure high degree} at least $d$.
	We let $\|\psi\|_1 \coloneqq  \sum_{x \in \bits^n} |\psi(x)|$, and refer to $\|\psi\|_1$ as the $\ell_1$-norm of $\psi$.
\end{definition}


The following theorem provides the aforementioned dual formulation of threshold degree.
\begin{theorem}\label{theorem:tdegdual}
	
	Let $f:D \to \{-1,1\}$ with $D \subseteq \{-1,1\}^n$. Then $\deg_{\pm}(f) > d$ if and only if there is a real function $\psi: \bits^n \to \R$ such that:
	\begin{enumerate}
		\item (Pure high degree): $\psi$ has pure high degree at least $d$.
		\item (Non-triviality): $\|\psi\|_1 > 0$.
		\item (Sign Agreement): $\psi(x) \cdot f(x) \ge 0$ for all $x \in \bits^n$.
		\item (Appropriate Support): $\psi(x)=0$ for all $x \in \bits^n \setminus D$.
	\end{enumerate}
\end{theorem}

We refer to functions mapping $\bits^n \to \R$ as \emph{dual polynomials}. 
For a function $f \colon D \to \bits$, we refer to any $\psi$ satisfying the
%
conditions %
of \expref{Theorem}{theorem:tdegdual} as a \emph{threshold
	degree dual polynomial for $f$}, or alternatively as a \emph{dual witness to the fact that $\deg_{\pm}(f) \geq d$}.
Along the way to constructing a threshold-degree %
 dual polynomial for $f$, we will often construct dual polynomials
that \emph{almost} satisfy the Sign Agreement and Appropriate Support conditions, and it will be convenient
to give names to the inputs on which these two conditions are violated. To this end,
given a dual polynomial $\psi \colon \bits^n \to \R$, let $\mathcal{E}(\psi, f) = \{x \in D \colon \psi(x) \cdot f(x) < 0\}$, 
and let $\mathcal{W}(\psi, f) \coloneqq  \{x \in \bits^n \setminus D \colon |\psi(x)| > 0\}$. For each input $x \in \mathcal{E}$,
we say that $\psi$ \emph{makes an error} on $x$.
We will let $E(\psi, f) \coloneqq  \sum_{x \in \mathcal{E}(\psi, f)} |\psi(x)|$
and $W(\psi, f) \coloneqq  \sum_{x \in \mathcal{W}(\psi, f)} |\psi(x)|$.


\subsection{Approximate degree and its dual formulation}
\begin{definition} \label{def:adeg}
	Let $D \subseteq \{-1, 1\}^n$, and let $f$ be a function mapping $D$ to $\{-1,1\}$.
	The $\eps$-\emph{approximate degree} of $f$, denoted $\adeg_{\eps}(f)$, is the least degree of a real polynomial $p \colon \bits^n \to \R$ such that $|p(x) -f(x)| \leq \eps$ when $x \in \mathcal{D}$.
	No constraint is placed on $p(x)$ for any $x \in (\bits^n \setminus D)$.%
%
\footnote{Prior work (\eg, \cite{aaronsonshi, bkt}) has also considered  %
another natural definition of approximate degree for promise problems, that does require $p(x)$ to be bounded in magnitude outside
of the domain $D$ on which $f$ is defined. For our purposes
in this paper, the definition that does not constrain
$p$ on inputs outside of $D$ is most useful.}
\end{definition}


The following theorem provides a standard dual formulation of approximate degree.
\begin{theorem}\label{theorem:adegdual}
	
	Let $f:D \to \{-1,1\}$ with $D \subseteq \{-1,1\}^n$. Then $\adeg_{\eps}(f) > d$ if and only if there is a real function $\psi: \bits^n \to \R$ such that:
	\begin{enumerate}
		\item (Pure high degree): $\psi$ has pure high degree at least $d$.
		\item (Non-triviality): $\|\psi\|_1 > 0$.
		\item (Correlation): $\sum_{x\in\bits^n} \psi(x) \cdot f(x) \geq \eps \cdot \|\psi\|_1$.
		\item (Appropriate Support): $\psi(x)=0$ for all $x \in \bits^n \setminus D$.
	\end{enumerate}
\end{theorem}

We will also frequently use the following simple fact, which is immediate from linearity.

\begin{fact}
	\label{sumfact}
	Suppose that $\psi_1, \dots, \psi_k \colon \bits^n \to \R$ all have pure high degree at least $d$. Then their sum, $\psi_1 + \dots + \psi_k$,
	has pure high degree at least $d$.
\end{fact}



\subsection{The Surjectivity function}


\begin{definition}\label{def:surj}
	For $N \ge R$, define the function $\SURJ_{R, N} : [R]^N \to \bits$ by $\SURJ_{R, N}(s_1, \dots, s_N) = -1$ iff for every $j \in [R]$, there exists an $i$ such that $s_i = j$.
\end{definition}

When $N$ and $R$ are clear from context, we will refer to the function $\SURJ$ without the explicit dependence on these parameters. It will often be convenient to think of the input to $\SURJ_{R, N}$ as a function
mapping $\bits^n \to \bits$ rather than $[R]^N \to \bits$. 
When needed, we assume that $R$ is a power of 2 and an element of $[R]$ is encoded in binary using $\log R$ bits.
In this case we will view $\Surjectivity$ as a function on $n=N\log R$ bits, \ie,  $\SURJ: \bits^{n} \to \bits$.

For technical reasons, when proving lower bounds, it will be more convenient to work with a variant of $\SURJ$ where the range $[R]$ is augmented by a ``dummy element'' $0$ that is simply ignored by the function. That is, while any of the items $s_1, \dots, s_N$ may take the dummy value $0$, the presence of a $0$ in the input is not required for the input to be deemed surjective. We denote this variant of Surjectivity by $\dumSURJ$. More formally:

\begin{definition}\label{def:dumsurj}
	For $N \ge R$, define the function $\dumSURJ_{R, N} : [R]_0^N \to \bits$ by $\dumSURJ_{R, N}(s_1, \dots, s_N) = -1$ iff for every $j \in [R]$, there exists an $i \in [N]$ such that $s_i = j$.
\end{definition}

The following simple proposition shows that a lower bound on the approximate degree of $\dumSURJ$ implies a lower bound for $\SURJ$ itself.

\begin{proposition}[Bun, Kothari, and Thaler~\cite{bkt}] \label{prop:dumsurj-reduction} \label{relatingdsurjtosurj}
	Let $\eps > 0$ and $N \ge R$. Then
	\[ \adeg_\eps(\dumSURJ_{R, N}) \le \adeg_\eps(\SURJ_{R + 1, N + 1}) \cdot \log(R+1)\,.\]
\end{proposition}

We will also use two additional simple properties of $\SURJ_{R, N}$. The first roughly says
that increasing the range size can only make $\mathsf{SURJECTIVITY}$ harder to approximate (so long as the domain
size remains significantly larger than the range size).

\begin{proposition} \label{yoursurjprop}
	Let $\eps > 0$. Then for any $R' \geq R$, $\adeg_{\eps}(\SURJ_{R, N}) \leq \adeg_{\eps}(\SURJ_{R', N + R' - R})$. 
\end{proposition}
\begin{proof}
	Let $p : \bits^{(N+R'-R) \cdot \lceil\log(R')\rceil} \to \bits$ be a polynomial of degree $d$ that $\eps$-approximates $\SURJ_{R', N +R' - R}$. We will use to $p$ to construct a polynomial of degree $d$ that $\eps$-approximates $\SURJ_{R, N}$. Recall that an input to $\SURJ_{R, N}$ takes the form $(s_1, \dots, s_N)$ where each $s_i$ is the binary representation of a number in $[R]$. 
	For every $(s_1, \dots, s_N) \in [R]^N$, observe that
	\[\SURJ_{R, N}(s_1, \dots, s_N) = \SURJ_{R', N+R'-R}(s_1, \dots, s_N, R+1, R+2, \dots, R')\,.\]
	Hence, the polynomial
	\[p(s_1, \dots, s_N, R+1, R+2, \dots, R')\]
	has degree at most $d$ and $\eps$-approximates $\SURJ_{R, N}$. This assumes
	that for each element $r$ of $[R]$, each bit of the encoding of $r$ in $\bits^{\lceil\log R'\rceil}$ (\ie, when $r$ is viewed as an element of $[R']$),
	is a degree-1 function of its encoding in $\bits^{\lceil \log R \rceil}$ (\ie, when $r$ is viewed as an element of $[R]$). This property 
	holds for the natural binary encoding. Even if some encoding were used that did not satisfy this %
	condition, %
	then $p(s_1, \dots, s_N, R+1, R+2, \dots, R')$ will have degree no larger than $d \cdot (\log(R') + 1)$.
\end{proof}

The second says that increasing the domain size can only make $\mathsf{SURJECTIVITY}$ harder to approximate. 
\begin{proposition} \label{yoursurjprop2}
	Let $\eps > 0$. If $N' > N$, then $\adeg_{\eps}(\SURJ_{R, N})  \leq \adeg_{\eps}(\SURJ_{R+1, N'})$.
\end{proposition}
\begin{proof}
	For every $(s_1, \dots, s_N) \in [R]^N$, observe that
	\[\SURJ_{R, N}(s_1, \dots, s_N) = \SURJ_{R+1, N'}(s_1, \dots, s_N, R+1, R+1, \dots, R+1),\]
	where the element $R+1$ is repeated $N'-N$ times. 
	Hence, if $p : \bits^{N' \cdot \lceil \log(R+1) \rceil} \to \bits$ is a polynomial of degree $d$ that $\eps$-approximates $\SURJ_{R+1, N'}$,
	then $p(s_1, \dots, s_N, R+1, R+1, \dots, R+1)$ is a degree $d$ polynomial that $\eps$-approximates $\SURJ_{R, N}$.
\end{proof}

\subsection{A useful auxiliary function}
The following lemma, due to Razborov and Sherstov~\cite{razborovsherstov}, gives a function that we will use
many times in this paper to ``zero out'' mass that intermediate dual witnesses $\psi$ place
on the ``bad'' sets $\mathcal{E}(\psi, f)$ and $\mathcal{W}(\psi, f)$. 

\begin{lemma}[%
	\!{\cite[Proof of Lemma 3.2]{razborovsherstov}}]\label{lem:rs}
	Let $D, n \in \N$ with $0 \le D \le n - 1$. Then for every $y \in \bits^n$ with $|y| > D$, there exists a function $\phi_y: \bits^n \to \mathbb{R}$ such that
	\begin{equation}\phi_y(y) = 1	\label{eqn:rs-point}\end{equation}
	\begin{equation}|x| > D, x \ne y \implies \phi_y(x) = 0	\label{eqn:rs-hamming}\end{equation}
	\begin{equation}\deg p \le D \implies \langle \phi_y, p \rangle = 0 \label{eqn:rs-phd} \end{equation}
	\begin{equation}\sum_{|x| \le D} |\phi_y(x)| \le 2^D {|y| \choose D}\,.   \label{eqn:rs-norm} \end{equation}
\end{lemma}
%

\subsection{A dual witness for \texorpdfstring{$\OR$}{OR}}
We will repeatedly make use of a dual witness for the high approximate
degree of the $\OR$ function. The dual witness itself was essentially first constructed by {\v S}palek~\cite{spalek}, 
but we require an analysis of it due to Bun, Kothari, and Thaler~\cite{bkt}.

\begin{proposition}[\!\cite{bkt}] \label{prop:or-sym-dual}  %
	There exist constants $c_1, c_2 \in (0, 1) $ for which the following holds. Let $T \in \N$ and $1/T \le \delta \le 1/2$. Then there is a function $\omega : [T]_0 \to \R$ such that
	\begin{align}
	&\omega(0) - \sum_{t = 1}^T \omega(t) \ge 1- \delta\,,  \label{eqn:or-sym-corr}\\
	& \omega(0) \ge (1-\delta)/2\,,\\
	  &\sum_{t = 0}^T |\omega(t)| = 1\,, \\
%
	& \text{for all univariate polynomials } q \colon \R \to \R\text{, } \deg q < c_1\sqrt{\delta T} \implies \sum_{t=0}^T \omega(t) \cdot q(t) = 0\,, \label{eqn:or-sym-phd}\\
	& |\omega(t)| \le \frac{170 \exp(-c_2 t \sqrt{\delta} /  \sqrt{T})}{\delta \cdot t^2} \qquad \forall t = 1, \dots, T\,.\label{eqn:or-sym-decay}\end{align}
\end{proposition}

\begin{proposition} \label{prop:or-dual}
	Let $T, N \in \N$ with $T \le N$, and let $\delta > 1/T$. Define $\omega$ as in \expref{Proposition}{prop:or-sym-dual}. Define the function $\psi: \bits^N \to \R$ by $\psi(x) = \omega(|x|) / {N \choose |x|}$ for $x \in \mathcal{H}_N^{\le T}$ and $\psi(x) = 0$ otherwise. Then $\psi$ satisfies:
	\begin{align} &\langle \psi, \OR_N \rangle \ge 1-\delta\,,  \label{eqn:or-corr} \\
& \psi(\mathbf{1}_{N}) \ge (1-\delta)/2\,, \label{eqn:or-loaded}\\
	& \|\psi\|_1 = 1\,,   \label{eqn:or-norm}\\
%
	& \text{for any polynomial } p \colon \bits^N \to \R\text{, } \deg p < c_1\sqrt{\delta T} \implies \langle \psi, p \rangle = 0\,, \label{eqn:or-phd}\\
	& \sum_{|x|=t} |\psi(x)| \leq \frac{170 \exp(-c_2 t \sqrt{\delta} /  \sqrt{T})}{\delta \cdot t^2} \qquad \forall t = 1, \dots, T\,.\label{eqn:or-decay}
	\end{align}
\end{proposition}

\subsection{Block composition of functions}
\begin{definition}
	For functions $f : Y^n \to Z$ and $g: X \to Y$, define the \emph{block composition} $f \circ g : X^n \to Z$ by $(f\circ g)(x_1, \dots, x_n) = f(g(x_1), \dots, g(x_n))$, for all $x_1, \dots, x_n \in X$.
\end{definition}

\subsection{Dual block composition}
\label{s:dualblockcomposition}
We now define the dual block method \cite{shizhu, lee, sherstovhalfspaces1} for combining dual witnesses for functions $f$, $g$ in order to obtain a dual witness
for $f \circ g$. We then state two basic properties
%
of the method.   %
\label{sec:ls}
\begin{definition}
	Let $\Psi : \bits^M \to \R$ and $\psi : \bits^m \to \R$ be functions that are not identically zero. Let $x=(x_1, \dots, x_M) \in \left(\bits^{m}\right)^M$. Define the \emph{dual block composition} of $\Psi$ and $\psi$, denoted $\Psi \ls \psi : (\bits^m)^M \to \R$, by
	\[(\Psi \ls \psi)(x_1, \dots, x_M) = 2^M \cdot \Psi(\dots, \sgn\left(\psi(x_i)\right), \dots) \cdot \prod_{i = 1}^M |\psi(x_i)|\,.\]
\end{definition}


\begin{proposition}[\!\cite{sherstovhalfspaces1,adegsurj}] \label{prop:ls}
	\label{prop12}
	The dual block composition %
has  %
the following properties:
	\paragraph{\emph{Preservation of $\ell_1$-norm}:} Assume that $\psi$ has pure high degree at least 1.
	If $\|\Psi\|_1 = 1$ and $\|\psi\|_1 = 1$, then
	\begin{equation}\|\Psi \ls \psi\|_1 = 1\,. \label{eqn:ls-norm}\end{equation}
	
	\paragraph{\emph{Multiplicativity of pure high degree}:} If $\langle \Psi, P \rangle = 0$ for every polynomial $P \colon \bits^M \to \bits$ of degree less than  $D$, and $\langle \psi, p \rangle = 0$ for every polynomial $p \colon \bits^m \to \bits$ of degree less than $d$, then
	for every polynomial $q \colon \bits^{m \cdot M} \to \bits$,
	\begin{equation}\deg q < D \cdot d \implies \langle \Psi \ls \psi, q \rangle = 0\,. \label{eqn:ls-phd} \end{equation}
	
	\paragraph{\emph{Associativity}:} For every $\zeta: \bits^{m_\zeta} \to \R$, $\varphi: \bits^{m_\varphi} \to \R$, and $\psi: \bits^{m_\psi} \to \R$, we have
	\begin{equation} (\zeta \ls \varphi) \ls \psi = \zeta \ls (\varphi \ls \psi)\,. \label{eqn:ls-assoc} \end{equation}
\end{proposition}

Given three dual witnesses $\zeta$, $\varphi$, and $\psi$, the associativity property allows us to write $\zeta \ls \varphi \ls \psi$ without ambiguity.

Next, we state an important lemma that follows directly from an analysis of Bun, Kothari, and Thaler \cite[Proof of Proposition 31]{bkt}.
This lemma shows that if $\psi$ is the dual witness for $\OR_N$
from \expref{Proposition}{prop:or-dual}, then for \emph{any} $\Phi \colon \bits^R \to \bits$, $\xi = \Phi \ls \psi$ places an exponentially small amount of mass on inputs outside of $\mathcal{H}^{\le N}_{N \cdot R}$.\footnote{Our \expref{Lemma}{prop:prelim-mass} follows directly from an intermediate calculation in the proof of Proposition 31 of~\cite{bkt}.}

\begin{lemma}[\cite{bkt}] \label{prop:prelim-mass}
	Fix a parameter $1 \leq \alpha \leq R^2$, and assume that $N \geq \lceil 20 \sqrt{\alpha} \rceil R$.
	For any $\beta \in (0,1)$
	assume that $\Phi : \bits^R \to \R$ satisfies $\|\Phi\|_1=1$. Furthermore, let $\psi : \bits^N \to \R$ be symmetric
	(meaning $\psi(x)$ depends only on $|x|$) and assume that $\psi$ satisfies the following conditions: 
	\begin{align}
&\sum_{t = 0}^N \sum_{|x|=t} \psi(x) = 0\,, \label{eqn:omega-balanced}\\
&\sum_{t = 0}^N \sum_{|x| = t} |\psi(x)| = 1\,, \label{eqn:omega-norm}\\
&\sum_{|x| = t} |\psi(x)| \le \alpha \exp (-\beta t) / t^2 \quad \forall t = 1, 2, \dots, N\,. \label{eqn:omega-decay}
	\end{align}
	Then for sufficiently large $R$, 
	\begin{equation}
	\sum_{x \notin \mathcal{H}^{\le N}_{N \cdot R}} |(\Phi \ls \psi)(x)| \le \frac{2}{\beta} \exp\left(-\frac{\beta N}{6 \ln R}\right)\,.
	\label{eqn:prelim-mass}
	\end{equation}
	
	In particular, if $\psi$ is the dual witness for $\OR_N$ obtained from \expref{Proposition}{prop:or-dual} with constant parameter $\delta \in (0, 1)$ and with $T=N/\log^3(N)$,
	then there is a constant $c_3 > 0$ such that
	\begin{equation}
	\sum_{x \notin \mathcal{H}^{\le N}_{N \cdot R}} |(\Phi \ls \psi)(x)| \le 2^{-c_3 \cdot \sqrt{N \log N}}\,. \label{eqn:prelim-mass-OR}
	\end{equation}
\end{lemma}


\subsection{Hard functions and Hamming weight promises}

In our %
earlier paper %
\cite{adegsurj}, we identified a natural class 
of functions satisfying the following %
condition: %
for any function $h \colon \bits^n \to \bits$ in the class,
there is a related function $H \colon \bits^{\poly(n)} \to \bits$ such that 
$\adeg_{\eps}(h) = \tilde{\Theta}(\adeg_{\eps}(H^{\leq N}))$ for some $N=\tilde{\Theta}(n)$.
This enables one to prove $\eps$-approximate degree (respectively, threshold degree) lower bounds on $h$ by lower bounding
the $\eps$-approximate degree (threshold degree) of the related function $H$, \emph{which is defined on significantly more
	variables than $h$ itself}, under the promise that the Hamming weight of $H$'s input is at most $N$.

In more detail, this connection is as follows. 
Fix $R, N \in \N$, let $f : \bits^R \to \bits$ and let $g : \bits^N \to \bits$. Suppose $g$ is a \emph{symmetric} function, in the sense that for any $x \in \bits^N$ and any permutation $\sigma : [N] \to [N]$, we have
\[g(x_1, \dots, x_N) = g(x_{\sigma(1)}, \dots, x_{\sigma(N)})\,.\]
Equivalently, the value of $g$ on any input $x$ depends only on its Hamming weight $|x|$.

The functions $f$ and $g$ give rise to two functions. The first, denoted by $F^{\mathsf{prop}} : [R]^N_0 \to \bits$, is a certain property of a list of numbers $s_1, \dots, s_N \in [R]_0$. The second, which we denote by $F^{\le N} : \mathcal{H}_{N \cdot R}^{\le N} \to \bits$, is the block composition of $f$ and $g$ restricted to inputs of Hamming weight at most $N$. Formally, these functions are defined as:
\begin{align*}
&F^{\mathsf{prop}}(s_1, \dots, s_N) = f(g(\mathds{1}[s_1 = 1], \dots, \mathds{1}[s_N = 1]), \dots, g(\mathds{1}[s_1=R], \dots,\mathds{1}[s_N=R])) \\
&F^{\le N}(x_1, \dots, x_R) = \begin{cases} f(g(x_1), \dots, g(x_R)) & \text{ if } x_1, \dots, x_R \in \bits^N,  |x_1| + \dots + |x_R| \le N\\
\text{undefined} & \text{otherwise}.
\end{cases}
\end{align*}

Observe that for $f=\AND_R$ and $g = \OR_N$, $F^{\mathsf{prop}}=\dumSURJ_{R, N}$. 
That is, $\dumSURJ_{R, N}$ can be expressed as an $\AND_R$ (over all range items $r \in [R]$) of the $\OR_N$ (over all inputs $i \in [N]$) of ``Is input $x_i$ equal to $r$?''.

The following proposition relates the approximate degrees of the two functions $F^{\mathsf{prop}}$ and $F^{\le N}$.

\begin{theorem}[Bun and Thaler~\cite{adegsurj}] \label{thm:main-reduction}
	Let $f : \bits^R \to \bits$ be any function and let $g : \bits^N \to \bits$ be a symmetric function. Then for $F^{\mathsf{prop}}$ and $F^{\le N}$ defined above, we have
	\[\adeg_{\eps}(F^{\mathsf{prop}}) \ge \adeg_{\eps}(F^{\le N})\,.\]
	In addition, \[\deg_{\pm}(F^{\mathsf{prop}}) \ge \deg_{\pm}(F^{\le N})\,.\]
\end{theorem}




\section{The threshold degree of Surjectivity: Upper bound}
\label{app:ub}
For completeness, we prove a tight upper bound on the threshold degree of $\SURJ_{R, N}$.
This upper bound follows from standard techniques.
\begin{claim}
	The function $\SURJ_{R, N}$ has threshold degree $O\left(\min\left\{R  \cdot \log R, N^{1/2} \cdot \log^{3/2} R\right\}\right)$.
\end{claim}
\begin{proof}
	We assume throughout that $R \leq N$, the claim being trivial otherwise.
	This claim follows from two well-known facts, stated as \expeqref{Expressions}{nomoremore} \expeqref{and}{nomoremore2} below.
	
	\begin{equation} \label{nomoremore}\deg_{\pm}(\SURJ_{R, N}) \leq \deg_{\pm}(\AND_R \circ \OR_N) \cdot \log R \end{equation}
	
	\begin{equation} \label{nomoremore2} \deg_{\pm}(\AND_R \circ \OR_N) = O\left(\min\left\{R, (N \log R)^{1/2}\right\}\right) \end{equation}
%
	
	\expeqref{Expression}{nomoremore} holds because, as mentioned in \expref{Section}{surjoverview}, 
	$\SURJ_{R, N}$ is equivalent to the $\AND_R$ (over all range items $r \in [R]$) of the $\OR_N$ (over all inputs $i \in [N]$) of ``Is input $s_i$ equal to $r$?'', and the quoted question is computed by a conjunction of width $\log R$ over the input bits. 
	That is, if $$y_{i,j} \coloneqq  \begin{cases} -1 \text{ if } s_i=j \\ 1 \text{ otherwise,}\end{cases}$$
	then $$\SURJ_{R, N}(x) = \AND_R(\OR_N(y_{1, 1} ,\dots, y_{N, 1}), \dots, \OR(y_{1, R}, \dots, y_{N, R})).$$
	Hence, if $p$ sign-represents $\AND_R \circ \OR_N$, then $p(y_{1, 1}, \dots,  y_{N, R})$ is a polynomial of degree $\deg(p) \cdot \log R$
	that sign-represents $\SURJ_{R, N}$. 
	
	\expeqref{Equation}{nomoremore2} holds by the following standard analysis. 
	To show that $$\deg_{\pm}(\AND_R \circ \OR_N) \leq O((N \log N)^{1/2}),$$ 
	we exploit the well-known fact that $\adeg_{1/(3R)}(\OR_N) = \Theta(\sqrt{N \log R})$ \cite{kahn, buhrman1999bounds}. 
	Let $p$ be a minimal degree polynomial that uniformly approximates $\OR$ to error $1/(3R)$.
	For inputs $x=(x_1, \dots, x_R) \in \left(\bits^{N}\right)^{R}$, 
	the polynomial $\sum_{j=1}^{R} p(x_j) + R-1$ sign-represents $\AND_R \circ \OR_N$,
	and has degree at most $\deg(p) = O(\sqrt{N \log R})$. 
	
	To show that $\deg_{\pm}(\AND_R \circ \OR_N) \leq O(R)$, we use the fact that there exist two linear functions $p$ and $q$ such that the ratio
	$$\frac{p(x)}{q(x)} = \frac{-\left(N - \sum_{i=1}^N x_i\right) + 1/(3R)}{\left(N - \sum_{i=1}^N x_i \right)+ 1/(3R)}$$
	satisfies $|p(x) / q(x) - \OR_N(x)| \le 1 / (3R)$ for every $x$. 
	Then the following sum of rational functions agrees in sign with $\AND_R \circ \OR_N$
	at all inputs $x=(x_1, \dots, x_R) \in \left(\bits^R\right)^N$: $\sum_{j=1}^R p(x_j)/q(x_j) + R-1= \sum_{j=1}^R p(x_j)q(x_j)/q^2(x_j) + R-1$.
	Placing everything over the non-negative common denominator $\prod_{j=1}^R q^2(x_j)$, it follows that the polynomial $$(R-1) \cdot \prod_{j=1}^R q^2(x_j) + \sum_{j=1}^R p(x_j) q(x_j) \prod_{k=1, \dots, R \colon k \neq j} q(x_j)^2$$
	sign-represents $\AND_R \circ \OR_N$. The degree of this polynomial is  $O(R)$.
\end{proof}


\section{The threshold degree of Surjectivity: Lower bound}
\label{s:surj}
The key technical result established in this section is \expref{Theorem}{surjthmintro}, restated here for the reader's convenience.
\surjthmintro*

By combining \expref{Theorems}{thm:main-reduction}, \ref{relatingdsurjtosurj}, \expref{and}{surjthm}, we obtain the desired lower bound on the threshold 
degree of $\SURJ$.
\begin{corollary} \label{surjcor}
	For $R=N^{1/2}$, $\deg_{\pm}(\SURJ_{R, N}) = \Omega(N^{1/2}/\log^{5/2} N)$. 
\end{corollary}

In fact, combining \expref{Corollary}{surjcor} with \expref{Proposition}{yoursurjprop} %
determines %
 the threshold degree of $\SURJ_{R, N}$ up to logarithmic factors,
for all settings of $R$ and $N$.

\begin{corollary}[Restatement of \expref{Theorem}{thm:surj}]
  For any constant $c<1$, if $1 < R < cN$, then
%
\[ \deg_{\pm}\left(\SURJ_{R, N}\right) = 
\tilde{\Omega}\left(\min(R, N^{1/2})\right) \,.\]
\end{corollary}
\begin{proof}
	If $N^{1/2} \leq R \leq cN$, let $N'=\Theta(N)$ be such that $N = N' + R - (N')^{1/2}$. Then \expref{Proposition}{yoursurjprop} implies that 
	$\deg_{\pm}\left(\SURJ_{R, N}\right) \geq \deg_{\pm}\left(\SURJ_{(N')^{1/2}, N'}\right) = \tilde{\Omega}(N^{1/2})$, where the final
	equality holds by \expref{Corollary}{surjcor}.
	
	If $R < N^{1/2}$, then \expref{Proposition}{yoursurjprop2} 
	implies that $\deg_{\pm}(\SURJ_{R, N}) \geq \deg_{\pm}(\SURJ_{R-1, (R-1)^2}) = \tilde{\Omega}(R)$,
	where the final equality holds by \expref{Corollary}{surjcor}.
\end{proof}

%
%
%
%
\begin{proof}[Proof of Theorem \ref{surjthm}]
	Let $d=N^{1/2}/(c \log^{3/2} N)$ for a sufficiently large constant $c$ to be chosen later. 
	Let 
	$\psi$ be the dual witness for $\OR_{N}$ from \expref{Proposition}{prop:or-dual} with $T=N/\log^3 N$ and $\delta=1/4$. 
	Then $\psi$ has pure high degree at least $d$. %
	Fix any positive integer $r$, and 
	define $\Phi \colon \bits^{r} \to \reals$ as follows:
	$$\Phi_r(w) = \begin{cases}
	-1/2 & \text{ if } w=(-1, -1, \dots, -1)\\
	1/2 & \text{ if } w= (1, 1, \dots, 1)\\
	0 & \text{otherwise}.
	\end{cases}$$
	
Observe that $\Phi_r$ has pure high degree 1, and also $\|\Phi_r\|_1=1$. 
We remark that the dual block composition $\Phi_r \ls \psi$ is essentially the same dual witness constructed in~\cite{bt14} to show that
$\adeg_{\eps}(\AND_r \circ\OR_N) = \Omega(N^{1/2})$ for
$\eps = 1-2^{-r}$ (see \expeqref{Inequality}{btprop3} below).
This dual witness was also used in subsequent
work \cite{sherstov14, bchtv, bkt}.
	
%
%
We observe that $\Phi_r \ls \psi$ has the following properties. %
%

	\begin{equation}  \label{btprop1}
	\|\Phi_r \ls \psi\|_1=1\,.
	\end{equation}
	%
	\begin{equation}  \label{btprop2}
	\Phi_r \ls \psi \text{ has pure high degree at least } d\,.
	\end{equation}
	%
	\begin{equation}  \label{btprop3}
	E(\Phi_r \ls \psi, \AND_r \circ \OR_N) \leq 2^{-r}\,.
	\end{equation}
	%
	\begin{equation}  \label{btprop4}
	\left(\Phi_r \ls \psi\right)\left(\mathbf{1}_{r \cdot N}\right) \geq \frac12 \cdot (3/4)^{r}\,.
	\end{equation}
	%
	\begin{equation} \label{btprop5}
	\left(\Phi_r \ls \psi\right)(y) \geq 0 \text{ for all } y \in \Hcal_{r \cdot N}^{\leq r}\,. 
	\end{equation}
	%
	\begin{equation} \label{btprop6}
\text{There is some constant } c_3>0 \text{ such that } W(\Phi_r \ls \psi, (\AND_r \circ \OR_N)^{\leq N}) \leq 2^{-c_3 \sqrt{N \log N}}\,.
	\end{equation}
	%
	
\medskip  %
	\noindent \textbf{Justification for \expeqref{Properties}{btprop1}-\eqref{btprop6}.} \expeqref{Equation}{btprop1} follows from 
	\expref{Proposition}{prop12} (see \expeqref{Equation}{eqn:ls-norm}) and the fact that $\|\Phi_r\|_1=\|\psi\|_1=1$.
	\expeqref{Equation}{btprop2} follows from \expref{Proposition}{prop12} (see \expeqref{Equation}{eqn:ls-phd}),
	and the fact that the pure high degree of $\Phi_r$ is at least 1, and the pure high degree of 
	$\psi$ is at least $d$.
	\expeqref{Inequality}{btprop3} is an immediate consequence of \cite[Theorem 1]{bt14}.
	\expeqref{Inequality}{btprop4} is immediate from the definition of dual block composition
	and the fact that $\psi(\mathbf{1}_N) \geq 3/8$ (see \expeqref{Inequality}{eqn:or-loaded}). 
	
	\medskip
	\expeqref{Property}{btprop5} holds because $\left(\Phi_r \ls \psi\right)(x) < 0$ only if $\psi(x_i) < 0$ for all $i=1, \dots, r$,
	and $\psi(x_i) < 0$ implies that $|x_i| \geq 1$ (see \expeqref{Inequality}{eqn:or-loaded}). Hence,  $\left(\Phi_r \ls \psi\right)(x) < 0 \Longrightarrow |x| \geq r$.
	
	\medskip
	Finally, 
	\expeqref{Property}{btprop6} is immediate from \expref{Lemma}{prop:prelim-mass} (see \expeqref{Inequality}{eqn:prelim-mass-OR}) by choosing $c_3$ sufficiently large.
	%
	
	\medskip
	\noindent \textbf{Definitions of auxiliary functions $\nu_x$.}
	Fix any $x \in \Hcal_{R \cdot N}^{\leq d}$. We think of $x$ as consisting of $R$ blocks, each of length $N$. Accordingly,  denote 
	$x=(x_1, \dots, x_{R})$, where each $x_i \in \bits^{N}$. 

	Let $\mathcal{I}=\{i \colon x_i = \mathbf{1}_N\}$, and let $r(x)=|\mathcal{I}|$. Enumerate the set $\mathcal{I}$ as $\mathcal{I} = \{i_1, i_2, \dots, i_{r(x)}\}$.
	Since $|x| \leq d$, $r(x) \geq R-d.$ 
	Define the function $\nu_x \colon \left(\bits^{N}\right)^R \to \reals$ by 
	\[
	\nu_x(y_1, \dots, y_R) = 
	\begin{cases}
	(\Phi_{r(x)} \ls \psi)(y_{i_1}, \dots, y_{i_{r(x)}}) & \text{ if } y_i = x_i \text{ for all } i \notin \mathcal{I}\\
	0 & \text{ otherwise}.
	\end{cases} 
	\]
	
%
%
%
We observe that $\nu_x$ has the following properties. %
%
	\begin{equation}
	\label{prop1}
	\|\nu_x\|_1=1\,.
	\end{equation}
	\begin{equation}
	\label{prop2}
	\nu_x \text{ has pure high degree at least } d\,.
	\end{equation}
	\begin{equation} \label{prop3}
	E\left(\nu_x, \AND_R \circ \OR_N\right) = E\left(\Psi_{r(x)} \ls \psi, \AND_{r(x)} \circ \OR_N\right) \leq 2^{-r(x)}\,.
	\end{equation}
	\begin{equation} \label{prop4} 
	\nu_x(x)\cdot \left(\AND_R \circ \OR_N\right)(x) = \left(\Phi_{r(x)} \ls \psi\right)\left(\mathbf{1}_{R \cdot N}\right) \geq 1/2 \cdot (3/4)^{r(x)} \geq 1/2 \cdot (3/4)^{R}\,.
	\end{equation} 
	\begin{equation}
	\label{prop5} 
	\nu_x(y) \geq 0 \text{ for all } y \in \Hcal_{RN}^{\leq d}\,.
	\end{equation}
	\begin{equation}
	\label{prop6} 
	W\left(\nu_x, \left(\AND_R \circ \OR_N\right)^{\leq N + d}\right) \leq 2^{-c_3 \sqrt{N \log N}}\,.
	\end{equation}
	
	\medskip
\noindent \textbf{Justification for \expeqref{Properties}{prop1}-\eqref{prop6}.} 
	\expeqref{Equations}{prop1} \expeqref{and}{prop2} are immediate from \expeqref{Equations}{btprop1} \expeqref{and}{btprop2} respectively.
	
	\medskip
	
	We now derive \expeqref{Expression}{prop3}.
	The first equality holds by the following
	reasoning. Since $x_i \neq \mathbf{1}_N$ for all $i \not \in \mathcal{I}$,
	$z_i = x_i \text{ for all } i \notin \mathcal{I} \Longrightarrow 
	\AND_R \circ \OR_N(z) = \left(\AND_{r(x)} \circ \OR_N\right)(z_{i_1}, \dots, z_{i_{r(x)}})$. 
	The inequality in \expeqref{Expression}{prop3} follows from \expeqref{Inequality}{btprop3}.
	
	
	\medskip \expeqref{Expression}{prop4}
	is immediate from the definition of dual block composition and the fact that $\sgn(\psi(\mathbf{1}_N)) > 0$ (see 
	\expeqref{Expression}{eqn:or-loaded}). 
	\expeqref{Property}{prop5} is immediate from \expeqref{Property}{btprop5} and the fact that $2d \le R$. \expeqref{Inequality}{prop6} is immediate from \expeqref{Inequality}{btprop6}, the definition of $\nu_x$, and the fact that $|x| \leq d$. 
	
	\medskip
	\noindent \textbf{Another intermediate dual witness.}
	Let $T=|\Hcal_{RN}^{\leq d}| = \sum_{i=0}^d {RN \choose i}$. Consider the dual witness $\eta=\frac1T \cdot \sum_{x \in \Hcal^{\leq d}_{RN}} \nu_x$. 
	Let $r_{\text{min}} = \min_{x \in \mathcal{H}_{RN}^{\leq d}} r(x) \geq R-d \geq R/2$. 
	
%
%
%
We observe that $\eta$ has the following properties. %
	\begin{equation}
	\label{etaprop1}
	\|\eta\|_1 \leq 1\,. 
	\end{equation}
	\begin{equation}
	\label{etaprop1prime}
	\|\eta\|_1 \ge  1-2^{1-r_{\text{min}}} \,.
	\end{equation} 
	\begin{equation}
	\label{etaprop2}
	\eta \text{ has pure high degree at least } d\,.
	\end{equation}
	\begin{equation} \label{etaprop3} E(\eta, \AND_R \circ \OR_N) \leq 2^{-r_{\text{min}}}\,.
	\end{equation}
	\begin{equation}
	\text{ For all } y \in \mathcal{H}^{\leq d}_{RN}, \eta(y) \geq (1/T) \cdot 1/2 \cdot (3/4)^{R}\,.
	\label{etaprop5}
	\end{equation} %
	\begin{equation} %
	W(\eta, (\AND_R \circ \OR_N)^{\leq N + d}) \leq 2^{-c_3 \sqrt{N \log N}}\,.
	\label{etaprop6}
	\end{equation}
	
	\medskip \noindent \textbf{Justification for \expeqref{Properties}{etaprop1}-\eqref{etaprop6}.}
	To justify \expeqref{Inequality}{etaprop1}, observe that
	\begin{align*} \sum_{y \in \bits^{RN}} |\eta(y)| \leq \sum_{y \in \bits^{RN}}\frac1T \sum_{x \in \Hcal^{\leq d}_{RN}} |\nu_x(y)|
	=\frac1T \sum_{x \in \Hcal^{\leq d}_{RN}}  \sum_{y \in \bits^{RN}} |\nu_x(y)|\\ =
	\frac1T \sum_{x \in \Hcal^{\leq d}_{RN}} \|\nu_x\|_1= 1\,,
	\end{align*}
	where the final equality is an immediate  consequence of \expeqref{Equation}{prop1}.
	
	To justify \expeqref{Inequality}{etaprop1prime}, observe that 
	\begin{align*}
	\sum_{y \in \bits^{RN}} |\eta(y)| = \frac1T \sum_{y \in \bits^{RN}} \left|\sum_{x \in \Hcal^{\leq d}_{RN}}  \nu_x(y)\right| \\
	\geq \frac1T \left( \left(\sum_{y \in \bits^{RN}} \sum_{\substack{x \in \Hcal^{\leq d} \colon\\ y \not \in \mathcal{E}(\nu_x, \AND_R \circ \OR_N)}} \left|\nu_x(y)\right|\right)\right)
	-  \left( \left(\sum_{y \in \bits^{RN}} \sum_{\substack{x \in \Hcal^{\leq d} \colon\\ y \in \mathcal{E}(\nu_x, \AND_R \circ \OR_N)}}  |\nu_x(y)|\right)\right)\\
	= \frac1T \left( \left(\sum_{x \in \Hcal^{\leq d}_{RN}} \sum_{\substack{y \in \bits^{RN} \colon \\ y \not \in \mathcal{E}(\nu_x, \AND_R \circ \OR_N)}} \left|\nu_x(y)\right|\right)\right)
	-  \left( \left(\sum_{x \in \Hcal^{\leq d}_{RN}}  \sum_{\substack{y \in \bits^{RN} \colon\\ y \in \mathcal{E}(\nu_x, \AND_R \circ \OR_N)}} \left|\nu_x(y)\right|\right)\right)\\
	= \frac1T \left( \left(\sum_{x \in \Hcal^{\leq d}_{RN}} \left(1-E\left(\nu_x, \AND_R \circ \OR_N\right)\right)\right) - \left( \sum_{x \in \Hcal^{\leq d}_{RN}} E\left(\nu_x, \AND_R \circ \OR_N\right)\right)\right)\\
	=  \frac1T \sum_{x \in \Hcal^{\leq d}_{RN}} 1-2 \cdot E\left(\nu_x, \AND_R \circ \OR_N\right) \geq 1-2^{1-r_{\text{min}}}\,,
	\end{align*}
where the final statement holds by \expeqref{Expression}{prop3}.
	
	\expeqref{Equation}{etaprop2} is immediate from \expeqref{Equation}{prop2} and \expref{Fact}{sumfact}.
	
	\expeqref{Expressions}{etaprop3} \expeqref{and}{etaprop6} are respectively immediate from \expeqref{Expressions}{prop3} \expeqref{and}{prop6} and the triangle inequality.
	
	\expeqref{Expression}{etaprop5} follows from \expeqref{Expressions}{prop4} \expeqref{and}{prop5}.
	
	\medskip
	Let $\zeta \coloneqq  \eta/\|\eta\|_1$. By the definition of $\zeta$, and \expeqref{Equations}{etaprop1}-\eqref{etaprop6}, $\zeta$ has the following properties. %
	
	\begin{equation}
	\label{zetaprop1} \|\zeta\|=1\,.
	\end{equation}
	\begin{equation}
	\label{zetaprop2}
	\zeta \text{ has pure high degree at least } d\,.
	\end{equation}
	\begin{equation}
	\label{zetapropr3}
	E\left(\zeta, \AND_R \circ \OR_N\right) \leq 2^{-r_{\text{min}}}/\left(1-2^{-r_{\text{min}}}\right) \leq  2^{1-r_{\text{min}}}\,.
	\end{equation}
	\begin{equation}
	\label{zetaprop4}
	W\left(\zeta, \left(\AND_R \circ \OR_N\right)^{\leq N+d}\right) \leq 2^{-c_3 \sqrt{N \log N}}/(1-2^{-r_{\text{min}}}) \leq 2^{1-c_3 \sqrt{N \log N}}\,.
	\end{equation}
	\begin{equation}
	\label{zetaprop5}
	\text{ For all } x \in \mathcal{H}^{\leq d}_{RN}, \zeta(x) \geq (1/T) \cdot 1/2 \cdot (3/4)^{R}\,.
	\end{equation}
	
	Let $s = (1/T) \cdot 1/2 \cdot (3/4)^R$ denote the right hand side of \expeqref{Expression}{zetaprop5}\,. 
Recalling that $T=\sum_{i=0}^{d} {RN \choose i} \leq (1+RN)^d$, and that $d=N^{1/2}/(c \log^{3/2} N)$ for a constant $c$ of our choosing, 
	taking $c$ sufficiently large ensures that 
	\begin{equation} \label{zeq} s \geq 2^{-R/5}\,. \end{equation}
	
	
	
	\medskip
	\noindent \textbf{The final dual witness.} Let $f=\left(\AND_R \circ \OR_N\right)^{\leq N+d}$. We now modify $\zeta$ using \expref{Lemma}{lem:rs} to zero out the mass on the ``bad sets'' $\mathcal{E}(\zeta, \AND_R \circ \OR_N)$ and $\mathcal{W}(\zeta, f)$,
	thereby constructing a %
	threshold-degree %
	 dual witness $\mu$ for $f$. 
	Specifically, let $B=\mathcal{W}(\zeta, f) \cup \mathcal{E}(\zeta, \AND_R \circ \OR_N)$.
	By \expeqref{Expression}{zetaprop5},
	and the fact that $f(x) =1$ for all $x \in \mathcal{H}_{RN}^{\leq d}$,
	it holds that 
	\begin{equation}
	\label{subseteq} B \subseteq \bits^{RN}\setminus \mathcal{H}_{RN}^{\leq d}\,.
	\end{equation}
	Hence, for every $x \in B$, we can invoke \expref{Lemma}{lem:rs} with $D=d$ to
	obtain a function $\phi_x$ having  %
	\expeqref{Properties}{eqn:rs-point}-\eqref{eqn:rs-norm}.
	For every $x \in B$, 
	define $\psi_{\mathsf{corr}, x} \coloneqq  \zeta(x) \cdot \phi_x$. %
	Our final dual witness $\mu$ is 
	\begin{equation} \mu \coloneqq  \zeta - \sum_{x \in B} \psi_{\mathsf{corr}, x}\,. \label{omg1}\end{equation}
	
	
	
	\medskip
	\noindent \textbf{Analysis of $\mu$: A brief overview.}
	Since the correction terms in the definition of $\mu$ are specifically designed to have pure high degree $d$
	and zero out all of the ``bad'' mass of $\zeta$, all that remains in order to show that
	$\mu$ witnesses a degree $d$ lower bound on the threshold degree of $f$ is show that the correction terms do not
	disturb the sign of any inputs in their support. To accomplish this,
	we start by observing that the total $\ell_1$-mass of all the correction objects is at most $(W(\zeta, f)+E(\zeta, \AND_R \circ \OR_N)) \cdot 2^{d} \cdot {RN \choose d}$. 
	Recalling that $R=N^{1/2}$ and $d = N^{1/2}/(c \log^{3/2} N)$ for a constant $c$ of our choosing, by \expeqref{Properties}{zetapropr3} \expeqref{and}{zetaprop4} we can set $c$ sufficiently large to ensure that both 
	$R-d - 1 > 4R/5$ and $c_3 \sqrt{N \log N} -1 > 4R/5$.
	Then \begin{equation}
	\label{akeyeq} \left(W(\zeta, f)+E\left(\zeta, \AND_R \circ \OR_N\right)\right) \cdot 2^{d} \cdot {RN \choose d} \leq 2^{-4R/5} \cdot (RN)^d \leq  2^{-3R/5} \leq s/2\,,\end{equation} 
	where the final inequality holds by \expeqref{Property}{zeq}.
	Hence, for each $x \in \Hcal_{RN}^{\leq d}$,
	$\mu(x) \cdot f(x) \geq \zeta(x) \cdot f(x) - s/2 > s/2$, where the final inequality holds by \expeqref{Property}{zetaprop5}. That is, the correction terms do not disturb the sign of any points in their support.
	
	\medskip
	\noindent \textbf{Analysis of $\mu$: Details.} 
	We need to prove that $\mu$ satisfies the conditions of \expref{Theorem}{theorem:tdegdual} with $f=(\AND_R \circ \OR_N)^{\leq N + d}$. 
	
	\begin{itemize}
	\item \textbf{(Pure high degree).} That $\mu$ has pure high degree at least $d$ follows from \expref{Fact}{sumfact} %
and the fact that $\zeta$ and each $\psi_{\mathsf{corr}, x}$ have pure high degree at least $d$	(\expeqref{Equations}{zetaprop2} \expeqref{and}{eqn:rs-phd}). 
		\item \textbf{(Sign-agreement).}
		To establish sign-agreement, we consider three cases:
		\begin{itemize}
			\item $y \in B$. In this case $\mu(y) = \zeta(y) - \psi_{\mathsf{corr}, y}(y) = \zeta(y) - \zeta(y) = 0$. Here,
the first equality holds because $B \subseteq \mathcal{H}_{RN}^{> d}$ 
(see \expeqref{Equation}{subseteq}), and for each
			$x \in B \setminus \{y\}$, the point $y$ is not in the support of $\psi_{\mathsf{corr}, x}$ (as \expeqref{Equation}{eqn:rs-hamming} states that the support of $\psi_{\mathsf{corr}, x}$ is a subset of $\mathcal{H}_{RN}^{\leq d} \cup \{x\}$).
			\item $y \not \in (B \cup \mathcal{H}^{\leq d}_{RN})$. In this case, $\mu(y) \cdot f(y) = \zeta(y) \cdot f(y) \geq 0$. The first equality holds 
			because $y \not \in B$, and since $y \not \in \mathcal{H}^{\leq d}_{RN}$, $y$ is not in the support of any $\psi_{\mathsf{corr}, x}$ for any $x \in B$ (see 
			\expeqref{Equation}{eqn:rs-hamming}). 
			\item $y \in \mathcal{H}^{\leq d}_{RN}$. In this case, \begin{align} 
			\mu(y) = \zeta(y) - \sum_{x \in B} \psi_{\mathsf{corr}, x}(y) \geq \zeta(y) - \left|\sum_{x \in B} \psi_{\mathsf{corr}, x}(y)\right|
			\geq \zeta(y) - \sum_{x \in B} \|\psi_{\mathsf{corr}, x}\|_1 \notag\\ \label{forsignrank} \geq
			\zeta(y) - \left(E\left(\zeta, \AND_R \circ \OR_N\right) +W\left(\zeta, f\right)\right) \cdot 2^d \cdot {RN \choose d} > s-s/2 >0\,.\end{align}
			
			Here, the third to last inequality follows from the definition of $\psi_{\mathsf{corr}, x}$ and \expeqref{Expression}{eqn:rs-norm}, and the penultimate
			inequality follows by \expeqref{Expressions}{zetaprop5} \expeqref{and}{akeyeq}.
			
		\end{itemize} 
		\item \textbf{(Appropriate support).}
		The first case considered in the sign-agreement analysis above also establishes the appropriate support condition, since it shows that $\mu(x)=0$ for all $x \in B$, and by definition of $B$, $\mathcal{H}^{> N+d}_{RN} \subseteq B$.
		\item \textbf{(Non-triviality).}
		The third case considered in the sign-agreement analysis also establishes non-triviality, since $\mu(x) > 0$ for all $x \in \mathcal{H}^{\leq d}_{RN}$. 
	\end{itemize}
	
	By \expref{Theorem}{theorem:tdegdual}, $\mu$ witnesses the fact that for $R=N^{1/2}$,
%
%
\[\deg_{\pm}\left(\left(\AND_R \circ \OR_N\right)^{\leq N+d}\right) =
\Omega\left(N^{1/2}/\log^{3/2} N\right)\,.\]
For any $t > 0$,
%
\[ \deg_{\pm}\left(\left(\AND_R \circ \OR_N\right)^{\leq t}\right) \]
is clearly nondecreasing with $N$. Hence, 
%
%
\[ \deg_{\pm}\left(\left(\AND_R \circ \OR_{N+d}\right)^{\leq N+d}\right) = \Omega\left(N^{1/2}/\log^{3/2} N\right) \,.\]
Since $d=o(N)$, setting $N'=N+d$ implies that
%
%
\[ \deg_{\pm}\left(\left(\AND_R \circ \OR_{N'}\right)^{\leq N'}\right)
 = \Omega\left((N')^{1/2}/\log^{3/2}(N')\right)\,,\]
as desired.  %
\end{proof}


\section{The large-error approximate degree of
  \texorpdfstring{AC$^0$}{AC0} is nearly linear}
\label{s:discproof}
\begin{theorem} \label{ac0thm}
	For any constant $\delta > 0$, there is a function $h \colon \bits^n \to \bits$ computed
	by an AC$^0$ circuit of depth $O(1/\delta)$ such that $\adeg_{\eps}(h) = \tilde{\Omega}(n^{1-\delta})$,
	for some $\eps=1-2^{-\tilde{\Omega}(n^{1-\delta})}$. 
\end{theorem}

\expref{Theorem}{ac0thm} is a consequence of the following hardness amplification result, which
shows how to transform any function $f$ that is hard to approximate by low-degree polynomials
into a function $F$ that is even harder to approximate, in terms of \emph{both} the degree
and the error that is achievable by polynomials of said degree.

\begin{theorem} \label{hardnessamp}
	Fix constants $k\geq 1, c \ge 0$. Let $f_n \colon \bits^{n\log^c n} \to \bits$ be any (infinite family of) functions satisfying $\adeg_{\eps}(f_n) \geq \Omega\left(n^{k/(k+1)}\right)$
	for some $\eps=1-2^{-\Omega(n^{k/(k+1)})}$ and sufficiently large $n$. Then there is an explicitly given sequence of functions
	$F_n \colon \bits^{n \log^{c+O(1)}n} \to \bits$ such that 
	$\adeg_{\eps}(F_n) \geq \Omega\left(n^{(k+1)/(k+2)}\right)$
	for some $\eps=1-2^{-\Omega\left(n^{(k+1)/(k+2)}\right)}$ and sufficiently large $n$. Moreover,
	if $f_n$ is computed by polynomial-size %
circuits of depth %
$\Delta$, then $F_n$ is computed by 
polynomial-size %
circuits of depth at most $\Delta+O(1)$. 
\end{theorem}

\begin{proof}[Proof of Theorem \ref{ac0thm} assuming Theorem \ref{hardnessamp}]
$\SURJ_{N^{1/2}, N}$ is a function on $n=\Theta(N \log N)$ input bits,
%
computed by a polynomial-size %
circuit of depth 3, and 
	\expref{Corollary}{surjcor} implies that $\deg_{\pm}(\SURJ_{N^{1/2}, N}) = \tilde{\Omega}(N^{1/2})$.
	Applying \expref{Theorem}{hardnessamp} to $f=\SURJ_{N^{1/2}, N}$ with $k=1$ yields a function $F_1$ 
	on $n \cdot \polylog(n)$ inputs that is computed by a 
polynomial-size %
circuit of depth $3+O(1)=O(1)$, and satisfies $\adeg_{\eps}(F_1) \geq \Omega(n^{2/3})$
	for some $\eps=1-2^{-\Omega(n^{2/3})}$.
	Applying \expref{Theorem}{hardnessamp} yet again, with $f=F_1$ and $k=2$
	yields a function $F_2$ in AC$^0$ that is also defined on $n \cdot \polylog(n)$ inputs, is computed
by a polynomial-size %
circuit of depth $O(1)$, and satisfies 
	$\adeg_{\eps}(F_2) \geq \Omega(n^{3/4})$
	for some $\eps=1-2^{-\Omega(n^{3/4})}$. %
	
	In general, for any constant $\delta \in (0, 1)$, let $k$ be a constant such that $1-\delta \le k/(k+1)$,
	\ie, $k \ge 1/\delta-1$. 
	Then iteratively applying \expref{Theorem}{hardnessamp} $k-1$ times, starting with $f=\SURJ_{N^{1/2}, N}$,
	yields a function $h$ computed by an AC$^0$ circuit of depth $O(k)=O(1/\delta)$, defined on $n' = n \log^{O(k)} n = \tilde{O}(n)$ input bits,
	such that $\adeg_{\eps}(h) = \Omega(n^{1-\delta})$
	for some $\eps=1-2^{-\Omega(n^{1-\delta})}$. \expref{Theorem}{ac0thm} follows.
\end{proof}

%
\begin{proof}[Proof of Theorem \ref{hardnessamp}] 
	\expref{Theorem}{thm:main-reduction} implies that, in order to prove \expref{Theorem}{hardnessamp}, 
	it is sufficient to identify a block-composed function $G$ defined on $\poly(n)$ inputs that is computed by a 
polynomial-size %
circuit of depth $\Delta+O(1)$ such that for some $\ell=n\cdot\polylog(n)$, we have $\deg_{\eps}(G^{\leq \ell}) \geq \Omega(n^{(k+1)/(k+2)})$
	for some $\eps=1-2^{-\Omega\left(n^{(k+1)/(k+2)}\right)}$. Indeed, if we accomplish this,
	then \expref{Theorem}{hardnessamp} follows by setting $F=G^{\mathsf{prop}}$.
	
	To define $G$, we need the following lemma, which follows from the techniques of~\cite{AB84} (see~\cite{Kop13} for an exposition).
	
	\begin{lemma}\label{lem:gapmaj}
  There exists a Boolean circuit $C_n \colon \bits^n \to \bits$ with $n$ 
inputs, depth 3, and size $\tO(n^2)$ satisfying the following two
conditions:
		\begin{itemize}
	\item $C_n(x) = 1$ for all $x$ of Hamming weight at most $n/3$.
	\item $C_n(x) = -1$ for all $x$ of Hamming weight at least $2n/3$.
		\end{itemize}
	\end{lemma}

	We refer to the function computed by the circuit $C$ of \expref{Lemma}{lem:gapmaj} as $\mathsf{GAPMAJ}$, short for a gapped majority function (such a function is sometimes also called an \emph{approximate majority} function).\footnote{In prior related work~\cite{bchtv}, $\mathsf{GAPMAJ}$ referred
		to the \emph{promise} function that equals $1$ for all inputs $x$ of Hamming weight at most $n/3$, equals  $-1$ for all inputs $x$ of Hamming weight at least $2n/3$,
		and is undefined otherwise. In contrast, we use $\mathsf{GAPMAJ}$ to refer to any total function in AC$^0$ that agrees with the partial function
		from~\cite{bchtv} at all points in the partial function's domain.} We remark that while the circuit $C$ from \expref{Lemma}{lem:gapmaj} is not explicitly constructed,
explicit constructions of AC$^0$ circuits satisfying the two bulleted
conditions 
of \expref{Lemma}{lem:gapmaj} are known~\cite{viola2009approximate}.
	
	\medskip
	\noindent \textbf{Definition of $G$.} Let $t=n^{1/(k+2)}$, $z=n^{(k+1)/(k+2)}$, $r=10 \log n$, and $m=n^{2/(k+2)}$. Let $N= t \cdot Z \cdot r \cdot M$, where $Z = z \log ^c z$, and $M = 10n\log^{c+4}n$.
	We define $G \colon \bits^N \to \bits$ to equal $\GAPMAJ_t \circ f_z \circ \AND_r \circ \OR_M$. Here, $f_z$ denotes 
	the function $f$ on $Z = z \log^c z$ variables whose existence is assumed by the hypothesis of the theorem.
	
	
	\medskip
	We now begin the process of constructing a dual witness $\mu$ showing that $\adeg_{\eps}(G^{\leq 10 n \log^{c + 4} n})=\Omega\left(n^{(k+1)/(k+2)}\right)$
	for some $\eps=1-2^{-\Omega\left(n^{(k+1)/(k+2)}\right)}$.
	
	For any appropriate constant $c > 0$, let $\phi' \colon \bits^Z \to \R$ be a dual witness for the fact that $\adeg_{\eps}(f_z) \geq C \cdot z^{k/(k+1)} = C \cdot n^{k/(k+2)}\coloneqq d$ (for some constant $C$) as per \expref{Theorem}{theorem:adegdual}.
	Then $\phi \coloneqq  \phi'/\|\phi\|_1$ has the following properties. %
	
	\begin{equation}
	\label{phiprop1}
	\|\phi\|_1 =1\,.
	\end{equation}
	\begin{equation}
	\label{phiprop2}
	\phi \text{ has pure high degree at least } d\,.
	\end{equation}
	\begin{equation}
	\label{phiprop3}
	\sum_{x\in\bits^Z} \phi(x) \cdot f_z(x) \geq 1-\delta' \text{ for some } \delta'=2^{-\Omega(d)}\,.
	\end{equation}
	
	Similar to the proof of \expref{Theorem}{surjthm}, let $\psi$ be the dual witness for $\OR_M$ from \expref{Proposition}{prop:or-dual} with $M=10n\log^{c+4}n$, $T = m$, and $\delta=1/4$,
	and let 
	$$\Phi_r(w) = \begin{cases}
	-1/2 & \text{ if } w=(-1, -1, \dots, -1)\\
	1/2 & \text{ if } w= (1, 1, \dots, 1)\\
	0 & \text{otherwise}.
	\end{cases}$$
	
	We remind the reader %
	that \expeqref{Equations}{btprop1}-\eqref{btprop3} from the proof of \expref{Theorem}{surjthm} guarantee %
	 that $\Phi_r \ls \psi$ has
	the following properties.
	\begin{equation}  \label{btprop1repeat}
	\|\Phi_r \ls \psi\|_1=1\,.
	\end{equation}
	%
	\begin{equation}  \label{btprop2repeat}
	\Phi_r \ls \psi \text{ has pure high degree at least that of } \psi \text{, which is } D'\coloneqq \Omega\left(m^{1/2}\right)=\Omega\left(n^{1/(k+2)}\right)\,.
	\end{equation}
	%
	\begin{equation}  \label{btprop3repeat}
	E\left(\Phi_r \ls \psi, \AND_r \circ \OR_M\right) \leq 2^{-r}=1/n^{10}\,.
	\end{equation}

	We now combine $\phi$ and $\Phi_r \ls \psi$ to obtain an intermediate dual witness $\gamma$, 
	which will witness the fact that $\adeg_{1-2^{-\Omega(d)}}(f_z \circ \AND_r \circ \OR_M) = \Omega(d \cdot \sqrt{m})$.
	The combining technique that we use is precisely the one introduced by Sherstov \cite[Theorem 6.1]{sherstovdp} 
	to establish a direct-sum theorem for approximate degree. Roughly speaking, \cite[Theorem 6.1]{sherstovdp}
	showed that for any Boolean functions $f$ and $g$, if $\adeg_{\delta_1}(f) \geq d$ and $\adeg_{1/n}(g) \geq d'$, then $\adeg_{\delta_2}(f \circ g) = \Omega(d \cdot d')$
	for some $\delta_1 \approx \delta_2$. The $\gamma$ that we construct below is 
	(a normalized version of) the dual witness constructed in the proof of \cite[Theorem 6.1]{sherstovdp} when
	applied with outer function $f_z$ and inner function $g=\AND_r \circ \OR_M$. 
	
	For a fixed \emph{even} integer $j \in (d/2-2, d/2]$, let $P_j$ be the degree $j$ univariate polynomial given by $P_j(a) = \prod_{i=1}^j (a-i)$,
	and define $p_j \colon \bits^{Z} \to \R$ to
	be the unique multilinear polynomial such that $p_j(x)= P_j(|x|)$. 
	We will need the following properties of $p_j$.
	%
	\begin{equation}
	\label{fourierpj} 
	\text{The sum of the absolute values of the Fourier coefficients of } p_j \text{ is at most } j! {Z + j \choose j}\,.
	\end{equation}
	%
	\begin{equation} \label{pjeq}
	\text{For any } c \in [1-1/Z^2, 1], \text{ it holds that } p_j(c, \dots, c) \geq \frac{1}{2}  j! \geq 1\,.
	\end{equation}
	
	\expeqref{Property}{fourierpj} is precisely Lemma 3.1 (specifically, Property 3.2) of~\cite{sherstovdp}.
	To see that \expeqref{Property}{pjeq} holds, first observe that $p_j(x) \geq 0$ for all $x \in \bits^{Z}$,
	owing to the fact that $j$ is even. Second, by the multilinearity of $p_j$, the value of $p_j(c, \dots, c)$ is the expected
	value of $p_j$ under the distribution over $x \in \bits^{Z}$ in which each coordinate $i$ of $x$ is chosen independently
	from $\{-1, 1\}$ such that the expected value of $x_i$ is $c$. Denoting this distribution as $B_c$, $\Pr_{x \sim B_c}[x = \mathbf{1}_{Z}] \geq 1/2$.
Hence, %
\begin{equation}
 p_j(c, \dots, c) \geq \Pr_{x \sim B_c}[x = \mathbf{1}_{Z}] \cdot p_j(\mathbf{1}_{Z}) \geq \frac12  j!\,.
\end{equation}	
	\medskip
	 
	Let $\delta''=2 \cdot E(\AND_r \circ \OR_M, \Phi_r \ls \psi) \leq 2 \cdot 2^{-r} = 2/n^{10}$.
	For $y \in \bits^{r \cdot M}$, define\footnote{Our definition of $\alpha$ specializes the definition in \cite[Page 39]{sherstovdp}
		to our setting: our definition and the definition in \cite[Page 39]{sherstovdp} coincide owing to the fact that (as established in \cite[Theorem 1]{bt14})
		$(\Phi_r \ls \psi)(y) < 0 \Longrightarrow \left(\AND_r \circ \OR_M\right)(y) < 0$.}:
	$$\alpha(y)=\begin{cases} 1 & \text { if } (\AND_r \circ \OR_M)(y) = \sgn((\Phi_r \ls \psi)(y)) = +1\\
	1-2\delta'' & \text { if } (\AND_r \circ \OR_M)(y) = \sgn((\Phi_r \ls \psi)(y)) = -1\\
	-1 & \text{ otherwise.} \end{cases}$$

	For an input $x$ to $f_z \circ \AND_{r} \circ \OR_M$, 
	write $x=(x_1, \dots, x_Z)$ with each $x_i \in \bits^{r \cdot M}$. 
	Define:
	$$\gamma'(x_1, \dots, x_Z) = p_j\left(\alpha(x_1), \dots, \alpha(x_Z)\right) \cdot \left( \phi \ls \Phi_r \ls \psi\right) (x),$$
	and let $\gamma = \gamma'/\|\gamma'\|_1$.
	
	We observe that $\gamma'$ has the following properties. %
	
	\begin{equation}
	\label{gammaprimeprop1}
	\|\gamma'\|_1 = p_j(1-2\delta'', \dots, 1-2\delta'') > 0\,.
	\end{equation}
	%
	\begin{equation}
	\label{gammaprimeprop2}
	\gamma' \text{ has pure high degree at least } (d/2) \cdot D'\,.
	\end{equation}
	\begin{equation}
	\label{gammaprimeprop3}
	\sum_{x \in \bits^{Z \cdot r \cdot M}} \gamma'(x) \cdot \left(f_z \circ \AND_r \circ \OR_M\right)(x) \geq p_j\left(\dots, 1-2\delta'', \dots\right) \cdot \left( (1-2\delta') - \frac{2 (\delta'')^{j+1}}{(1-\delta'')^n} \cdot {n \choose j+1}\right)\,.
	\end{equation}
	%
	\expeqref{Property}{gammaprimeprop1} is precisely Claim 6.2 from~\cite{sherstovdp}.
	\expeqref{Equation}{gammaprimeprop2} is immediate from \cite[Equation 6.7]{sherstovdp}, combined
with the fact that $\Phi_r \ls \psi$ has pure high degree at least $D'$
     (\expeqref{Equation}{btprop2repeat})
	and $\phi$ has pure high degree at least $d$ 
	(\expeqref{Equation}{phiprop2}). 
	\expeqref{Property}{gammaprimeprop3} is precisely Claim 6.3 from~\cite{sherstovdp}. 
	
	Recall that we defined $\gamma = \gamma'/\|\gamma'\|_1$. \expeqref{Properties}{gammaprimeprop1}-\eqref{gammaprimeprop3} imply that $\gamma$ 
satisfies the following conditions: %
	
	\begin{equation}
	\label{gammaprop1}
	\|\gamma\|_1 = 1\,.
	\end{equation}
	\begin{equation}
	\label{gammaprop2}
	\gamma \text{ has pure high degree at least } (d/2) \cdot D'\,.
	\end{equation}
	\begin{align}
	\sum_{x \in \bits^{Z \cdot r \cdot M}} \gamma(x) \cdot \left(f_z \circ \AND_r \circ \OR_M\right)(x) &\geq 
	(1-2\delta') - \frac{2 (\delta'')^{j+1}}{(1-\delta'')^n} \cdot {n \choose j+1} \notag\\ &\geq 
	\left(1-2\delta' - 2^{-3d}\right) = 1-2^{-\Omega(d)}\,.
	\label{gammaprop3}
	\end{align}
	In \expeqref{Property}{gammaprop3}, the penultimate inequality holds because $\delta'' \leq 2 \cdot 2^{-r} =2/n^{10}$, and $j > d/2-2$. 
	
	\medskip \noindent \textbf{The final steps in the construction of $\mu$.}
	Next, define $\zeta \colon \bits^{ t \cdot Z \cdot r \cdot M} \to \R$ via:
	$$\zeta = \Phi_t \ls \gamma.$$  
	
	For every $x \in \mathcal{W}\left(\zeta, G^{\leq 10 n \log^{c + 4} n}\right)$, invoke \expref{Lemma}{lem:rs} with $D=C' \cdot n^{(k+1)/(k+2)}$ (for a sufficiently
small constant $C'$ to be chosen later) to obtain a function $\phi_x$
satisfying  \expeqref{Conditions}{eqn:rs-point}-\eqref{eqn:rs-norm}, and
define $\psi_{\mathsf{corr}, x}
\coloneqq %
\zeta(x) \cdot \phi_x$\,.
	We define our final dual witness to be $\mu \coloneqq  \zeta - \sum_{x \in \mathcal{W}\left(\zeta, G^{\leq 10 n \log^{c + 4} n}\right)} \psi_{\mathsf{corr}, x}$\,.
	
\medskip \noindent \textbf{Analysis of $\zeta$ and $\mu$.} %
%
%
%
We observe that $\zeta$ has the following properties.
	
	\begin{equation}\label{zeta2prop1}
	\|\zeta\|_1=1\,.
	\end{equation}
	
	\begin{equation}\label{zeta2prop2}
	\zeta \text{ has pure high degree at least that of } \gamma \text{, which is  at least } D''\coloneqq (d/2) \cdot D'\,.
	\end{equation}
	
	\begin{equation}\label{zeta2prop3}
	\sum_{x\in\bits^{t \cdot Z \cdot r \cdot M}} \zeta(x) \cdot \left(\GAPMAJ_t \circ f_z \circ \AND_{r} \circ \OR_M\right)  \geq 1-2^{-\Omega(d\cdot t)}\,.
	\end{equation}
	
	\begin{equation} \label{zeta2prop4}
	W(\zeta, G^{\leq 10 n \log^{c + 4} n}) \leq 2^{-\omega(n^{(k+1)/(k+2)} \log n)}\,.
	\end{equation}
	
	%
	\medskip \noindent \textbf{Justification for \expeqref{Properties}{zeta2prop1}-\eqref{zeta2prop4}.}
	\expeqref{Equation}{zeta2prop1} follows from \expref{Proposition}{prop:ls} (see  \expeqref{Equation}{eqn:ls-norm}), and 
	the facts that $\|\Phi_t\|=1$ and $\|\gamma\|=1$ (see  \expeqref{Equation}{gammaprop1}), and $\gamma$ has pure high degree at least 1.
	\expeqref{Equation}{zeta2prop2} follows from \expref{Proposition}{prop:ls} (see  \expeqref{Equation}{eqn:ls-phd}), 
	the fact that $\Phi_t$ has pure high degree 1, and the fact that $\gamma$ has pure high degree at least $(d/2) \cdot D'$ (see  \expeqref{Equation}{gammaprop2}).
	The validity of \expeqref{Inequality}{zeta2prop3} is established via a standard, but somewhat lengthy, analysis that we
	defer to \expref{Section}{app:zeta2prop3}.
	%
	
	We justify \expeqref{Inequality}{zeta2prop4} as follows. Denote
	an input $x$ in $\bits^{t \cdot Z \cdot r \cdot M}$ as $x=(\dots, x_{i, s}, \dots)$ 
	where $i$ ranges over $1, \dots, t$, $s$ ranges over $1, \dots, Z$,
	and each $x_{i,s} \in \bits^{r \cdot M}$.  Since $p_j$ is non-negative on all inputs in $[-1, 1]^Z$ (this follows from the same reasoning as in the proof of \expeqref{Property}{pjeq}),
	\begin{align}
	\left|\zeta(x)\right| &= \left|\left(\Phi_t \ls \phi \ls \Phi_r \ls \psi\right)(x) \cdot \left(\|\gamma'\|_1\right)^{-t}  \cdot \prod_{i=1}^{t} p_j\left(\alpha(x_{i, 1}), \dots, \alpha(x_{i, Z})\right) \right| \notag \\
	\notag &=  \left|\left(\Phi_t \ls \phi \ls \Phi_r \ls \psi\right)(x) \cdot p_j(\dots, 1-2\delta'', \dots)^{-t} \cdot \prod_{i=1}^{t} p_j(\alpha(x_{i, 1}), \dots, \alpha(x_{i, Z}))\right|\\
	\notag &\leq  \left|\left(\Phi_t \ls \phi \ls \Phi_r \ls \psi\right)(x) \cdot \prod_{i=1}^{t} p_j(\alpha(x_{i, 1}), \dots, \alpha(x_{i, Z}))\right|\\
	&\leq \left| \left(\Phi_t \ls \phi \ls \Phi_r \ls \psi\right)(x) \right| \left(j! {Z+j \choose j}\right)^t \notag \\
	&\leq \left| \left(\Phi_t \ls \Phi_r \ls \psi\right)(x) \right| \cdot Z^{O(jt)} \notag\\
	&\leq 
	\left| \left(\Phi_t \ls \phi \ls \Phi_r \ls \psi\right)(x) \right| \cdot 2^{O(d t \log n)}\,. \label{fineq} 
	\end{align}
	Here, the second equality holds by \expeqref{Equation}{gammaprimeprop1}, the first inequality holds by \expeqref{Property}{pjeq}, and the second inequality follows from \expeqref{Property}{fourierpj} and the fact that $|\alpha(x_{i, s})| \leq 1$ for all $i=1, \dots, t$ and $s=1, \dots, Z$. 
	
	Now invoke \expref{Lemma}{prop:prelim-mass} (and associativity of dual block composition, see 
	 \expeqref{Equation}{eqn:ls-assoc}) to conclude that 
	\begin{align} \label{reallyfineq} W\left( \left(\Phi_t \ls \phi \ls \Phi_r\right) \ls \psi, \left(\GAPMAJ_t \circ f_z \circ \AND_r \circ \OR_M\right)^{\leq 10 n \log^{c+4} n}\right) \notag \leq\\
	2^{-\Omega\left(n \log^2 n / \sqrt{m}\right)} \leq 2^{-\Omega(n^{(k+1)/(k+2)} \log^2 n)}\,.\end{align}
	
	Combining \expeqref{Inequalities}{fineq} \expeqref{and}{reallyfineq}, we conclude that
	\begin{align*} W\left( \zeta, \left(\GAPMAJ_t \circ f_z \circ \AND_r \circ \OR_M\right)^{\leq 10 n \log^{c+4} n}\right) &
	\leq 2^{-\Omega\left(n^{(k+1)/(k+2)} \log^2 n\right)} \cdot 2^{O\left(d t \log n\right)} \\
	&\leq 2^{-\Omega\left(n^{(k+1)/(k+2)} \log^2 n\right)} \cdot 2^{O\left(n^{(k+1)/(k+2)} \log n\right)}\\
	&\leq 2^{-\Omega\left(n^{(k+1)/(k+2)} \log^2 n\right)}\,,
	\end{align*}
	where the final inequality holds for large enough $n$. 
	
%
%
%
%
Finally, we observe the following properties of $\mu$. %
	
	\begin{equation}\label{muprop1}
	\|\mu\|_1\leq 1 + 2^{-\omega(n^{(k+1)/(k+2)})}\,.
	\end{equation}
	
	\begin{equation}\label{muprop2}
	\mu \text{ has pure high degree at least }  \min(D'', D) = \Omega\left(n^{(k+1)/(k+2)}\right)\,.
	\end{equation}
	
	\begin{equation}\label{muprop3}
	\sum_{x\in\bits^{t \cdot Z \cdot r \cdot M}} \mu(x) \cdot G(x)  \geq 1-2^{-\Omega(t \cdot d)} = 1-2^{-\Omega\left(n^{(k+1)/(k+2)}\right)}\,.
	\end{equation}
	
	\begin{equation}\label{muprop4}
	\mu(x)=0 \text{ for all } x \in \mathcal{H}_{t \cdot Z \cdot r \cdot M}^{> 10 n \log^{c+4} n}\,.
	\end{equation}
	
	\medskip \noindent \textbf{Justification for \expeqref{Properties}{muprop1}-\eqref{muprop4}.}
	For \expeqref{Inequality}{muprop1}, observe that 
	\begin{align*}\|\mu\|_1 &\leq \|\zeta\|_1 + \sum_{x \in \mathcal{W}(\zeta, G^{\leq 10 n \log^{c+4} n})} \|\psi_{\mathsf{corr}, x}\|_1\\
	&\leq 1 + W(\zeta, G^{\leq 10 n \log^{c+4} n}) \cdot 2^D \cdot {t \cdot Z \cdot r \cdot M \choose D}  \\&\leq
	1 + 2^{-\omega\left(n^{(k+1)/(k+2)} \log n\right)}\,,
	\end{align*}
	where the penultimate inequality holds by \expeqref{Expressions}{eqn:rs-norm} \expeqref{and}{zeta2prop1}, and the final inequality holds by \expeqref{Inequality}{zeta2prop4}.
	
\expeqref{Equation}{muprop2} follows from \expref{Fact}{sumfact}, \expeqref{Equation}{zetaprop2}, and the fact that each term $\psi_{\mathsf{corr}, x}$ has pure high degree at least $D$ (\expeqref{Equation}{eqn:rs-phd}).
	
	\expeqref{Expression}{muprop3} holds because \begin{align*} \sum_{x\in\bits^{t \cdot Z \cdot r \cdot M}} \mu(x) \cdot G(x)
	&\geq \sum_{x\in\bits^{t \cdot Z \cdot r \cdot M}} \zeta(x) \cdot G(x) - \sum_{x \in \mathcal{W}\left(\zeta, G^{\leq 10 n \log^{c+4} n}\right)} \|\psi_{\mathsf{corr},x}\|_1\\
	&\geq 1 - 2^{-\Omega(td)} -  2^{-\omega(n^{(k+1)/(k+2)})} \geq 1 - 2^{-\Omega(td)}\,,\end{align*}
	where the penultimate inequality invokes \expeqref{Expressions}{eqn:rs-norm} \expeqref{and}{zeta2prop4}.
	
	\expeqref{Expression}{muprop4} holds because for any $x \in \mathcal{H}_{t \cdot Z \cdot r \cdot M}^{> 10 n \log^{c+4} n}$, we have $\mu(x) = \zeta(x) - \psi_{\mathsf{corr}, x}(x) = \zeta(x) - \zeta(x) = 0$. Here, the first equality holds because \expeqref{Property}{eqn:rs-hamming} ensures that, for any $x' \in \mathcal{W}\left(\zeta, G^{\leq 10 n \log^{c+4} n}\right)$ such that $x \neq x'$, we have $\psi_{\mathsf{corr}, x'}(x)=0$. 
	
	It is immediate from \expeqref{Properties}{muprop1}-\eqref{muprop4} that $\mu$ %
	satisfies the conditions 
	 required by \expref{Theorem}{theorem:adegdual}
	to witness the fact that $\adeg_{\eps}(G^{\leq 10 n \log^{c+4} n})  = \Omega(n^{(k+1)/(k+2)})$ for some $\eps=1-2^{-\Omega(n^{(k+1)/(k+2)})}$. 
\end{proof}

\subsection{Justification for \expeqref{Expression}{zeta2prop3}}
\label{app:zeta2prop3}
All notation in this section the same as in the proof of \expref{Theorem}{hardnessamp}.
Let $F=f_z \circ \AND_r \circ \OR_M$.
Recall that $\zeta = \Phi_t \ls \gamma,$ where $\gamma$ 
has the following properties.
\begin{equation} \|\gamma\|_1=1\,.\end{equation}
\begin{equation} \gamma
\text{ has pure high degree at least } (d/2)\cdot D' \geq 1\,.\label{equationtwotwotwo}\end{equation}
\begin{equation} \sum_{y \in \bits^{Z \cdot r \cdot M}} \gamma(y) \cdot F(y) \geq \eps \coloneqq 1-2^{-\Omega(d)}\,.\label{killmenow}\end{equation}
\expeqref{Equation}{equationtwotwotwo} implies that 
\begin{equation}
\label{lame} \sum_{y \colon \sgn(\gamma(y))=1} |\gamma(y)| = \sum_{y \colon \sgn(\gamma(y))=-1} |\gamma(y)|=1/2\,.\end{equation}
Let $E_1(\gamma, F) = \sum_{y \colon \gamma(y) > 0 \text{ and } F(y) < 0} |\gamma(y)|$ and $E_{-1}(\gamma, F) = \sum_{y \colon \gamma(y) < 0 \text{ and } F(y) > 0} |\gamma(y)|$.
Combining \expeqref{Expressions}{killmenow} \expeqref{and}{lame}
implies that $E_{1}(\gamma, F)$ and $E_{-1}(\gamma, F) \leq 2^{-\Omega(d)}$. 

Let $\tau$ be the product distribution on $\left(\bits^{Z \cdot r \cdot M}\right)^t$
given by $\tau(x_1, \dots, x_t) = \prod_{i=1}^{t} |\gamma(x_i)|$. 
Given a string $b=(b_1, \dots, b_t) \in \bits^t$, let $\tau_b$ be the distribution over $x \in \bits^{t \cdot Z \cdot r \cdot M}$
equal to $\tau$ conditioned on $(\dots, \sgn(x_i), \dots, ) = b$. 
\expeqref{Equation}{lame} implies that when $x=(x_1, \dots, x_t)$ is drawn from $\tau$, the string $(\dots, \sgn(\gamma(x_i), \dots)$
is uniformly distributed in $\bits^t$. 

Moreover,
for any given $b \in \{-1, 1\}^{t}$, the following two random variables are identically distributed:

\begin{itemize}
	\item The string $(\dots, F(x_i), \dots)$ when one chooses $(\dots, x_i, \dots)$ from the conditional distribution 
	$\tau_b$. 
	\item The string $(\dots, \delta_ib_i, \dots)$, where $\delta \in \{-1, 1\}^{t}$ is a random string whose $i$th bit independently
	takes on value $-1$ with probability $2 \sum_{x \in E_{b_i}} |\psi(x)| < 2^{-\Omega(d)}$. 
\end{itemize}



Thus,
\begin{align*}
  \sum_{x \in \bits^{t \cdot Z \cdot r \cdot M}} \zeta(x) \cdot \GAPMAJ(\dots, F(x_i), \dots) %
%
\\
=
2^{t} \cdot \sum_{b \in \bits^t} \Phi_t(b) \cdot \mathbb{E}_{x \sim \tau_{b}}[\GAPMAJ(\dots, F(x_i), \dots)]
\\
=
2^t \left( \frac{1}{2} \mathbb{E}_{x \sim \tau_{\mathbf{1}_t}}[\GAPMAJ(\dots, F(x_i), \dots)]  - \frac{1}{2} \mathbb{E}_{x \sim \tau_{-\mathbf{1}_t}}[\GAPMAJ(\dots, F(x_i), \dots)\right) \\
=
\frac12 \mathbb{E}_{\delta^{(1)}}[\GAPMAJ_t(\delta^{(1)}_1, \dots, \delta^{(1)}_t)] - \frac12  \mathbb{E}_{\delta^{(-1)}}[\GAPMAJ_t(-\delta^{(-1)}_1, \dots, -\delta^{(-1)}_t)]\,,
\end{align*}
where for $j \in \bits$, $\delta^{(j)} \in \bits^t$ is a random string whose $i$th coordinate takes value $-1$ with probability $2 \sum_{x \in E_{j}} |\psi(x)| < 2^{-\Omega(d)}$. 

Note that $\mathbb{E}_{\delta^{(1)}}[\GAPMAJ_t(\delta^{(1)}_1, \dots, \delta^{(1)}_t)] \geq 1-2^{-\Omega(td)}$:
since each bit of $\delta^{(1)}$ equals $1$ with probability $1-2^{-\Omega(d)}$, the probability that more than $t/3$
of the coordinates of $\delta^{(1)}$ are $-1$ is at most ${t \choose t/3} \cdot 2^{-\Omega(td)} = 2^{-\Omega(td)}$.
Similarly, $\mathbb{E}_{\delta^{(-1)}}[\GAPMAJ_t(-\delta^{(-1)}_1, \dots, -\delta^{(-1)}_t)] \leq -1+2^{-\Omega(td)}$. 
We conclude that the correlation of $\zeta$ with $\GAPMAJ_t \circ F$ is at least $1-2^{-\Omega(td)}$ as claimed.

\section{Algorithmic and complexity-theoretic applications}
\label{s:apps}
\subsection{Complexity measures}
Throughout this section, for a function $F \colon \bits^n \times \bits^n \to \bits$,
we also view $F$ as a $2^n \times 2^n$ matrix whose $(x, y)$'th entry is given by $F(x, y)$.

\medskip \noindent \textbf{Approximate rank.} For a matrix $F \in \bits^{N \times N}$, 
the $\eps$-approximate rank of $F$, denoted $\text{rank}_{\eps}(F)$, is the least rank
of a matrix $A \in \R^{N \times N}$ such that $|A_{ij} - F_{ij}| \leq \eps$ for all $(i, j) \in [N] \times [N]$.
Approximate rank is a fundamental notion in learning theory and communication complexity.
Meanwhile, the \emph{sign-rank} of $F$ is the least rank of a matrix  $A \in \R^{N \times N}$ such that $|A_{ij} - F_{ij}| < 1$ for all $(i, j) \in [N] \times [N]$.
Clearly, for \emph{any} $F$, the \emph{exact} (not just approximate) rank of $F$ is at most $N$. 

%

\medskip \noindent \textbf{Discrepancy.} The discrepancy of $F$, denoted
$\mathrm{disc}(F)$, is a combinatorial measure of the complexity of $F$ that roughly
captures $F$'s correlation with combinatorial rectangles (small discrepancy
corresponds to high complexity).

\begin{definition} \label{def:disc}
	A combinatorial rectangle of $X \times Y$ is a set of the form $A \times B$ with $A \subseteq X$ and $B \subseteq Y$. 
	For a distribution $\mu$ over $X \times Y$, the discrepancy of $F$ with respect to $\mu$ is defined to be the maximum over all rectangles $R$ of the \emph{bias} of $F$ on $R$. That is:
	$$\mathrm{disc}_\mu(F) = \max_{R} \left|\sum_{(x, y) \in R} \mu(x, y) F(x, y)\right|.$$ The discrepancy of $F$, $\mathrm{disc}(F)$ is defined to be $\min_\mu \mathrm{disc}_\mu(F)$.
\end{definition}

It is known that for any function $F \colon \bits^n \times \bits^n \to \bits$, the discrepancy of $F$ is at least $2^{-O(n)}$. 
Discrepancy plays a central role in communication complexity because it characterizes
$\mathsf{PP}^{\mathsf{cc}}$, the
class of communication problems efficiently solvable by small-bias protocols~\cite{klauck2001lower}. Discrepancy is also important in circuit complexity, where it lower bounds
the size of Majority-of-Threshold circuits computing $F$ \cite{ghr, hajnal, nisan, sherstovmajmaj}. 
Finally, discrepancy
is known to be equivalent (up to a constant factor) to margin complexity~\cite{linial2009learning}, 
which is itself a fundamental notion in learning
theory. 

\medskip \noindent \textbf{Threshold weight.}
The weight of an $n$-variate polynomial $p$
is the sum of the absolute value of its coefficients. The length of $p$
is the number of non-zero coefficients of $p$.
The threshold weight (respectively, length) of $F \colon \bits^N \times \bits^N \to \bits$ is the least weight (respectively, length) of a polynomial 
$p$ with integer coefficients such that $p(x, y) \cdot F(x,y) > 0$ for all $(x, y) \in \bits^N \times \bits^N$ (note that
no restriction is placed on the degree of $p$). The threshold weight of $F$ is always at most $2^{O(n)}$,
since Parseval's inequality implies that every Boolean function is always \emph{exactly} computed
by a polynomial of weight $2^{O(n)}$. 


\subsection{Nearly optimal bounds on discrepancy, threshold weight, and more}
Via well-known techniques, we now translate our %
approximate-degree lower bounds into approximate-rank, discrepancy, and threshold-weight bounds %
in a black-box manner.

We start with the following theorem (from which \expref{Theorem}{thm:discintro} from \expref{Section}{s:discintro} follows).

\begin{theorem} \label{thm:diediedie}
	For any constant $\delta > 0$, there is an AC$^0$ function $F \colon \bits^n \times \bits^n \to \bits$ with discrepancy at most $\exp(-\Omega(n^{1-\delta}))$
	and $\eps$-approximate rank $\exp(\Omega(n^{1-\delta}))$ for some $\eps = 1-2^{-\Omega(n^{1-\delta})}$. 
\end{theorem}
\begin{proof}
	Let $f$ be the AC$^0$ function with $\eps$-approximate degree at least $n^{1-\delta}$ for some $\eps=1-2^{-\Omega(n^{1-\delta})}$ whose existence is guaranteed
	by \expref{Theorem}{ac0thm}. The pattern matrix method \cite[Theorem 8.1]{patmat} implies that for some $C=O(1)$,
	the function 
	$F \colon \bits^{Cn}\times \bits^{Cn} \to \bits$ given by
	$$F(x, y) = f\left(\dots, \vee_{j=1}^C \left(x_{i, j} \wedge y_{i, j}\right) \dots\right)$$
	satisfies
	$\text{rank}_{\eps'}(F) \geq \exp(\Omega(n^{1-\delta}))$,
	for some $\eps'=1-2^{-\Omega(n^{1-\delta})}$.  
	Moreover, by \cite[Theorem 7.3]{patmat} $F$ also satisfies $\mathrm{disc}(F) \leq \exp(-\Omega(n^{1-\delta}))$.\footnote{\label{foot}\cite[Theorem 7.3]{patmat}
		is actually expressed in terms of the degree $d$ threshold weight of $f$, rather than the $\eps$-approximate degree of $f$,
		but our lower bound on the $\eps$-approximate degree of $f$ is easily seen to imply the lower bound on the degree $d$ threshold weight of $f$
required to apply  \cite[Theorem 7.3]{patmat} (see, \eg, \cite[Lemma 20]{bt14}).} 
	Moreover, if $f$ is computed by a Boolean circuit of depth $k$ and polynomial size,
	then $F$ is computed by a Boolean circuit of polynomial size and depth $k+2$. 
	This completes the proof of the theorem for approximate rank and discrepancy.
	
	Meanwhile, a result of Krause~\cite{krause} implies that the following
	function $F'$ on $3n$ inputs has threshold weight $2^{\Omega(n^{1-\delta})}$: 
	$F'(x, y, z)= f(\dots, (x_i \wedge \overline{z}_i )\vee  (y_i \wedge z_i), \dots )$ is at least $2^{\Omega(n^{1-\delta})}$.\footnote{As
		in \expref{Footnote}{foot}, Krause's result is expressed in terms of the degree $d$ 
		threshold weight of $f$, rather than the $\eps$-approximate degree of $f$,
		but our lower bound on the $\eps$-approximate degree of $f$ is easily seen to imply the lower bound on the degree $d$ threshold weight of $f$
		required to apply Krause's result (see, \eg, \cite[Lemma 20]{bt14}).}
	Clearly, since $f$ is computed by a Boolean circuit of depth $k$ and polynomial size,
	$F'$ is computed by a Boolean circuit of polynomial size and depth $k+2$. 
\end{proof}

Via established applications of discrepancy, we conclude from \expref{Theorem}{thm:diediedie}
that for any constant $\delta > 0$, there is an AC$^0$ function $F$ with margin complexity $\exp(\Omega(n^{1-\delta}))$~\cite{linial2009learning}, 
Majority-of-Threshold circuit size $\exp(\Omega(n^{1-\delta}))$ \cite{nisan, sherstovmajmaj}, and  $\mathsf{PP}^{\mathsf{cc}}$
communication complexity $\Omega(n^{1-\delta})$~\cite{klauck2001lower}. These bounds are all nearly tight, in the sense
that any function has margin complexity $2^{O(n)}$, is computed by Majority-of-Threshold circuits of size $2^{O(n)}$, 
and has $\mathsf{PP}^{\mathsf{cc}}$ communication complexity at most $n$.

\subsection{Lower bounds for sign-rank and threshold length}
\subsubsection{Threshold length}
\begin{theorem}
There is a function $F \colon \bits^n \to \bits$ computed by a 
polynomial-size  %
circuit of depth 3 and logarithmic bottom fan-in
	such that the threshold 
	length of $F$ is $\exp(\Omega(n^{1/2}))$. 
\end{theorem}
\begin{proof}
Recall (\expref{Corollary}{surjcor}) that for $n=N \log(N^{1/2})$, 
we have  %
$\SURJ_{N^{1/2}, N} \colon \bits^n \to \bits$, and 
$\deg_{\pm}(\SURJ_{N^{1/2}, N}) = \Omega(n^{1/2})$. 
Moreover, $\SURJ_{N^{1/2}, N}$ is computed by a 
quadratic-size %
circuit of depth 3, with an $\AND$ gate at the top, 
	and logarithmic bottom fan-in.
	
	Krause and Pudl{\'{a}}k showed that if $\deg_{\pm}(f) \geq d$, then the function $F \colon \bits^{3n} \to \bits$ defined
	via $F(x, y, z) \coloneqq  f(\dots, (x_i \wedge \overline{z}_i )\vee  (y_i \wedge z_i), \dots )$
	has threshold length $2^{\Omega(d)}$. 
	If $f=\SURJ$, then $F$ is clearly computed by a 
polynomial-size %
circuit of depth 5, with an $\AND$ where 
	gates at the bottom two layers (just above the inputs) have fan-in 2, and gates at the third-to-bottom layer have fan-in 
	$O(\log n)$. Hence, each gate at the third layer from the bottom computes a function of just $O(\log n)$
	inputs, and any such function can be computed by a 
polynomial-size %
DNF and logarithmic width. Replacing each gate at the third-to-bottom
layer with an equivalent polynomial-size %
DNF of polynomial size and logarithmic width, and collapsing the two adjacent layers of $\OR$
gates, yields a depth-three %
circuit of logarithmic bottom  fan-in.
\end{proof}

\subsubsection{Sign-rank}
Recall that \expref{Corollary}{cor:signrank} from the 
%
Introduction %
asserted the existence of an AC$^0$ function
%
\[ F(x, y) \colon \bits^n \times \bits^n \to \bits \]
such that the sign-rank of the matrix $[F(x, y)]$ is
 $\exp(\tilde{\Omega}(n^{1/2}))$.
We prove \expref{Corollary}{cor:signrank} below.

Via established applications of sign-rank, we conclude as a consequence that there is a communication problem in AC$^0$ 
with $\UPP^{\text{cc}}$ communication complexity $\tilde{\Omega}(n^{1/2})$~\cite{paturisimon}, and such that all
Threshold-of-Majority circuits computing the function have size $2^{\tilde{\Omega}(n^{1/2})}$~\cite{forstertm}. 


%
\begin{proof}[Proof of Corollary \ref{cor:signrank}]
Suppose that $\deg_{\pm}(f) \geq d$. 
In order to establish sign-rank lower bounds for a certain matrix $A$ derived from $f$, Razborov and Sherstov~\cite{razborovsherstov} extended a lemma of Forster~\cite{forster} to show that it is enough to 
construct a dual witness $\mu$ for the fact that $\deg_{\pm}(f) \geq d$ that additionally satisfies a \emph{smoothness} condition. Specifically,
to show that the sign-rank of $A$ is $\exp(\Omega(d))$, it suffices to show that there is a threshold
degree dual witness $\mu$ for $f$ satisfying $\mu(x) = \exp(-O(d))$ for
all but an $\exp\left(-O(d)\right)$ fraction of inputs in $x \in \{-1, 1\}^n$. Formally, we have the following
theorem, which is implicit in~\cite{razborovsherstov} (the statement here is taken from~\cite[Theorem 4.1]{bt16}). 


\begin{restatable}[{Implicit in~\cite[Theorem 1.1]{razborovsherstov}}]{theorem}{thmrs} \label{thm:rs}
	Let $h \colon \bits^n \to \bits$ be a Boolean function, and suppose there exists a function $\tau \colon \bits^n \to \bits$ of pure high degree
	at least $d$
	such that $\tau(x) \cdot h(x) \geq 0$ for all $x \in \bits^n$, and $\|\tau\|_1=1$. 
	Moreover, suppose that $\tau(x) \ge 2^{-cd} \cdot 2^{-n}$ for all but a $2^{-cd}$ fraction of inputs $x \in \bits^n$. Then there exists a constant $C$ (depending only on $c$) such that if $F(x, y) \coloneqq  h(\dots, \land_{j = 1}^C (x_{ij} \lor y_{ij}), \dots),$
	then the matrix $[F(x, y)]_{x, y}$ has sign-rank $\exp(\Omega(d))$.
\end{restatable}

Let $R=N^{1/2}$, and assume $R+1$ is a power of 2. Letting $n=N \log(R+1)$, 
recall $\dumSURJ_{R, N}$ is a function on $n$ bits, where any $x \in \bits^n$
is interpreted as specifying a list of $N$ numbers from $[R]_0$. Here, 0 denotes
a ``dummy item'' that is ignored by $\dumSURJ$. 
We assume without loss of generality in this section that:
\begin{equation} 
\label{dummyencoding}
\text{The dummy element is represented by the string } \mathbf{1}_{\log(R+1)}\,.
\end{equation}
This ensures that strings $x \in \bits^n$ that encode mostly dummy items have low Hamming weight.

Recall that within the proof of \expref{Theorem}{thm:surj} and \expref{Corollary}{surjcor}, we proved that $\deg_{\pm}(\dumSURJ_{R, N}) \geq d$ for some $d=\Omega(N^{1/2}/\log^{3/2} N)$, via a two-step process. 
First, we borrowed a result (\expref{Theorem}{thm:main-reduction})
from our prior work~\cite{adegsurj}, which showed that for any range size $R$, $\deg_{\pm}(\dumSURJ_{R, N})$ is equivalent to $\deg_{\pm}\left(\left(\AND_R \circ \OR_N\right)^{\leq N}\right)$. Second, we constructed a dual witness $\mu \colon \mathcal{H}_{NR}^{\leq N} \to \R$ showing that the latter quantity is at least $d$. %

Unfortunately, the construction of 
$\mu$ is not sufficient to apply \expref{Theorem}{thm:rs} to $\dumSURJ_{R, N}$, for two reasons.
First, 
to apply \expref{Theorem}{thm:rs} to $\dumSURJ_{R, N}$, we need to give a smooth dual witness for $\dumSURJ_{R, N}$ itself, rather than for $\left(\AND_R \circ \OR_N\right)^{\leq N}$. Note that
$\dumSURJ_{R, N}$ is defined over the domain $\bits^n$ where $n=N \log(R+1)$, while $\left(\AND_R \circ \OR_N\right)^{\leq N}$ is defined over the domain $\mathcal{H}_{NR}^{\leq N}$. Second, the dual witness $\mu$ for  $\left(\AND_R \circ \OR_N\right)^{\leq N}$ constructed in the proof of \expref{Corollary}{surjcor}
is not smooth in the sense required by Razborov and Sherstov, as it is only ``large'' on inputs of Hamming weight at most $d$ (see the first paragraph of \expref{Section}{igivethefrigup}
for further discussion of this point). 


We will address both of the above issues as follows. First, we will show how to turn $\mu$ into a dual witness $\hat{\sigma}$
for the fact that $\deg_{\pm}\left(\dumSURJ_{R, N}\right)\geq d$, such that $\hat{\sigma}$ inherits the ``largeness'' property of $\mu$ on inputs of Hamming weight at most $d$. 
Second, we transform $\hat{\sigma}$ into a dual witness $\tau$ for the fact that $\deg_{\pm}\left(\dumSURJ_{R, N} \circ \AND_{\log^2 n} \circ \PARITY_{\log^3 n} \right)\geq d$, such
that $\tau$ satisfies the smoothness condition required to apply \expref{Theorem}{thm:rs}.

\medskip
\noindent \textbf{Construction and analysis of $\hat{\sigma}$.}
In the proof of \expref{Theorem}{surjthm}, we constructed a dual witness $\mu$ for $F=\AND_R \circ \OR_N$ satisfying the following 
conditions.

\begin{equation} \label{asdf1}
\mu(x) \cdot F(x) \geq 0 \text{ for all } x \in \bits^{R \cdot N}\,.
\end{equation}
\begin{equation} \label{asdf2}
\mu(x) = 0  \text{ for all } x \in \mathcal{H}^{> N}_{R \cdot N}\,.
\end{equation}
\begin{equation} \label{asdf3}
\mu \text{ has pure high degree at least } d\,.
\end{equation}
\begin{equation} \label{nomorelabelsbrah}
\mu(x) \geq 2^{-1-R/5} \text{ for all } x \in \mathcal{H}_{NR}^{\leq d}\,.
\end{equation}
\begin{equation} \label{asdf4}
\mu(x) \text{ has } \ell_1\text{-norm at most } 2\,.
\end{equation}

\expeqref{Expressions}{asdf1}-\eqref{asdf3} are explicitly established in the proof of \expref{Theorem}{surjthm} (specifically, the part entitled ``Analysis of $\mu$'').  %
\expeqref{Expression}{nomorelabelsbrah} is immediate from the penultimate inequality in 
\expeqref{Expression}{forsignrank}.
\expeqref{Equation}{asdf4} holds because we defined $\mu= \zeta - \sum_{x \in B} \psi_{\mathsf{corr}, x}$ (\expeqref{Expression}{omg1}),
and hence $$\|\mu\|_1 \leq \|\zeta\|_1 + \sum_{x \in B} \|\psi_{\mathsf{corr}, x}\|_1 \leq
1 + \left(E\left(\zeta, \AND_R \circ \OR_N\right) +W\left(\zeta, f\right)\right) \cdot 2^d \cdot {RN \choose d}
\leq 2.$$ 


Defining $\mu'=\mu/\|\mu\|_1$ yields a function that also satisfies
\expeqref{Equations}{asdf1}-\eqref{asdf3}, and such that:

\begin{equation} \label{nomorelabelsbrahsweet}
\mu'(x) \geq 2^{-R/4} \text{ for all } x \in \mathcal{H}_{NR}^{\leq d}\,.
\end{equation}
\begin{equation} \label{asdf42}
\mu'(x) \text{ has } \ell_1\text{-norm } 1\,.
\end{equation}



Let $\mathcal{Z}$ denote the subset of $\left([N]_0\right)^R$ defined as follows: $\mathcal{Z} \coloneqq  \{z \in \left([N]_0\right)^R \colon z_1 + \dots +z_R \leq N\}$. 
Define $G(z_1, \dots, z_R) \colon \mathcal{Z} \to \bits$ to equal $-1$
if and only if $z_i \geq 1$ for all $i=1, \dots, R.$ 


We will now turn $\mu'$ (which is defined over domain $\bits^{NR}$) into a function over domain $\mathcal{Z}$. 


Define $$\sigma(z_1, \dots, z_R) = \sum_{x=(x_1, \dots, x_R) \in \left(\bits^{N}\right)^R: |x_i|=z_i \text{ for all } i} \mu'(x).$$
We claim that $\sigma$ has
the following properties.
\begin{equation} \label{mysouliscrushed1}
\sigma(z) \cdot G(z) \geq 0 \text{ for all } z \in \mathcal{Z}.
\end{equation}
\begin{equation}
\sum_{z \in \mathcal{Z}} |\sigma(z)|=1\,. \label{brr2}
\end{equation}
\begin{equation} \label{phdsigmasigh}
\deg(q) \leq d \Longrightarrow \sum_{z \in \mathcal{Z}} \sigma(z) \cdot q(z) =0\,.
\end{equation}
\begin{equation} \label{zrp2}
\sigma(z) \geq 2^{-R/4} \text{ for all } z \in \mathcal{Z} \text{ such that } \sum_{i=1}^{R} z_i \leq d\,.
\end{equation}

\expeqref{Expression}{mysouliscrushed1} is immediate from \expeqref{Expression}{asdf1}. 
\expeqref{Equation}{brr2} is immediate from \expeqref{Equations}{asdf42} \expeqref{and}{asdf1}, and the fact that $F(x)=F(x')$ if $|x_i|=|x'_i|$ for $i=1, \dots, R$. 
\expeqref{Equation}{phdsigmasigh} holds because
$$\sum_{z \in \mathcal{Z}} \sigma(z) q(z) = \sum_{x=(x_1, \dots, x_R) \in \bits^{R \cdot N}} \mu'(x) \cdot q(|x_1|, \dots, |x_R|)
= 0,$$ 
where the last equality holds because $q(|x_1|, \dots, |x_R|)$ is a polynomial over $\left(\bits^{N}\right)^R$ of degree at most $\deg(q)$. 
\expeqref{Expression}{zrp2} is immediate from \expeqref{Expression}{nomorelabelsbrahsweet}. 

Finally, we are in a position to define our desired
dual witness $\hat{\sigma} \colon \bits^n \to \R$,
which will witness the fact that $\deg_{\pm}(\dumSURJ_{R, N}) \geq d/\log R$. 
For an $x \in \bits^n$ interpreted as a sequence $(s_1, \dots, s_N)$ in $[R]_0^N$, let $z_i(x)$ denote $\left| \{j \colon s_j=i\}\right|$,
and let $z(x)=(z_1(x), \dots, z_R(x)) \in \mathcal{Z}$. For a $z^* \in \mathcal{Z}$, let $N(z^*) = |\{x \in \bits^n \colon z(x)=z^*\}|$. 
Define $\hat{\sigma} \colon \bits^n \to \mathbb{R}$ via: 
%
\begin{equation} \label{hatsigmadef} \hat{\sigma}(x) = \frac{1}{N(z(x))} \cdot \sigma(z(x))\,.\end{equation}

%
We observe that $\hat{\sigma}$ has the following properties. %

\begin{equation} \label{mysouliscrushedhat}
\hat{\sigma}(x) \cdot \dumSURJ(x) \geq 0 \text{ for all } x \in \bits^n\,.
\end{equation}
\begin{equation} \label{brrr}
\hat{\sigma}(x) \text{ has } \ell_1\text{-norm } 1\,.
\end{equation}
\begin{equation} \label{zrp}
\hat{\sigma}(x) \geq 2^{-R/3} \text{ for all } x \in \mathcal{H}^{\leq d}_{n}\,.
\end{equation}
\begin{equation} \label{phdsigmasighhat}
\hat{\sigma}(x) \text{ has pure high degree at least } d/\log R\,.
\end{equation}

\expeqref{Expression}{mysouliscrushedhat} follows from \expeqref{Expression}{mysouliscrushed1} and the definition of $\hat{\sigma}$.
\expeqref{Equation}{brrr} is immediate from the definition of $\hat{\sigma}$ and \expeqref{Equation}{brr2}. 
To see that \expeqref{Expression}{zrp} holds, observe that if $x \in \mathcal{H}^{\leq d}_n$, then
\expeqref{Equation}{dummyencoding} implies that $\sum_{i=1}^{R} z_i(x) \leq d$. Hence,
\expeqref{Expression}{zrp2} implies that $\sigma(z(x)) \geq 2^{-R/4}$. Moreover, for any $x \in \mathcal{H}^{\leq d}_{n}$,
$N(z(x)) \leq {n \choose d}$ (this is because for any $x, x'$, if $z(x) = z(x')$, then $|x|=|x'|$, 
so $N(z(x)) \leq {n \choose |x|} \leq {n \choose d}$).
Combining these two inequalities with the definition of $\hat{\sigma}$ (\expeqref{Equation}{hatsigmadef}) yields
that $\hat{\sigma}(x) \geq 2^{-R/4} \cdot {n \choose d}^{-1} \geq 2^{-R/3}$. 


\expeqref{Equation}{phdsigmasighhat} follows from \expref{Proposition}{keysrprop} below, combined with \expeqref{Equation}{phdsigmasigh}.

\begin{proposition}\label{keysrprop} Let $\sigma \colon\colon \mathcal{Z} \to \R$
	satisfy \expeqref{Equation}{phdsigmasigh}, and let $\hat{\sigma}$ be as per \expeqref{Equation}{hatsigmadef}. 
	Then  
	$\hat{\sigma}$ has pure high degree at least $d/\log R$. 
\end{proposition}
\expref{Proposition}{keysrprop} is essentially a dual formulation of an important lemma of Ambainis~\cite{ambainissmallrange}, as we now explain. 

%
\begin{lemma}[Ambainis~\cite{ambainissmallrange}] \label{ambainislemma}
	Let $n=N \log(R+1)$. Let $p \colon \bits^n \to \bits$ be any polynomial of degree at most $d$. 
	Then there is a polynomial $q \colon [R]_0^N \to \R$ of degree at most $d \log R$
	such that for all $z^* \in \mathcal{Z}$, $q(z^*) = \mathbb{E}_{x \colon z(x)=z^*}[p(x)]$.
\end{lemma}
\expref{Proposition}{keysrprop} follows from \expref{Lemma}{ambainislemma} by the following reasoning.
If $p \colon \bits^n \to \bits$ is polynomial of degree at most $d / \log R$,
then \begin{equation} \sum_{x \in \bits^n} \hat{\sigma}(x) \cdot p(x) = \sum_{z \in \mathcal{Z}} \sigma(z) \cdot q(z) =0\,, \label{sighsighsighsigh}\end{equation}
where the polynomial $q$ is as in \expref{Lemma}{ambainislemma}, and the second equality holds by \expeqref{Equation}{phdsigmasigh}.

\medskip \noindent \textbf{Construction and analysis of $\tau$.}
Now that we have constructed a dual witness $\hat{\sigma}$ for the high threshold degree of $\dumSURJ$ 
(captured by \expeqref{Equations}{mysouliscrushedhat}, \eqref{brrr}, \expeqref{and}{phdsigmasighhat}), that additionally %
%
satisfies  %
the extra condition
of ``largeness on low-Hamming-weight inputs'' (\expeqref{Equation}{zrp}),
we can turn to constructing a dual witness for $\dumSURJ \circ \AND_{\log^2 n} \circ \PARITY_{\log^3 n}$
that satisfies the smoothness condition needed by Razborov's and Sherstov's sign-rank analysis.

Specifically, let $a=\log^2 n$, and define $\psi \colon \bits^{a} \to \R$ via:
$$\psi(x) = \begin{cases} -1/2 \text{ if } x=\mathbf{-1}_{a}\\
1/(2 \cdot (2^a-1)) \text{ otherwise}. \end{cases}$$

Clearly, $\psi$ has pure high degree at least $1$, has $\ell_1$-norm $1$, and $\psi(x) \cdot \AND_a(x) \geq 0$
for all $x \in \bits^a$.
Consider the dual witness $\zeta = \hat{\sigma} \ls \psi$. 
Then $\zeta$ has the following properties:
\begin{equation} \label{umissed1bro}
\zeta \text{ has } \ell_1\text{-norm } 1\,.
\end{equation}
\begin{equation} \label{thismuststopbrah}
\zeta \text{ has pure high degree at least } d/\log R\,.
\end{equation}
\begin{equation} \label{asdfasdf}
\zeta(x) \cdot \left(\dumSURJ \circ \AND_{\log^2 n}\right)(x) \geq 0 \text{ for all } x \in \bits^{n \cdot a}\,. %
\end{equation}
Here, \expeqref{Equation}{umissed1bro} follows from \expeqref{Equation}{brrr}, the fact that $\psi$ has $\ell_1$-norm $1$, and \expeqref{Property}{eqn:ls-norm}. \expeqref{Equation}{thismuststopbrah} follows from \expeqref{Equation}{phdsigmasighhat}, \expeqref{Equation}{eqn:ls-phd},
and the fact that $\psi$ has pure high degree at least $1$.
\expeqref{Expression}{asdfasdf} follows from the definition of dual block composition, \expeqref{Expression}{mysouliscrushedhat},
and the fact that $\psi(x) \cdot \AND_a(x) \geq 0$
for all $x \in \bits^a$.

Let $S$ be the set of all inputs $x=(x_1, \dots, x_n) \in \left(\bits^a\right)^{n}$ such that 
fewer than $d$ 
of the $x_i$'s are equal to $-\mathbf{1}_a$.
By \expeqref{Expression}{zrp} and the definition of dual block composition,
\begin{equation} \label{eq:tautautau}
\text{For any input } x\in S\text{, } |\zeta(x)| \geq 2^{-R/3} \cdot (2^{a}-1)^{-n} \geq 2^{-R/3}  \cdot 2^{-an}\,.\end{equation}
Moreover, 
\begin{equation}
\label{eq:tautautau2} \Pr_{x \sim \left(\bits^a\right)^n}[x \not \in S] 
\leq 2^{-a d} \cdot {n \choose d} \leq n^{-\Omega(d \log n)} \leq 2^{-\Omega(R)}\,. \end{equation}

Finally, letting $\ell=\log^3 n$, let $\eta \colon \bits^{\ell} \to \R $ be defined by $\eta(x) = 2^{-\ell} \cdot \PARITY_{\ell}(x)$. 
Observe that $\eta$ has $\ell_1$-norm 1, has pure high degree $\ell$, and $\sgn(\eta(x))=\sgn(\PARITY_{\ell}(x))$ for all $x \in \bits^\ell$.
Then \expeqref{Equation}{eqn:ls-norm} implies that $\tau \coloneqq  \zeta \ls \eta$ has $\ell_1$-norm $1$, \expeqref{Equation}{eqn:ls-phd} implies that $\tau$ has pure high degree at least $(d/\log R) \cdot \ell \geq R$,
and it follows from \expeqref{Equations}{eq:tautautau} \expeqref{and}{eq:tautautau2} that 
\begin{equation}
\Pr_{x \sim \left(\left(\bits^{\ell}\right)^a\right)^n}[|\tau(x)| < 2^{-R/3} \cdot 2^{-\ell \cdot a \cdot n}] \leq 2^{-\Omega(R)}\,.
\end{equation}

Hence, we can apply \expref{Theorem}{thm:rs} to the function $h=\dumSURJ \circ \AND_{\log^2 n} \circ \PARITY_{\log^3 n}$, which is defined on $\tilde{O}(n)$
variables, to obtain an AC$^0$ function $F(x, y)$ that is also defined on $\tilde{O}(n)$ variables, such that $[F(x, y)]_{x, y}$ has sign-rank $\exp(\tilde{\Omega}(n^{1/2}))$\,.
\end{proof}





\subsection{Secret sharing schemes}
Bogdanov et al.~\cite{viola} observed that for any
$f \colon \bits^n \to \bits$ and integer $d>0$, any dual polynomial $\mu$
for the fact that $\adeg_{\eps}(f) \geq d$ leads to a scheme
for sharing a single secret bit $b \in \bits$ among $n$ parties as follows.
Decompose $\mu$ as $\mu_{+} - \mu_{-}$, where $\mu_{+}$ and $\mu_{-}$
are non-negative functions with $\|\mu_{+}\|_1 = \|\mu_{0}\|_1 = 1/2$. Then in order to split $b$
among $n$ parties,
one draws an input $x=(x_1, \dots, x_n) \in \bits^n$ from the distribution $2 \cdot \mu_b$, and
gives bit $x_i$ to the $i$th party.
In order to reconstruct $b$, one applies $f$ to $(x_1, \dots, x_n)$.

Because $\mu$ is $\eps$-correlated with $f$, the probability of correct reconstruction if the bit is 
chosen at random
is at least $(1+\eps)/2$ (and the the \emph{reconstruction advantage}, 
defined to equal $\Pr_{x \sim \mu_{+}}[f(x) = 1] - \Pr_{x\sim\mu_{-}}[f(x) = 1]$, is at least
$\eps$). 
The fact that $\mu$ has pure high degree at least $d$
means that any subset of shares of size less than $d$ provides no information 
about the secret bit $b$. 
Hence, an immediate corollary of \expref{Theorem}{ac0thm} is the following.

\begin{corollary} \label{cor:secret-sharing}
	For any arbitrarily small constant $\delta > 0$, there is a secret sharing scheme 
	that shares a single bit $b$ among $n$ parties by assigning a bit $x_i$ to each party $i$. The scheme has the following properties.
\begin{enumerate}
\item The reconstruction procedure is computed by an $\operatorname{AC}^0$ circuit.
\item The reconstruction advantage is at least $1-2^{-\Omega(n^{1-\delta})}$.
\item Any subset of shares of size less than $d = \Omega(n^{1-\delta})$ provides no information about the secret bit $b$.
\end{enumerate}
\end{corollary}

The best previous results could only guarantee security against subsets of shares of size $d=\Theta(n^{1/2})$,
or could only guarantee reconstruction advantage bounded away from 1 \cite{viola, adegsurj}. Cheng et al.~\cite{cheng2017near} recently considered
a relaxed notion of security, where even very small subsets of shares are allowed to 
provide (a bounded amount of) information about the secret bit $b$. Under this relaxed notion of security, they achieved perfect reconstruction
and security against subsets of size $\Omega(n)$. 

\medskip \noindent \textbf{Acknowledgments.} The authors are grateful to Robin Kothari, Nikhil Mande, Jonathan Ullman, and the anonymous reviewers
for valuable comments earlier versions of this manuscript. 




%% Theory of Computing
%% File: v017a007.bbl
%% Publication date: September 23, 2021
%% Authors: Mark Bun and Justin Thaler
\newcommand{\prelim}{Preliminary version in\ }
% Define empty bibhead if not already defined
\providecommand{\bibhead}[1]{}
% Define tocrefpdfbookmark if not already defined
\expandafter\ifx\csname pdfbookmark\endcsname\relax%
  \providecommand{\tocrefpdfbookmark}{}
\else\providecommand{\tocrefpdfbookmark}{%
   \hypertarget{tocreferences}{}%
   \pdfbookmark[1]{References}{tocreferences}}%
\fi

\tocrefpdfbookmark
\begin{thebibliography}{10}

\bibitem{aaronsonshi}\bibhead{aaronsonshi}
{\sc Scott Aaronson and Yaoyun Shi}: Quantum lower bounds for the collision and
  the element distinctness problems.
\newblock {\em J. ACM}, 51(4):595--605, 2004.
\newblock [\epfmtdoi{10.1145/1008731.1008735}]

\bibitem{AB84}\bibhead{AB84}
{\sc Miklos Ajtai and Michael Ben-Or}: A theorem on probabilistic constant
  depth computations.
\newblock In {\em Proc. 16th STOC}, pp. 471--474. ACM Press, 1984.
\newblock [\epfmtdoi{10.1145/800057.808715}]

\bibitem{allender}\bibhead{allender}
{\sc Eric Allender}: A note on the power of threshold circuits.
\newblock In {\em Proc. 30th FOCS}, pp. 580--584. IEEE Comp. Soc., 1989.
\newblock [\epfmtdoi{10.1109/SFCS.1989.63538}]

\bibitem{ambainissmallrange}\bibhead{ambainissmallrange}
{\sc Andris Ambainis}: Polynomial degree and lower bounds in quantum
  complexity: Collision and element distinctness with small range.
\newblock {\em Theory of Computing}, 1(3):37--46, 2005.
\newblock [\epfmtdoi{10.4086/toc.2005.v001a003}]

\bibitem{babaietal}\bibhead{babaietal}
{\sc L{\'{a}}szl{\'{o}} Babai, P{\'e}ter Frankl, and Janos Simon}: Complexity
  classes in communication complexity theory (preliminary version).
\newblock In {\em Proc. 27th FOCS}, pp. 337--347. IEEE Comp. Soc., 1986.
\newblock [\epfmtdoi{10.1109/SFCS.1986.15}]

\bibitem{beame}\bibhead{beame}
{\sc Paul Beame and Widad Machmouchi}: The quantum query complexity of
  {AC}$^0$.
\newblock {\em Quantum Inf. Comput.}, 12(7-8):670--676, 2012.
\newblock [\epfmtdoi{10.26421/QIC12.7-8-11}]

\bibitem{viola}\bibhead{viola}
{\sc Andrej Bogdanov, Yuval Ishai, Emanuele Viola, and Christopher Williamson}:
  Bounded indistinguishability and the complexity of recovering secrets.
\newblock In {\em Proc. 36th Ann. Internat. Cryptology Conf. (CRYPTO'16)}, pp.
  593--618. Springer, 2016.
\newblock [\epfmtdoi{10.1007/978-3-662-53015-3\_21}, \epfmt{eccc}{TR15-182}]

\bibitem{bogicalp}\bibhead{bogicalp}
{\sc Andrej Bogdanov and Christopher Williamson}: Approximate bounded
  indistinguishability.
\newblock In {\em Proc. 44th Internat. Colloq. on Automata, Languages, and
  Programming (ICALP'17)}, pp. 53:1--11. Schloss Dagstuhl--Leibniz-Zentrum fuer
  Informatik, 2017.
\newblock [\epfmtdoi{10.4230/LIPIcs.ICALP.2017.53}]

\bibitem{bchtv}\bibhead{bchtv}
{\sc Adam Bouland, Lijie Chen, Dhiraj Holden, Justin Thaler, and
  Prashant~Nalini Vasudevan}: On the power of statistical zero knowledge.
\newblock {\em SIAM J. Comput.}, 49(4):FOCS17:1--58, 2019.
\newblock \prelim \href{https://doi.org/10.1109/FOCS.2017.71}{FOCS'17}.
\newblock [\epfmtdoi{10.1137/17M1161749}, \epfmt{eccc}{TR16-140},
  \epfmt{arxiv}{1609.02888}]

\bibitem{buhrman1999bounds}\bibhead{buhrman1999bounds}
{\sc Harry Buhrman, Richard Cleve, Ronald~de Wolf, and Christof Zalka}: Bounds
  for small-error and zero-error quantum algorithms.
\newblock In {\em Proc. 40th FOCS}, pp. 358--368. IEEE Comp. Soc., 1999.
\newblock [\epfmtdoi{10.1109/SFFCS.1999.814607}, \epfmt{arxiv}{cs/9904019}]

\bibitem{bvdw}\bibhead{bvdw}
{\sc Harry Buhrman, Nikolai~K. Vereshchagin, and Ronald~de Wolf}: On
  computation and communication with small bias.
\newblock In {\em Proc. 22nd IEEE Conf. on Comput. Complexity (CCC'07)}, pp.
  24--32. IEEE Comp. Soc., 2007.
\newblock [\epfmtdoi{10.1109/CCC.2007.18}]

\bibitem{bkt}\bibhead{bkt}
{\sc Mark Bun, Robin Kothari, and Justin Thaler}: The polynomial method strikes
  back: {T}ight quantum query bounds via dual polynomials.
\newblock {\em Theory of Computing}, 16(10):1--71, 2020.
\newblock \prelim \href{https://doi.org/10.1145/3188745.3188784}{STOC'18}.
\newblock [\epfmtdoi{10.4086/toc.2020.v016a010}]

\bibitem{bt13}\bibhead{bt13}
{\sc Mark Bun and Justin Thaler}: Dual lower bounds for approximate degree and
  {M}arkov--{B}ernstein inequalities.
\newblock {\em Inform. Comput.}, 243:2--25, 2015.
\newblock \prelim
  \href{https://doi.org/10.1007/978-3-642-39206-1\_26}{ICALP'13}.
\newblock [\epfmtdoi{10.1016/j.ic.2014.12.003}]

\bibitem{bt14}\bibhead{bt14}
{\sc Mark Bun and Justin Thaler}: Hardness amplification and the approximate
  degree of constant-depth circuits.
\newblock In {\em Proc. 42nd Internat. Colloq. on Automata, Languages, and
  Programming (ICALP'15)}, pp. 268--280. Springer, 2015.
\newblock [\epfmtdoi{10.1007/978-3-662-47672-7\_22}, \epfmt{arxiv}{1311.1616}]

\bibitem{bt16}\bibhead{bt16}
{\sc Mark Bun and Justin Thaler}: Improved bounds on the sign-rank of {AC}$^0$.
\newblock In {\em Proc. 43rd Internat. Colloq. on Automata, Languages, and
  Programming (ICALP'16)}, pp. 37:1--14. Schloss Dagstuhl--Leibniz-Zentrum fuer
  Informatik, 2016.
\newblock [\epfmtdoi{10.4230/LIPIcs.ICALP.2016.37}]

\bibitem{dlweight}\bibhead{dlweight}
{\sc Mark Bun and Justin Thaler}: Approximate degree and the complexity of
  depth three circuits.
\newblock In {\em Proc. 22nd Internat. Workshop on Randomization and
  Computation (RANDOM'18)}, pp. 35:1--18. Schloss Dagstuhl--Leibniz-Zentrum
  fuer Informatik, 2018.
\newblock [\epfmtdoi{10.4230/LIPIcs.APPROX-RANDOM.2018.35},
  \epfmt{eccc}{TR16-121}]

\bibitem{conf-version}\bibhead{conf-version}
{\sc Mark Bun and Justin Thaler}: The large-error approximate degree of
  {AC}$^0$.
\newblock In {\em Proc. 23rd Internat. Workshop on Randomization and
  Computation (RANDOM'19)}, pp. 55:1--16. Schloss Dagstuhl--Leibniz-Zentrum
  fuer Informatik, 2019.
\newblock [\epfmtdoi{10.4230/LIPIcs.APPROX-RANDOM.2019.55}]

\bibitem{adegsurj}\bibhead{adegsurj}
{\sc Mark Bun and Justin Thaler}: A nearly optimal lower bound on the
  approximate degree of {AC}$^0$.
\newblock {\em SIAM J. Comput.}, 49(4):FOCS17:59--96, 2019.
\newblock \prelim \href{https://doi.org/10.1109/FOCS.2017.10}{FOCS'17}.
\newblock [\epfmtdoi{10.1137/17M1161737}]

\bibitem{cheng2017near}\bibhead{cheng2017near}
{\sc Kuan Cheng, Yuval Ishai, and Xin Li}: Near-optimal secret sharing and
  error correcting codes in $\mathsf{AC}^0$.
\newblock In {\em Proc. Theory of Cryptography Conf. (TCC'17)}, pp. 424--458.
  Springer, 2017.
\newblock [\epfmtdoi{10.1007/978-3-319-70503-3\_14}]

\bibitem{feldmanevolve}\bibhead{feldmanevolve}
{\sc Vitaly Feldman}: Evolvability from learning algorithms.
\newblock In {\em Proc. 40th STOC}, pp. 619--628. ACM Press, 2008.
\newblock [\epfmtdoi{10.1145/1374376.1374465}]

\bibitem{forster}\bibhead{forster}
{\sc J{\"{u}}rgen Forster}: A linear lower bound on the unbounded error
  probabilistic communication complexity.
\newblock {\em J. Comput. System Sci.}, 65(4):612--625, 2002.
\newblock \prelim \href{https://doi.org/10.1109/CCC.2001.933877}{CCC'01}.
\newblock [\epfmtdoi{10.1016/S0022-0000(02)00019-3}]

\bibitem{forstertm}\bibhead{forstertm}
{\sc J{\"{u}}rgen Forster, Matthias Krause, Satyanarayana~V. Lokam, Rustam
  Mubarakzjanov, Niels Schmitt, and Hans~Ulrich Simon}: Relations between
  communication complexity, linear arrangements, and computational complexity.
\newblock In {\em Proc. 21st Found. Softw. Techn. Theoret. Comp. Sci. Conf.
  (EFSTTCS'01)}, pp. 171--182. Springer, 2001.
\newblock [\epfmtdoi{10.1007/3-540-45294-X\_15}]

\bibitem{ghr}\bibhead{ghr}
{\sc Mikael Goldmann, Johan H{\aa}stad, and Alexander~A. Razborov}: Majority
  gates {vs.} general weighted threshold gates.
\newblock {\em Comput. Complexity}, 2(4):277--300, 1992.
\newblock [\epfmtdoi{10.1007/BF01200426}]

\bibitem{hajnal}\bibhead{hajnal}
{\sc Andr{\'{a}}s Hajnal, Wolfgang Maass, Pavel Pudl{\'{a}}k, Mario Szegedy,
  and Gy{\"{o}}rgy Tur{\'{a}}n}: Threshold circuits of bounded depth.
\newblock {\em J. Comput. System Sci.}, 46(2):129--154, 1993.
\newblock [\epfmtdoi{10.1016/0022-0000(93)90001-D}]

\bibitem{kahn}\bibhead{kahn}
{\sc Jeff Kahn, Nathan Linial, and Alex Samorodnitsky}: Inclusion-exclusion:
  {E}xact and approximate.
\newblock {\em Combinatorica}, 16(4):465--477, 1996.
\newblock [\epfmtdoi{10.1007/BF01271266}]

\bibitem{klauck2001lower}\bibhead{klauck2001lower}
{\sc Hartmut Klauck}: Lower bounds for quantum communication complexity.
\newblock {\em SIAM J. Comput.}, 37(1):20--46, 2007.
\newblock \prelim \href{https://doi.org/10.1109/SFCS.2001.959903}{FOCS'01}.
\newblock [\epfmtdoi{10.1137/S0097539702405620},
  \epfmt{arxiv}{quant-ph/0106160}]

\bibitem{ksdnf}\bibhead{ksdnf}
{\sc Adam~R. Klivans and Rocco~A. Servedio}: Learning {DNF} in time
  {$2^{\tilde{O}(n^{1/3})}$}.
\newblock {\em J. Comput. System Sci.}, 68(2):303--318, 2004.
\newblock [\epfmtdoi{10.1016/j.jcss.2003.07.007}]

\bibitem{Kop13}\bibhead{Kop13}
{\sc Swastik Kopparty}: {AC}\textsuperscript{0} lower bounds and
  pseudorandomness, 2013.
\newblock
  \href{https://sites.math.rutgers.edu/~sk1233/courses/topics-S13/lec4.pdf}{Lecture
  notes of ``Topics in Complexity Theory and Pseudorandomness (Spring 2013)''
  at Rutgers University}.

\bibitem{krause}\bibhead{krause}
{\sc Matthias Krause}: On the computational power of {B}oolean decision lists.
\newblock {\em Comput. Complexity}, 14(4):362--375, 2006.
\newblock [\epfmtdoi{10.1007/s00037-005-0203-0}]

\bibitem{kp}\bibhead{kp}
{\sc Matthias Krause and Pavel Pudl{\'{a}}k}: On the computational power of
  depth-2 circuits with threshold and modulo gates.
\newblock {\em Theoret. Comput. Sci.}, 174(1--2):137--156, 1997.
\newblock [\epfmtdoi{10.1016/S0304-3975(96)00019-9}]

\bibitem{lee}\bibhead{lee}
{\sc Troy Lee}: A note on the sign degree of formulas, 2009.
\newblock [\epfmt{arxiv}{0909.4607}]

\bibitem{linial2009learning}\bibhead{linial2009learning}
{\sc Nati Linial and Adi Shraibman}: Learning complexity vs communication
  complexity.
\newblock {\em Combin. Probab. Comput.}, 18(1--2):227--245, 2009.
\newblock [\epfmtdoi{10.1017/S0963548308009656}]

\bibitem{mp}\bibhead{mp}
{\sc Marvin Minsky and Seymour Papert}: {\em Perceptrons: {A}n Introduction to
  Computational Geometry}.
\newblock MIT Press, 1969.
\newblock [\epfmtdoi{10.7551/mitpress/11301.001.0001}]

\bibitem{nisan}\bibhead{nisan}
{\sc Noam Nisan}: The communication complexity of threshold gates.
\newblock In {\em Combinatorics, Paul Erd\H{o}s is Eighty}, pp. 301--315.
  J{\'a}nos Bolyai Math. Soc., Budapest, 1994.

\bibitem{osnewbounds}\bibhead{osnewbounds}
{\sc Ryan O'Donnell and Rocco~A. Servedio}: New degree bounds for polynomial
  threshold functions.
\newblock {\em Combinatorica}, 30(3):327--358, 2010.
\newblock \prelim \href{https://doi.org/10.1145/780542.780592}{STOC'03}.
\newblock [\epfmtdoi{10.1007/s00493-010-2173-3}]

\bibitem{paturisimon}\bibhead{paturisimon}
{\sc Ramamohan Paturi and Janos Simon}: Probabilistic communication complexity.
\newblock {\em J. Comput. System Sci.}, 33(1):106--123, 1986.
\newblock [\epfmtdoi{10.1016/0022-0000(86)90046-2}]

\bibitem{podolskii2007perceptrons}\bibhead{podolskii2007perceptrons}
{\sc Vladimir~V. Podolskii}: Perceptrons of large weight.
\newblock {\em Probl. Info. Transmission}, 45:46--53, 2009.
\newblock \prelim \href{https://doi.org/10.1007/978-3-540-74510-5\_33}{CSR'07}.
\newblock [\epfmtdoi{10.1134/S0032946009010062}]

\bibitem{razborovsherstov}\bibhead{razborovsherstov}
{\sc Alexander~A. Razborov and Alexander~A. Sherstov}: The sign-rank of
  {AC}$^0$.
\newblock {\em SIAM J. Comput.}, 39(5):1833--1855, 2010.
\newblock [\epfmtdoi{10.1137/080744037}]

\bibitem{sherstovmajmaj}\bibhead{sherstovmajmaj}
{\sc Alexander~A. Sherstov}: Separating {AC}$^0$ from depth-2 majority
  circuits.
\newblock {\em SIAM J. Comput.}, 38(6):2113--2129, 2009.
\newblock [\epfmtdoi{10.1137/08071421X}]

\bibitem{patmat}\bibhead{patmat}
{\sc Alexander~A. Sherstov}: The pattern matrix method.
\newblock {\em SIAM J. Comput.}, 40(6):1969--2000, 2011.
\newblock \prelim \href{https://doi.org/10.1145/1374376.1374392}{STOC'08}.
\newblock [\epfmtdoi{10.1137/080733644}]

\bibitem{sherstovdp}\bibhead{sherstovdp}
{\sc Alexander~A. Sherstov}: Strong direct product theorems for quantum
  communication and query complexity.
\newblock {\em SIAM J. Comput.}, 41(5):1122--1165, 2012.
\newblock \prelim \href{https://doi.org/10.1145/1993636.1993643}{STOC'14}.
\newblock [\epfmtdoi{10.1137/110842661}]

\bibitem{sherstovhalfspaces1}\bibhead{sherstovhalfspaces1}
{\sc Alexander~A. Sherstov}: The intersection of two halfspaces has high
  threshold degree.
\newblock {\em SIAM J. Comput.}, 42(6):2329--2374, 2013.
\newblock \prelim \href{https://doi.org/10.1109/FOCS.2009.18}{FOCS'09}.
\newblock [\epfmtdoi{10.1137/100785260}]

\bibitem{sherstov14}\bibhead{sherstov14}
{\sc Alexander~A. Sherstov}: Breaking the {M}insky-{P}apert barrier for
  constant-depth circuits.
\newblock {\em SIAM J. Comput.}, 47(5):1809--1857, 2018.
\newblock \prelim \href{https://doi.org/10.1145/2591796.2591871}{STOC'14}.
\newblock [\epfmtdoi{10.1137/15M1015704}]

\bibitem{sherstov15}\bibhead{sherstov15}
{\sc Alexander~A. Sherstov}: The power of asymmetry in constant-depth circuits.
\newblock {\em SIAM J. Comput.}, 47(6):2362--2434, 2018.
\newblock \prelim \href{https://doi.org/10.1109/FOCS.2015.34}{FOCS'15}.
\newblock [\epfmtdoi{10.1137/100785260}]

\bibitem{algopoly}\bibhead{algopoly}
{\sc Alexander~A. Sherstov}: Algorithmic polynomials.
\newblock {\em SIAM J. Comput.}, 49(6):1173--1231, 2020.
\newblock \prelim \href{https://doi.org/10.1145/3188745.3188958}{STOC'18}.
\newblock [\epfmtdoi{10.1137/19M1278831}, \epfmt{arxiv}{1801.04607}]

\bibitem{SherstovWu19}\bibhead{SherstovWu19}
{\sc Alexander~A. Sherstov and Pei Wu}: Near-optimal lower bounds on the
  threshold degree and sign-rank of {AC}$^0$.
\newblock {\em SIAM J. Comput.}, online:STOC19:1--86, 2021.
\newblock \prelim \href{https://doi.org/10.1145/3313276.3316408}{STOC'19}.
\newblock [\epfmtdoi{10.1137/20M1312459}, \epfmt{arxiv}{1901.00988}]

\bibitem{shizhu}\bibhead{shizhu}
{\sc Yaoyun Shi and Yufan Zhu}: Quantum communication complexity of
  block-composed functions.
\newblock {\em Quantum Inf. Comput.}, 9(5):444--460, 2009.
\newblock [\epfmtdoi{10.26421/QIC9.5-6-7}]

\bibitem{evolve}\bibhead{evolve}
{\sc Leslie~G. Valiant}: Evolvability.
\newblock {\em J. ACM}, 56(1):3:1--21, 2009.
\newblock [\epfmtdoi{10.1145/1462153.1462156}]

\bibitem{viola2009approximate}\bibhead{viola2009approximate}
{\sc Emanuele Viola}: On approximate majority and probabilistic time.
\newblock {\em Comput. Complexity}, 18(3):337--375, 2009.
\newblock \prelim \href{https://doi.org/10.1109/CCC.2007.16}{CCC'07}.
\newblock [\epfmtdoi{10.1007/s00037-009-0267-3}]

\bibitem{spalek}\bibhead{spalek}
{\sc Robert \v{S}palek}: A dual polynomial for {OR}, 2008.
\newblock [\epfmt{arxiv}{0803.4516}]

\end{thebibliography}

\end{document}
